% 本文件由 LaTeX 排版 Agent 根据有效 translation.json 自动生成。
% author=Archelaus; work_id=0616; title=与摩尼辩论行实（阿刻拉俄斯）

\chapter{与\Person{摩尼}辩论行实（\Person{阿刻拉俄斯}）}

% structure: p0001 source=h1 target=omitted reason=page-title-already-rendered-by-parent

1. 真\WorkTitle{宝库}；即，在\Place{美索不达米亚}的\Place{卡尔查}城进行的辩论，以\Person{曼尼普斯}、\Person{埃吉阿雷乌斯}、\Person{克劳狄乌斯}和\Person{克莱奥布鲁斯}为裁判。在这\Place{美索不达米亚}的城里，有一个人名叫\Person{马尔塞卢斯}，因其生活方式、追求和出身而被视为最值得尊崇的人物，也因其明智和品格的高尚而同样受人敬重；他也拥有丰厚的资财；而最重要的是，他以最深的虔诚敬畏\OriginalTerm{天主}{God}，并总是以应有的敬意倾听论及\Person{基督}的事。简而言之，那人身上没有缺乏任何优良品质，因此他在全城备受尊崇；另一方面，他也以慷慨频繁的施舍行为，周济穷人，慰藉受苦者，帮助困厄之人，来丰厚地回报他城邑的善意。但我们就此说这么多已足够了，免得因我们言辞的薄弱，反而减损了那人的德行，而不是述说出配得上其光辉的事。因此，我要转入构成我主题的工作。有一次，当一大群俘虏被驻扎在那里的士兵送给主教\Person{阿刻拉俄斯}时，人数约七千七百人，他因士兵索取释放俘虏所需的大笔赎金而陷入极度的焦虑。而他无法隐藏自己的忧虑，全都燃烧着虔敬之情与对天主的敬畏，终于急忙去见\Person{马尔塞卢斯}，向他解释了事情的严重和困难。当那位虔敬的典范\Person{马尔塞卢斯}听了他的叙述后，毫不迟延地走进自己的家中，按照那些掳掠俘虏的人所定的价值，提供了索取的赎价；又打开了他财物的库藏，立刻将虔敬的赠礼分给士兵，丝毫不考虑数量或区分，以致它们更像是赠礼而非赎金。那些士兵对这人的虔诚和慷慨的宏大充满了惊奇和钦佩，并为这怜悯的榜样所震撼和感动；因此，他们中很多人归依于我们的主\Person{耶稣基督}的\OriginalTerm{信德}{faith}，并丢弃了军役的腰带，而另一些人只带走不过四分之一赎金便返回营地，其余的人甚至没有拿到足够支付旅途费用的数目就离开了。\par

正如所料，马尔切洛对这些事件极为满意；他召来一名叫科尔提尼乌斯的囚犯，向他询问战争的起因，以及他们是如何被击败并被捆绑在囚禁的锁链中的。被问话的人获得发言自由后，开始如下表达：\Quote{我主马尔切洛，我们只信奉永生的天主。我们有一种我将要描述的习俗，这习俗通过我们\Emph{信仰}中的弟兄的传统传承给我们，并且我们一直遵守直到今日。这个做法是，每年我们都带着妻子儿女走出城外，向唯一不可见的天主献上祈求，恳求祂为我们降下雨水，灌溉田地与庄稼。而当我们按通常的时间和惯常的方式庆祝此规矩时，我们在那里逗留，仍守着斋戒，夜晚不知不觉降临。因此，我们感受到人们必须忍受的两件最艰难之事的压力——即守斋与缺乏睡眠。但约在午夜时分，睡意嫉妒而不合时宜地侵袭了我们，脖颈低垂无力，头垂下，使我们的脸撞到膝盖上。这事发生，因为按天主的判断，我们该当承受我们应得的惩罚的时刻近了，无论我们是在无知中犯罪，还是明知有罪却仍不弃绝我们的罪。因此，在那个时刻，一大群士兵突然包围了我们，他们以为——依我判断——我们在那里设下了埋伏，并且是经验丰富、善于作战的人；他们未详细查问我们聚集的原因，便以战争相威胁，不是用言语，而是立即用刀剑。尽管我们是从未学过伤害任何人的男子，他们却用投射武器无情地打伤我们，用长矛刺穿我们，用剑割开我们的喉咙。这样，他们杀了我们当中大约一千三百人，又伤了另外五百人。当天色大亮时，他们将我们中的幸存者作为囚犯带到此处，且方式完全显露出他们对我们毫无怜悯。因为他们驱赶我们在马前行走，用长矛的戳刺催逼我们，用马头推挤使我们前行。那些有足够耐力的人确实挺住了；但很多人却在他们残暴主人的面前倒下，在那里断气身亡；母亲们手臂疲乏，因负重而完全无力，又被身后人的威胁搞得心烦意乱，任凭依偎在怀中的幼儿跌落在地；至于所有年老的人，一个接一个地瘫倒在地，被劳累压垮，因缺乏食物而精疲力尽。然而，那些骄横的士兵却以人们不断丧命的血腥场面为乐，仿佛这是一种娱乐，他们看见有些人倒在土上，绝望地匍匐着，另一些人因极度口渴、舌头的筋络完全干涸而丧失说话能力，还有一些人眼睛不断向后张望，为垂死儿女的命运呻吟，而这些幼儿又不断以哭声向他们极为不幸的母亲哀求，母亲们则因强盗的暴行而发疯，只能以哀哭回应，这的确是她们唯一能自由做的事情。她们中那些与儿女心肠最柔软相连的人，自愿选择与孩子们一同遭遇同样过早的死亡命运；而另一方面，那些尚有些耐力的人则与我们一同被当作囚犯带到这里。如此，过了三天，在此期间我们从未被允许休息片刻，连夜间也不行，我们被押送到此地，而此地此后发生的事情，您自己更加清楚。}\par

3. 当\Person{马塞勒斯}——这位至为虔诚的人——听完这段叙述，不禁泪如泉涌，对如此繁多而深重的灾祸深感怜悯。他毫不迟延，立即为受难者准备食物，亲手服侍疲惫的人；效法我们的圣祖\Person{亚巴郎}：圣祖曾于某次盛情款待天使时，不满足于吩咐奴仆从牛群中牵一头牛犊来，而是尽管年事已高，亲自将牛犊放在肩上扛进去，亲手预备食物，摆放在天使面前。同样，\Person{马塞勒斯}履行类似的职分，吩咐他们十人一桌，作为宾客入席；备齐了七百张桌子后，他满心喜乐地使全体俘虏恢复了精力，以致那些有力量在经受如此遭遇后幸存下来的人，忘却了自己的劳苦，忘记了一切不幸。到了第十五天，\Person{马塞勒斯}仍然慷慨地供应囚犯一切所需时，他觉得，除了那些因需要照料伤口而留下的人之外，应当让所有人都能返回原籍。他为这些人提供适当的治疗，并吩咐其余的人回到自己的家乡和亲友那里。在这些善举之上，\Person{马塞勒斯}更增添了更大的虔诚事迹：他带着一大队自己的仆从，去料理那些在行军中丧生者的遗体安葬事宜；凡他所能寻见的，无论何种身份，他都为他们安排了得体的安葬。办完这事以后，他回到\Place{哈兰}，准许受伤的人待健康充分恢复后，从那里返回家乡，并为他们的旅途备办了极丰富的供应。确实，对这一事迹的评断，为\Person{马塞勒斯}其他高贵行为的\Emph{声誉}增添了一份壮丽。因为在整个地区，\Person{马塞勒斯}虔诚的名声如此盛传，以致众多来自各城的人，被最强烈的渴望所点燃，想要一睹此人并与他结识，特别是那些先前未曾遭受过贫困的人。这位卓越的人物，效法古时的一位\Person{马塞勒斯}，对所有这些人都极为慷慨地施予援助，以致他们都说，没有比这个人更杰出地虔诚的人了。是的，所有信主的寡妇也都投奔他，而软弱无助的人可指望从他手中获得最确定的帮助以应付境遇；同样，孤儿都由他扶养，以致他的家被称为陌路人和贫乏者的收容所。除此之外，他还以卓越而独特的方式持守对信德的热忱，将自己的心建立在不可动摇的磐石上。\par

4. 因此，随着这个人的名声在各处日益广传，甚至越过\Place{斯特兰加河}，他名字的美好声誉传入了\Place{波斯}境内。在此地住着一个人，名叫\Person{摩尼}，当这人的名声传到他耳中时，他心中大费踌躇，思量如何能将\Person{马塞勒斯}诱入自己教义的罗网，希望\Person{马塞勒斯}能成为他教理的支持者。因为他盘算，只要能先把这样一个人笼络住，自己就能成为全省的主人。然而，在实行这计划时，他内心疑虑不安，不知是该立刻亲自去见那人，还是先尝试写信给他。因为他担心，倘若他突然意外地现身，恐怕会遭遇某种不测。最后，出于更为狡黠的计策，他决心写信；召来他的一名门徒，名叫\Person{图尔波}，曾受教于\Person{阿达斯}，把一封书信交给他，命他动身送给\Person{马塞勒斯}。这名追随者于是接信，将信送交到\Person{摩尼}所托付的人那里，整个旅途在五天内完成。上述的\Person{图尔波}在这趟旅途上确实行色匆匆，其间也饱受艰辛劳苦。因为每当他作为一名外乡旅客来到一所客栈时——这些客栈乃是\Person{马塞勒斯}自己出于广施款待而设立的——客栈管理人问他来自何处，他是谁，或是谁派来的，他就回答：\Quote{我属于\Place{美索不达米亚}地区，但目前从\Place{波斯}来，是基督信徒中一位教师\Person{摩尼}派来的。}但他们毫不愿意接待一个陌生的名字，有时甚至将\Person{图尔波}赶出客栈，连饮水都不给他。他不得不日复一日地忍受各地负责馆驿和招待处的人如此的对待，甚至比这更坏的对待，若不是最后他显明自己是给\Person{马塞勒斯}送信的，\Person{图尔波}本会在这旅途中丧命。\par

5. \Person{马塞勒斯}接到信后，当即打开，在当地的\Person{阿尔刻劳斯}主教面前朗读。信的内容如下：——\par

\Person{摩尼}，耶稣基督的宗徒，以及与我同在的所有圣人和贞女，致我亲爱的儿子\Person{马塞勒斯}：愿恩宠、仁慈和平安，来自天主圣父和我们的主耶稣基督，与你同在；愿\OriginalTerm{光明的右手}{right hand of light}保护你，脱离这现今邪恶的世界及其灾难，脱离恶者的罗网。阿们。\par

我极为欣喜地看到你们所怀有的爱，这爱确实极其深厚。但使我感到忧虑的，是你们的信德，因为它并不符合正确的准则。因此，既然我受派遣去寻求提升人类，又怜惜那些陷于欺诈与错谬的人，我认为有必要将这一封信寄给你们，首先是为了你们自己灵魂的救恩，其次也是为了与你们同在者的灵魂的救恩，好使你们得以防备可疑的见解，尤其是防备那些引导心思单纯者的人向他们灌输的看法。他们宣称善与恶具有相同的本源存在，并且为它们设定同一个开端，却不在光明与黑暗之间、善与恶或无价值者之间、以及内在之人和外在之人之间（正如我们先前所述）作出任何区分或辨别，反而不断地将彼此混合混淆。可是，我的儿子啊，你不可如此轻率地，以大众所遵循的那种无理性而愚蠢的方式，将这两样东西混为一谈，也不可将这种混淆归因於美善的天主。因为这些人竟将那些灾祸的起始、终结与渊薮都归於天主自己——\Quote{他们的结局临近诅咒。}因为他们不信我们的救主、主\Person{耶稣基督}在福音中所说的话，即\BibleQuote{好树不能结坏果子，坏树也不能结好果子。}\BibleRef{玛 7:18}至於他们何以竟敢大胆称天主为\Person{撒殚}及其恶行的制造者和策划者，这件事令我极其惊讶。是的，我甚至希望他们愚蠢的努力仅止於此，而未宣称那位从圣父怀中降来的独生子\Person{基督}\BibleRef{若 1:18}是某个女人\Person{玛利亚}的儿子，由血、肉以及妇女特有的各种不洁所生！然而，在此信中我不愿写得过多，也不愿过分打扰你的善良本性——更何况我本无口才的天赋——我将满足於以上所言。但当我亲自与你同在时，你将对整个问题有充分的认识，只要你确实仍关心自己的救恩。因为我不像大众中那些不够审慎的人那样\BibleQuote{给人设下圈套}\BibleRef{格前 7:35}。我最尊贵的儿子，请思想我所说的。\par

6. 读罢此信，\Person{玛尔塞鲁斯}以最亲切的态度，殷勤招待了送信人的需要。另一方面，\Person{阿尔克莱乌斯}却对信中所读的内容不很悦纳，反而\Quote{像被链锁的狮子一般咬牙切齿}，急於让人把信的作者交给他。然而，\Person{玛尔塞鲁斯}劝他保持平静；应许自己会负责确保此人到场：於是，\Person{玛尔塞鲁斯}决定回覆他收到的信，并写下了一封包含以下内容的信函：—\par

\Person{玛尔塞鲁斯}，有声望的人，致通过他的信函让我知晓其人的\Person{摩尼}，问安。\par

你写的信已到我手中，我亦以我惯常的善意接待了\Person{图尔柏}；但你的信的意思，我完全没有领会，也不可能领会，除非你亲自来到我们中间，以面谈的方式详细解释信的内容，正如你本人在信中所提议的那样。再会。\par

他封了这封信，交给\Person{图尔柏}，嘱咐他送给那个他已从他那里送过类似文件的人。然而，这位信差极不情愿回到他主人那里去，因为他记起在旅途中所要忍受的苦楚，便恳求另派他人代替，拒绝再回到\Person{摩尼}那里去，或与他再有任何往来。但\Person{玛尔塞鲁斯}召来了一个名叫\Person{卡利斯图斯}的青年，命他前往那个地方。这个青年毫不耽搁，立刻动身前往；三天之后，他来到了\Person{摩尼}那里，发现他在一个堡垒中，就是\Place{阿拉比翁}的堡垒，并把信交给了\Person{摩尼}。读信之后，\Person{摩尼}很高兴看到自己被\Person{玛尔塞鲁斯}邀请；他毫不迟疑地踏上了旅程；但他预感到\Person{图尔柏}未能回来是不祥之兆，便若有所思地前往\Person{玛尔塞鲁斯}那里去。至於\Person{图尔柏}，他根本没有想过离开\Person{玛尔塞鲁斯}的家；他也不放过任何机会与主教\Person{阿尔克莱乌斯}交谈。因为这两者都非常热心於调查\Person{摩尼}的行为，渴望知道他是谁，从哪里来，以及他说话的方式如何。於是\Person{图尔柏}清楚地讲述了整个状况，以下列方式叙述并阐述了他的信仰信条：—\par

如果你渴望由我领受摩尼的信仰，就请留心听我片刻。那人崇拜两位神明，无始、自有、永恒、彼此对立。他称其中一位为善，另一位为恶；将前者的名定为\Emph{光明}，后者的名定为\Emph{黑暗}。他还声称，人的灵魂是\Emph{光明}的一部分，而身体和物质的形成则是\Emph{黑暗}的部分。他更进一步主张，二者之间发生了一种混合或掺合，其方式即将陈述，并用以下类比来作说明：即，他们的关系可比作两个自始就互相敌对、各有自己地位、各按其序的国王之间的冲突。因此他认为，\Emph{黑暗}越出了自己的边界，与\Emph{光明}进行类似的争斗；但那时，善父察觉到\Emph{黑暗}已来到他的地上寄居，便从自身发出一种能力，称为\OriginalTerm{生命之母}{the Mother of Life}；然后这能力又从自身发出\Emph{\OriginalTerm{初人}{the first man}}和五元素。这五元素是风、光、水、火、物质。这位原始人既已拥有这些，并以此武装，好像准备作战，便降到这些低处，与\Emph{黑暗}交战。但\Emph{黑暗}的王子们反过来与他作战，消耗了他全副武装中那称为灵魂的部分。于是，那个\Emph{初人}在底下被\Emph{黑暗}重重伤害；若非善父听见他的祈祷，派遣第二种能力——那也是从祂自身发出，称为\Emph{\OriginalTerm{活灵}{the living Spirit}}——降下来，向他伸出右手，将他从\Emph{黑暗}的掌握中提拔上来，那初人早在太古之时就有全盘倾覆的危险。因此，从那时起，他便把灵魂遗留在底下。为此，摩尼教徒相遇时互相伸出右手，作为他们已从黑暗得救的标记；因为他认为所有异端都位于黑暗中。然后，活灵创造了世界；他自身带着另外三种能力，降下来，带走众王子，将他们安置在\OriginalTerm{气层}{firmament}中，那是他们的身体（虽称为）\OriginalTerm{球体}{sphere}。接着，活灵又创造了\OriginalTerm{光体}{luminaries}，那是灵魂的碎片，他使它们如此在气层周围旋转；他又创造了八类大地。\OriginalTerm{荷持者}{the Omophorus}在下肩负其重担；当他负重疲乏时便会颤抖，就这样使大地在不循固定周期的时间发生震动。为此，善父从祂自己的怀中将子派到地的中心，及大地的最深之处，以便为他取得合宜的纠正。每当地震发生，他或因疲乏而颤抖，或将重担从一边肩膀移到另一边。此后，物质自己又产生出各种生长物；当这些生长物被一些王子掳去为掠物时，他召聚了所有为首的王子，从他们个个取了又取能力，造成了按那\Emph{初人}形象的人，并将灵魂（连同这些能力）结合在他里面。这就是关于他的构成是如何设计出来的叙述。\par

8. 但当\OriginalTerm{生活的父}{the living Father}察觉灵魂在肉体中受苦，因满怀仁慈与怜悯，便派遣自己的爱子，为灵魂的救恩而来。这一点，连同\OriginalTerm{奥莫弗鲁斯}{Omophorus}之事，即是他派遣他的缘由。于是子来了，将自己变化成人的样式，以人的形象向人显现，而他实际上并不是人，人却以为他是被生的。他如此来临，预备了用以完成灵魂救恩的工作，并为此目的制造了一个带有十二个瓮的器具，这器具由天球推动旋转，并将垂死者的灵魂一起提取上来。更大的光体接收这些灵魂，以其光线净化它们，然后将它们传递给月亮；如此，我们所谓的月轮，便得以充满。因为他说这两个光体是船只或渡船。然后，如果月亮变满，它就将其乘客渡往东风方向，并借此在卸下货物后使自己亏损。如此，它持续进行渡运，并一再卸下由瓮提取上来的灵魂之货物，直到它救出自己份额内的灵魂。此外，他主张每一个灵魂，是的，每一个活动的活物，都分享了善父的实体。因此，当月亮将灵魂的货物交付给父的\OriginalTerm{爱安}{aeons}时，它们就居住在那被称为\OriginalTerm{完美之气}{perfect air}的\OriginalTerm{光荣之柱}{pillar of glory}中。而这气是\OriginalTerm{光之柱}{pillar of light}，因为它充满了正在被净化的灵魂。这便是灵魂得救的机制。但是，以下则是人死亡的缘由：有一位\OriginalTerm{童贞女}{the virgin}，容貌美丽，衣着华美，言辞极有说服力，企图掠夺那些被\OriginalTerm{生活的灵}{the living Spirit}提升并钉在穹苍上的掌权者；她向掌权者显现为美丽的女性，但向女掌权者则显现为清秀迷人的青年男子。掌权者们看着她辉煌的形象，被爱刺中；由于无法占有她，他们便因爱欲之火猛烈燃烧，失去一切理性。当他们如此追逐童贞女时，她从视野中消失。于是，\OriginalTerm{大君}{the great prince}在愤怒中从自己发出乌云，意图使全世界陷入黑暗。然后，如果他感到极度压抑，他的疲惫便以出汗表达，恰如人劳苦时出汗；他的汗形成雨。与此同时，\OriginalTerm{收获之君}{harvest-prince}，如果他恰好也被童贞女迷住，便在全地散布瘟疫，意图置人于死地。这个（人的）身体也被称为一个\Emph{\OriginalTerm{科斯摩斯}{cosmos}}，即小宇宙，相对于大的\Emph{\OriginalTerm{科斯摩斯}{cosmos}}，即宇宙的大宇宙；所有人都有根，下面的根与上面的根相连。因此，当这掌权者被童贞女的魅力迷住，他便开始剪断人的根；当他们的根被剪断，瘟疫便开始爆发，人便如此死亡。如果他猛烈摇动根的上部，便发生地震，这是奥莫弗鲁斯所遭受的震动所导致的结果。这就是（死亡现象的）解释。\par

9. 我也要向你们说明，灵魂是如何被转注入五体的。首先，在这过程中，它的一小部分得以净化；然后它被转注入狗、骆驼或其他某种动物的身体。但如果灵魂犯了杀人罪，就被转换到\OriginalTerm{凯勒菲}{celephi}的身体中；若被发现从事了切割，就被转入\Emph{哑巴之体}。这些就是灵魂的称谓——即理智、反省、明智、思考、推理。此外，收割的收割者被比作那些自起初就在黑暗中的\OriginalTerm{掌权者}{princes}，因为他们消耗了\OriginalTerm{原人}{first man}的一部分全副武装。为此，他们必然被转换到干草、豆类、大麦、谷子或蔬菜中，以便在这些形态中，他们也以同样的方式被收割和切割。再者，若有人吃面包，他也必须成为面包并被吃掉。若有人杀了一只鸡，他自己就将成为一只鸡。若有人杀了一只老鼠，他自己也将成为一只老鼠。再者，若有人在此世富有，则必须在他离开自己身体的帐幕时转入一个乞丐的身体，四处求乞，然后他将进入永罚。而且，由于这身体属于掌权者和\OriginalTerm{物质}{matter}，栽种\OriginalTerm{珀耳塞树}{persea}的人必须经过许多身体，直至那棵珀耳塞树被压倒。若有人为自己建造房屋，他就会被分裂并分散在所有的身体中。若有人沐浴，就会冻结他的灵魂；若有人拒绝向自己的\OriginalTerm{选民}{elect}表示虔诚的敬重，他将会世世代代受罚，并被转换到\OriginalTerm{慕道者}{catechumens}的身体中，直到他献上许多虔诚的贡物；为此，他们把自己食物中最好的部分献给那些选民。当他们要吃面包时，他们首先献上祈祷，对着面包说这样的话：\Quote{我未曾收割你，未曾碾磨你，未曾揉捏你，也未曾把你放进烤炉里；而是别人做了这些事，把你带到我这里，我吃了你却没有过错。}当他对面包自言自语说了这些话之后，他就对慕道者说：\Quote{我为你祈祷了；}然后那人就此离去。因为，正如我刚才向你们指出的，若有人收割，他必被收割；同样，若有人把谷子投入磨中，他自己也必同样被投入；若有人揉面，他必被揉捏；若有人烤制，他必被烤制；为此，他们被禁止从事任何此类工作。此外，还有某些其他的世界，那些\OriginalTerm{光体}{luminaries}在我们的世界落下后就在那些世界上照耀起来。又若有人在这地上行走，他就伤害了大地；若他移动他的手，他就伤害了空气；因为空气就是人和一切活物的灵魂（生命），无论是飞鸟、游鱼，还是爬虫。至于在这世上的每一个人，我已告诉你们，他的这个身体不属于天主，而是属于\OriginalTerm{物质}{matter}，本身就是黑暗，因此必须被投入黑暗。\par

10. 现在，关于乐园，它并不被称为一个\Emph{宇宙}。其中的树是贪欲和其他诱惑，败坏那些人的理性官能。而在乐园中那棵使人得以认识善的树，就是耶稣自己，\Emph{或是}对在世上的他的认识。谁若参与其中，就能辨明善与恶。然而，世界本身并非天主的\Emph{工程}，而是由一部分物质构建而成，因此万物在其中都会腐朽。而掌权者从原人那里夺来的战利品，就使月亮盈满，也正日复一日地从世界中净化出去。如果灵魂离世时尚未获得真理的知识，就会被交给魔鬼，以便他们在火中\OriginalTerm{革哈拿}{Gehennas}里制服它；经历了这番管教之后，它被转入各种身体，以求使之顺服，并以此方式被投入烈火之中，直到最终的结局。再者，关于你们中间的众先知，他这样说：他们的精神乃是起初兴起的黑暗的邪恶或非法。他们受这精神欺骗，没有宣讲\Emph{真理}；因为那掌权者蒙蔽了他们的心神。若有人听从他们的话语，便永远死亡，被捆锁于土块中，因为他没有习得\OriginalTerm{帕拉克莱特}{Paraclete}的知识。他又独独将训令交给他的选民，数目不超过七个。那训令是这样的：\Quote{你们停止进食时，要祈祷，并将一枚橄榄放在你们的头上，并凭着呼求许多名字起誓，以证实此信仰。}然而，那些名字并没有为我知晓；因为只有这七个人使用它们。再者，你们所尊崇大有能力的\OriginalTerm{撒巴奥特}{Sabaoth}这名字，他宣告说，它乃是人的本性，又是\OriginalTerm{欲望}{desire}的源头；因此，愚昧人敬拜欲望，并视之为神明。再说关于亚当的创造方式，他告诉我们说，那位说\BibleQuote{来，让我们照我们的肖像，按我们的模样造人}~\BibleRef{创 1:26}，或说\Quote{照我们所见的形式}的，就是那位掌权者；他向其他掌权者说话，其意旨可如此解读：\Quote{来，将我们领受的光分给我，让我们照我们掌权者的形式造人，就是照我们所见的形式，也就是原人。}他就以这种方式创造了人。他们也以类似的方式创造了厄娃，将她自己的一部分贪欲赋予她，为要欺骗亚当。世界便通过这些手段从掌权者的运作中衍生而来。\par

11. 他还主张，天主与世界本身毫无关系，也不喜悦它，因为世界一开始就被掌权者所掳掠，并且因那临到世界的邪恶。因此，他日复一日地差遣并通过光体——就是日月——，取走那些属他自己的灵魂；日月掌管着整个世界和一切受造物。再者，他宣称那位曾与\Person{梅瑟}、犹太人及司祭们说话的，乃是\OriginalTerm{黑暗之王}{prince of the darkness}；所以基督徒、犹太人和外邦人其实是同一个身体，敬拜着同一位天主：因为这天主以自己的情欲诱惑他们，并非\OriginalTerm{真理之天主}{God of truth}。为此缘故，所有寄望于那位与\Person{梅瑟}及先知们说话之天主的人，都必须与那神祇一起被捆绑，因为他们没有寄望于\OriginalTerm{真理之天主}{God of truth}；因那神祇是按照他们自己的情欲同他们说话的。此外，在这一切之后，他论到结局时如此说，也如此写着：\Quote{当长者展示他的形象后，\OriginalTerm{欧模弗罗斯（负载者）}{Omophorus}就任地离开他们，于是大火得释放，吞灭全世界。然后，他又让大地随\OriginalTerm{新永世}{new aeon}而去，使罪人们的灵魂永远被捆绑。}这些事将要发生在那人的形象临在之时。而所有这些由天主所发出的众能力——就是在较小船上的耶稣，与生命之母，与十二位舵手，与光之童女，与在较大船上的第三位长者，与活灵，与大火之墙，与风墙，与气，与水，与内在活火——都驻在较小的光体中，直到火吞灭全世界为止：这事将发生在许多年之内，然而其确实的数目，我尚未查明。在此之后，将有二性体的恢复；掌权者将各归其本位的下界，而圣父将占据上界，收回他当得的产业。——他将这全部的教义传授给他的三个门徒，并吩咐他们各自去往不同的地界。\Place{东方}地区归于\Person{阿达斯}；\Person{多默}承受了\Place{叙利亚}地区为业；还有另一位，就是\Person{赫梅亚斯}，往\Place{埃及}去了。直到今日，他们仍在那里侨居，为要确立这教义中所包含的各项命题。\par

12. 当\Person{图尔博}陈述完毕，\Person{阿基劳斯}极为激动；但\Person{玛尔凯路斯}仍不为所动，因为他期待天主会来援助他的真理。然而，\Person{阿基劳斯}更因挂虑民众而忧心如焚，如同牧人担忧自己的羊群，当隐秘的危险从狼群中威胁他们一样。于是\Person{玛尔凯路斯}给了\Person{图尔博}极为丰厚的礼物，并吩咐他留在主教\Person{阿基劳斯}的家中。但就在同一天，\Person{摩尼}到了，带着一些被拣选的青年男女，共二十二位。他首先在\Person{玛尔凯路斯}家门口寻找\Person{图尔博}；没有找到，便进去向\Person{玛尔凯路斯}问安。\Person{玛尔凯路斯}一见到他，起初因他的装束而惊讶不已。因为他穿着一双通常俗称\Quote{四底鞋}的鞋子；还披着一件彩色外氅，显得颇为飘逸；手中握着一根非常结实的乌木杖；左臂下夹着一本巴比伦文书籍；双腿裹着颜色各异的裤子，一条为红色，另一条如韭菜般翠绿；整个神态就像一位老波斯将领和指挥官。于是\Person{玛尔凯路斯}立刻派人去请\Person{阿基劳斯}，他来得极快，几乎话未出口人已到，一进门便几乎忍不住要当面发作，因为光是看到他的装束和容貌就被激怒了，更因为他自己退居静处时已反复思考从\Person{图尔博}叙述中学到的各种事情，因而作好了周全的准备。但\Person{玛尔凯路斯}深谋远虑，抑制了一切单为争辩的热忱，决定听取双方的陈述。为此，他邀请了城中的领袖人物；并从他们当中挑选了四位信奉外邦宗教的人士，作为\Emph{辩论的}裁判。这些公断人的名字如下：\Person{玛尼颇斯}，精通文法与修辞之术；\Person{埃吉亚劳斯}，极杰出的医师，且以博学享誉最高；\Person{克劳狄}和\Person{克莱奥博路斯}，两位以修辞家而闻名的兄弟。于是，一次盛大的集会就这样召集起来了；人数如此众多，以致\Person{玛尔凯路斯}那极大无比的寓所，都挤满了应邀前来聆听的人。当准备彼此对辩的双方在众人面前就位后，被选为裁判的人便就座于高出众人的席位上：而开始辩论的任务便交给了\Person{摩尼}。于是，在全场肃静之后，他以如下的言辞开始了讨论：——\par

13. 我的弟兄们，我确实是基督的门徒，并且是\Person{耶稣}的宗徒；我之所以急忙来到这里，是由于\Person{马塞禄}的极大善意，为要清楚地指示他应当如何持守神圣宗教的体系，使所说的\Person{马塞禄}，如今他如同一个将自己交付为囚的人，置身于\Person{阿尔刻劳}的教导之下者，不致像那些没有理智、不明白自己所作所为的哑巴畜生一样，因缺少进一步正确地遵行神圣敬礼的方法，而遭受致命打击，导致灵魂的毁灭。再者，我知道并且确信，一旦\Person{马塞禄}被纠正，你们所有人也很有可能获得救恩；因为你们的城悬于他的判断。如果你们每个人都能弃绝虚妄的自恃，并以真正对真理的热爱聆听我所要宣讲的，你们将要领受来世的产业和天国。我实在是护慰者，就是我昔日由\Person{耶稣}所预告将来到的那一位，要来\BibleQuote{指证世界关于罪恶和正义}。正如在我之前被派遣的\Person{保禄}论及自己时所说，\BibleQuote{他现今所知道的，只是局部的，所预言的，也只是局部的}\BibleRef{格前 13:9}，同样，我为自己保留了完全的事，为要除去那局部的。所以你们要接受这第三个见证，证实我是基督所选派的宗徒；如果你们选择接受我的话，你们必将获得救恩；但若拒绝，永火必将吞灭你们。因为如同\Person{依默纳约}和\Person{亚历山大}被交付于撒殚，为叫他们学习不要亵渎，你们所有人也都要被交付于刑罚之王，因为你们伤害了基督的圣父，声称祂是一切恶的原因、不义的创立者、一切邪恶的创造者。借着这样的教义，你们实实在在地从同一水源中涌出甜水和苦水——这无论如何也是不可能做到或理解的。应当相信谁呢？是你们那些以肉身为享乐、以极奢华享用养肥自己的师傅们，还是我们的救主\Person{耶稣}基督，正如福音书中所记载的，祂说：\BibleQuote{好树不能结坏果子，坏树也不能结好果子}，又在另一处向我们保证说，\BibleQuote{魔鬼的父从起初就是杀人的凶手，也是撒谎者}，又告诉我们，人们爱黑暗，因此不愿跟随那从起初由光中发出的圣言，且（又给我们指出）那与之作对的人，就是仇敌，撒稗子者，今世的神和王子，他蒙蔽了人的理智，使他们不服从基督福音中的真理？那位不希望自己人获得拯救的天主是善的吗？为避免涉及许多其他事情而浪费大量时间，我可以把真实教义的阐明推迟到另一个机会；暂且认为我已就此事说得够多了，现在回到眼前的事，尽力满意地证明这些人教义的荒谬，并说明这一切都不可归于我们的主和救主的天主圣父，而我们必须把撒殚视为一切灾祸的原因。当然，这一切都要归于他，因为这类灾祸都是由他所生。而且先知书和法律书上所写的那些事，也同样要归于他；因为正是他那时在先知们内说话，将许多对天主的愚昧观念，以及诱惑和偏情，灌输进他们的心里。他们也提出那吞食血肉者；而撒殚和他的先知们，正适于承受那些他企图转移给基督的圣父的一切事，因为他准备写一些看似真理的话，好借此也使他的其他谎言获得信任。因此，我们最好完全不接受那些从古时直到\Person{若翰}所写的一切，而只拥抱从他那日子以来在福音中所宣讲的天国；因为他们实在是自嘲，引入了一些荒谬可笑之事，保存了法律中一些含混地勾勒出的小小言论，却不明白，如果好事与恶事混杂，结果便是这些恶事败坏，连那些好事也一同毁灭了。而且，如果有人能证明法律维护正义，那么那法律就当遵守；但如果我们能指明它是恶的，那么就应废除并拒绝它，因为它包含那刻在石头上的属死的职务，这职务又掩盖并销毁了\Person{梅瑟}面容上的荣光。因此，你们任何人若将新约与法律和先知书一同教导，仿佛它们同出一源，那可不是没有危险的；因为我们救主的知识日日更新\Emph{那一个}，而另一个却渐旧渐衰，几乎趋于全然毁灭。这是那些能运用辨别力的人明显可见的事实。正如树的枝子枯老，或树干不再结果，就被砍下；又如身体的肢体生了坏疽，便被切除，因为坏疽的毒素从这些肢体扩散至全身，除非医师的技艺找到医治此病的良药，全身都将被损坏；同样，你们若不理解法律的起源就接受了它，你们必将毁灭自己的灵魂，丧失救恩。因为\BibleQuote{法律及先知到\Person{若翰}为止}；但自从\Person{若翰}以来，真理的法律、恩许的法律、天上的法律、新的法律，已向人类彰显了。确实，只要没有人向你们展示我们主\Person{耶稣}的这一最真实的知识，你们就没有罪。然而现在，你们既看见又听见，却仍愿在无知中行走，为要持守那已被摧毁和废弃的法律。\Person{保禄}，他在我们中间被认为是最受称许的宗徒，也在他的一封书信中表达了同样的意思，他说：\BibleQuote{如果我重建我所拆毁的，我就表明我是罪犯}。他说这话，是把他们当作外邦人，因为他们在信仰的圆满来到之前，还屈服于世俗的原理，那时他们只相信法律和先知。\par

14. \Emph{裁判官们说}：如果你还有任何更清楚的陈述需要作出，就向我们解释一下你的教义的性质和你信仰的名称吧。\Emph{\Person{摩尼}回答说}：我认为有两种本性，一种是善的，另一种是恶的；善的本性确实居住在它自己特定的某些部分中，而恶的本性就是这个世界，以及其中所有的一切，它们如同被囚禁在那恶者领域内的物件一般被安置在那里，正如 \Person{若望} 所说，\BibleQuote{全世界都屈服于恶者}，而不是在天主内。因此，我们一直主张有两个场域——一个是善的，另一个在这善的场域之外，这样，\Emph{在他的}场域内有空间，就能将那受造物，即世界的\Emph{创造物}，接收到它自身之中。因为如果我们说只有一本性的唯一统治，天主充满万物，在祂之外没有场域，那么受造物，即\Emph{创造物}的支撑者又是什么呢？火的\OriginalTerm{革哈辣}{Gehenna}将在哪里？外面的黑暗在哪里？哀号在哪里？我难道要说在祂自身内吗？天主不许！不然，祂自己也将要与这些一同受苦了。你们谁若关心自己的救恩，就不要怀有这种幻想；因为我要给你们举一个例子，好使你们能更充分地理解真理。世界是一个容器；如果\OriginalTerm{天主的本体}{substance of God}已经充满了这整个容器，那么现在怎么可能还有任何别的东西能被放进这同一个容器里呢？如果它是满的，那么，除非把容器的某部分倒空，否则它又怎能接纳放进其中的东西呢？否则，那要倒出去的东西将去往哪里呢？因为那里并没有场域容纳它。那么，大地在哪里？诸天在哪里？深渊在哪里？星辰在哪里？居所又在哪里？众德能呢？\OriginalTerm{掌权者}{princes}呢？外面的黑暗呢？是谁为这些建立了基础？又在哪里？没有人能告诉我们这些而不堕入亵渎之言。再者，如果没有自存的物质，祂又怎能造出这些受造物呢？因为如果祂是从无中造了它们，那么这些可见的受造物就应当是优越的，并且充满一切德能。可是，如果这些东西里面有邪恶、死亡、败坏，以及一切与善相对立的事物，我们又怎能说它们是由与它们自身不同的本性所形成的呢？然而，如果你们思考人之子的生育方式，你们就会发现，人的创造者不是主，而是另一个存有，这存有本身也是出于非受生的本性，他既无创始者，也无创造者，也无制造者，而是照他本然的样子，仅由自己的邪恶所生出。与此相应，你们男人与你们的妻子交媾，这种交媾是借着下述本性的事件而发生在你们身上的。当你们当中任何人吃饱了属肉的食物和其他各种食物之后，私欲偏情的冲动就在他里面升起，这样，生育儿子的快乐就加增了；这件事的发生，并非好像它是源于任何德行、哲学，或任何其他心灵禀赋，而是仅仅源于饱食，以及色欲与邪淫。又怎能有人告诉我，我们的父亲 \Person{亚当} 是按天主的肖像所造，并同祂的相似，且 \Person{亚当} 与创造他的天主相似呢？又怎能说，我们这些由他而生的人都与他相似呢？是啊，相反地，我们不是有着极大的形貌差异，并且带着不同的面容印记吗？我将用比喻向你们说明，这是多么真实。例如，看看一个人，他想要封存一件宝物或某个其他物件，你们会看到，当他弄到一点蜡或泥土时，他就设法用自己佩戴的戒指，在上面盖上他自己面容的印记；可是，如果另一个面容也以类似的方式，在物件上印上它自己的形象，印记还会相似吗？绝不相似，尽管你们可能不愿承认这是真的。然而，如果我们不是在共同的印记中彼此相似，而是在我们里面有各种差异，那么这岂不证明我们乃是\OriginalTerm{掌权者}{princes}和物质的作品吗？因为正是依照着他们的形体、相似和肖像，我们也就以各种不同的形体存在。如果你们愿意充分受教，了解那起初发生的交媾，以及这事发生的方式，我就会把这事向你们说明。\par

\Emph{法官们说}：我们首先需要看到已证明有两种本原性的原则，然后才需探究那原初的交合是以何种方式发生的。因为，一旦明确存在两种\OriginalTerm{非受生}{unbegotten}的本性，我们可能也会接受你的其他断言，即使它们中有一些似乎与可信之事不太相符。既然审判的权力已托付给了我们，我们就将宣告那在我们心中变得显明的事情。然而，我们也可准许\Person{阿尔克劳}自由地对你的这些陈述发言，这样，通过比较你们各人所说的话，我们就能按照真理作出裁决。\Emph{\Person{阿尔克劳}说}：然而，这\OriginalTerm{敌手}{adversary}的意图充满了极大的放肆和亵渎。\Emph{\Person{摩尼}说}：法官们，请听，他论及\Emph{\OriginalTerm{敌手}{adversary}}时说了什么。那么，他承认存在两个对象。\Emph{\Person{阿尔克劳}说}：在我看来，此人满是疯狂而非明智，他今天竟要与我挑起争论，只因我偶尔提到了\OriginalTerm{敌手}{adversary}。不过你这个异议只需寥寥数语便可消除，尽管你从我的这话中猜测我将容许有这两种本性。你提出了一种极其荒谬的教义；因为你所作的断言没有一条站得住脚。因为完全有可能，一个并非由于本性、而是由于决断而成为\OriginalTerm{敌手}{adversary}的人，可以被结交为朋友，从而不再是敌手；这样，当我们中的一方已默许了另一方，我们两个就显得仿佛是同一个对象了。这一说明也表明，理性的受造物被赋予了自由意志，正是因着这自由意志，他们才也能转变。因此，不可能存在\Emph{两种}\OriginalTerm{非受生}{unbegotten}的本性。那么，你怎么说呢？这两种本性是不可转变的吗？还是可转变的？又或是其中一种会转变？\Emph{\Person{摩尼}}却退缩了，因为他找不到合适的答复；他正在心里琢磨，若作两种回答中的任何一种，会得出什么样的结论，如此深思着：如果我说它们可转变，他就会用\BibleBook{}{福音}中关于树的记载来反驳我；但如果我说它们不可转变，他必然要问我它们相互混杂的情况和原因。此时，稍加拖延之后，\Emph{\Person{摩尼}答道}：就相反者而言，两者确实都不可转变；但就属性而言，它们却是可转变的。\Emph{\Person{阿尔克劳}于是说}：依我看，你似乎已经失去理智，并且忘记了你自己的命题；是的，你甚至似乎认不出你所学的那些词本身的能力或性质。因为你既不明白转变是什么，也不明白\Emph{\OriginalTerm{非受生}{unbegotten}}的含义，不明白二元性意味着什么，不明白什么是过去、什么是现在、什么是未来，我正是从你刚才所表达的意见中获知这些的。因为你曾断言，就相反者而言，这两种本性各自不可转变，但就属性而言，它们却是可转变的。但我却说，一个在属性中变动的人，并不走出自身，而是持存于那同样的属性中，并在这属性中永远不可转变；而对于一个能接受转变的人来说，其效果却是：他被置于属性之外，而进入\OriginalTerm{附性}{accidents}的领域中。\par

16. \Emph{法官们说：}\Quote{\OriginalTerm{可转变性}{convertibility}将它所临到的人转变成为另一个人；例如，我们可以说，如果一个犹太人决心成为基督徒，或者反过来，如果一个基督徒决定成为外邦人，这就是一种可转变性，并且是其原因。但是，同样，如果我们设想一个外邦人保持他所有\Emph{异教}的特质，向他的神明献祭，并且像往常一样在殿宇中服事，那么当他仍保持着他的特质并继续在其中时，你肯定不会认为他可以被说成是转变了，对吗？那么，你们怎么说？它们是否具有可转变性呢？}\Emph{摩尼犹豫时，阿凯劳斯接着说：}\Quote{如果他说两种本性都是可转变的，那么有什么能阻止我们认为它们是同一个事物呢？因为如果它们是无可转变的，那么在同样无可转变且同样\OriginalTerm{非受生}{unbegotten}的本性中，就没有任何区别，也不能将其中一种识别为善或恶。但是，如果两者都是可转变的，那么结果可能是善变为恶，恶变为善。然而，如果这是可能的结果，为什么我们不只说一个非受生的呢？这将是一个更符合真理推算的概念。因为我们必须考虑那恶者最初是如何成为这样的，或者在天地形成之前他对什么对象施行了他的邪恶。当诸天尚未出现，大地尚未存在，也没有人也没有动物时，他对谁施行他的邪恶呢？他不义地压迫了谁？他抢夺并杀害了谁？但是如果你说他首先在他的恶本性中向自己的同类显现，那么毫无疑问你证明了他是出于善的本性。再者，如果所有这些都是恶的，撒殚又怎能驱逐撒殚呢？但当你在这个问题上陷入困境时，你可能会在讨论中改变立场，说善遭受了恶的强暴。但即使这么说，对你来说仍非没有危险，因为你肯定了光的被战胜；因为被战胜的离毁灭不远了。因为天主的话是怎样说的呢？\BibleQuote{或者，一个人怎样进入壮士的家，抢走他的财物？除非他先把壮士捆起来，然后才能抢走他的家。}\BibleRef{玛 12:29}但是，如果你声称他首先以恶本性向人显现，并且仅从那时起才公开显示他邪恶的标记，那么在此之前他是善的，并且他具有了这种转变的性质，因为人的受造被发现有成为他邪恶的原因。但是，最后，让他告诉我们他所理解的恶是什么，免得他或许只是在为一个虚名辩护或设立。如果他所谈论的不是恶的名号而是恶的实体，那么让他将这种邪恶和不义的果实摆在我们面前，因为树的本质只能由它的果实得知。}\par

17. \Emph{摩尼说：}先请你们承认，有一个天主未曾种植的\OriginalTerm{外来的邪恶之根}{alien root of wickedness}，然后我便告诉你们它的果实。\Emph{阿尔赫劳斯说：}真理的核算并无这样的要求；我也不会向你们承认存在任何此类恶树的根，因为没有人尝过它的果实。但正如一个人想要购买物品时，除非他先通过品尝确定那东西是干的还是湿的种类，否则他不会拿出钱来，同样，除非首先展示其果实的品质，否则我也不会向你们承认那树是恶的且全然败坏；因为经上记载说：\BibleQuote{树是由果子认出来的。}\BibleRef{玛 7:20} 因此，摩尼啊，请告诉我们，那被称为恶的树结出什么果实，或说它是什么性质，又是什么品质，好叫我们也能与你们一同相信那树的根，就是你们所描述的那种性质。\Emph{摩尼说：}那根确实是恶的，那树极尽败坏，但其生长并非来自天主。再者，奸淫、通奸、谋杀、贪婪，以及一切恶行，都是那邪恶之根所结的果实。\Emph{阿尔赫劳斯说：}要我们相信你说的这些是那邪恶之根的果实，就让我们尝尝这些东西吧；因为你已宣称这树的实体是\OriginalTerm{非受生的}{ungenerate}，其果实是按自己的模样产生的。\Emph{摩尼说：}存在于人内的不义本身就提供了证据，而在贪婪中你也可以尝到那邪恶之根。\Emph{阿尔赫劳斯说：}好吧，那么，照你所提出的论点，在人间盛行的这些不义就是这树的果实。\Emph{摩尼说：}正是如此。\Emph{阿尔赫劳斯接着说道：}那么，如果这些是果实，就是说人的邪恶行为，其结果将是人自己占据了根和树的位置；因为你声明他们产生这种性质的果实。\Emph{摩尼说：}\Emph{这就是我的说法}。\Emph{阿尔赫劳斯回答说：}你这么说可不太对，说什么\Quote{这就是我的说法}：因为这肯定不可能是你的说法；否则，当人停止犯罪时，这棵邪恶之树就会显得没有果实了。\Emph{摩尼说：}你说的不可能发生；因为即使一个、两个或几个人停止犯罪，仍然还会有别人继续作恶。\Emph{阿尔赫劳斯说：}如果正如你所承认的，一个、两个或几个人有可能不犯罪，那么所有人也都有可能同样做到；因为他们都出自同一祖先，都是由同一团泥土造成的人。而且，我不会就此轻易放过你那些混乱提出的、充满荒谬的断言，我将用一些无可辩驳的反证来结束对它们的驳斥。你是否断言，那邪恶之根和邪恶之树所结的果实是人的行为，也就是说，奸淫、通奸、发虚誓、谋杀，以及其他类似的事？\Emph{摩尼说：}我断言。\Emph{阿尔赫劳斯说：}好吧，那么，如果人类从地面上灭绝，以至于他们再也不能犯罪，那树的实体就会消亡，它也就不再结果实了。\Emph{摩尼说：}你所说的这件事何时会发生？\Emph{阿尔赫劳斯说：}将来的事我不知道，因为我不过是个人；尽管如此，我不会不对你这些话加以审察。关于人类，你怎么说？它是\OriginalTerm{非受生的}{unbegotten}，还是被造的产物？\Emph{摩尼说：}它是被造的产物。\Emph{阿尔赫劳斯说：}如果人是产物，那么奸淫、通奸之类的事，其父母是谁？这又是谁的果实？在人被造之前，有谁犯奸淫、通奸或谋杀呢？\Emph{摩尼说：}但如果人是那邪恶本性所塑造的，显然他就是这样的果实，他无论犯罪与否，都如此；由此，人类的名称和族类一劳永逸地、绝对地具有这种特性，不管他们行义还是行不义。\Emph{阿尔赫劳斯说：}好了，我们也可以注意那件事。如果正如你所断言，那邪恶者亲自造了人，为什么他还要对人施行他的恶意呢？\par

18. \Emph{法官们说：}我们希望你就此点向我们说明，\Person{玛尼凯乌斯}，即你断言他是邪恶的，究竟是什么意思。你是说自从人被造之时起他就如此，还是在那个时期之前呢？因为你必须从他成为邪恶的时刻起，就证明他的邪恶，你要确信，除非先尝一尝，否则无法断定酒的好坏；你要明白，同样，每一棵树都是凭它的果实来认识的。那么，你怎么说呢？这个人格是从何时起成为邪恶的？因为我们觉得有必要解释这个问题。\Emph{\Person{玛尼}说：}他始终如此。\Emph{\Person{阿尔柯劳斯}说：}那么，最优秀的朋友们和最明辨的听众们，我也要从这一点表明，他的说法绝非正确。因为，以铁为例，铁并非自始便是邪恶的事物，而是自人类存在以来，并且自从人的技艺将它误用于虚假的用途，从而将其变成邪恶以来；而且，每一个罪都是自人类存在以来才出现的。就连那条大蛇本身，在人类之前也并非邪恶，只是在人类之后，他才在其中显露出他邪恶的果实，因为他自己愿意如此。因此，如果邪恶之父是在人类\Emph{出现}之后才出现在我们面前，根据圣经，那么他既然是在人（人本身是一个造物）之后才被构成恶的，又怎能是\OriginalTerm{非受生的}{unbegotten}呢？然而，再问一下，为什么他偏偏是从你假设他亲自创造人类的那个时刻起，才表现出自己是邪恶的呢？他在人里面渴望得到什么？如果人的整个身体都是他自己的作品，他在人里面热切地追求什么呢？因为一个热切追求或渴望什么的人，他所追求的是某种不同的、更好的东西。的确，如果人是出于他而源自恶的本性，那么正如我经常表明的，我们看到人是属于他自己的。因为如果人是属于他自己的，那么他自己也是邪恶的，这正如我们关于相似的树和相似的果子的例子一样；因为，照你说的，坏树结坏果子。既然所有的人都是邪恶的，那么他还有什么可渴望的，或者他能在何事上显出他邪恶的开始呢，如果从人类形成时起，人就是他的邪恶的原因呢？此外，既然律法和诫命已经赐给了人自己，人绝无能力顺从蛇和蛇所说的话；倘若人当时不服从他，那么他会有什么邪恶的诱因呢？然而，若邪恶是\OriginalTerm{非受生的}{unbegotten}，那么人有时被发现比邪恶更强大，这怎么可能呢？因为，人顺从了天主的律法，就会常常胜过一切邪恶的根源；然而，一个仅仅是造物的人，竟被发现比\OriginalTerm{非受生的}{unbegotten}更强大，那将是荒谬的。再者，那条附有诫命的律法，我指的是那条赐给人的诫命，是谁的呢？毫无疑问，人们会承认它属于天主。那么，律法怎能赐给外邦人呢？或者，谁能将自己的诫命赐给仇敌呢？或者，说到那领受诫命的人，他怎能与魔鬼争战呢？也就是说，根据这种假设，他怎能反抗自己的创造者呢？就好像儿子虽然因善行而亏欠于父亲，却偏要伤害父亲一样。因此，你只是标明人在这一方面的无用，假如你以为他凭律法和诫命在反抗创造他的那位，并试图胜过他。诚然，我们还不得不设想，魔鬼自己竟会愚昧到如此地步，以致未曾察觉在造人时为自己造了一个敌对者，既没有考虑他未来的情况，也没有预见到自己行为的实际后果；而就算在我们这些仅是造物的存在者身上，也至少有一些细小的知识恩赐、一份辨别能力、一种适度的考虑，有时这种考虑还是相当可靠的。那么，我们怎能相信，在\OriginalTerm{非受生的}{unbegotten}那里，连一点点辨别能力、考虑，或理智都没有呢？或者，根据你的断言，我们怎能做出相反的假设，即发现他最愚拙无知，心最迟钝，简而言之，其天性构成更近似于禽兽呢？然而，如果情况果真如此，那么，既然人在心智能力和知识上拥有绝非微不足道的力量，又怎能从那位在众生中最为无知、领悟最为迟钝的主接受自己的\OriginalTerm{实体}{substance}呢？谁会轻率地宣称，人就是这种品格者的作品呢？然而，如果人是由灵魂和身体构成的，而不仅仅是无灵魂的身体，并且如果二者不能彼此分离而独存，那么你为什么要说这二者是彼此对立、相互矛盾的呢？因我们的主耶稣基督，在我看来，似乎在祂的比喻中谈到了这些，当祂说：\BibleQuote{也没有人把新酒装入旧皮囊里的；不然，皮囊一裂，酒也流了，皮囊也坏了；而是应把新酒装在新皮囊里。} 因为皮囊和酒确实有同一个主。因为，尽管\OriginalTerm{实体}{substance}可能不同，然而通过这两个实体，凭借各自正当的能力，并维持它们彼此之间适当的关系，人这一位格便得以成立。我们当然不是说，灵魂与身体是同一\OriginalTerm{实体}{substance}，而是肯定它们各有自己的特性；正如这比喻中皮囊和酒应用于一类、一种人类，同样，真理的推论要求我们承认，人是由唯一的天主所完全造成的：因为灵魂在身体中喜乐，并爱它、珍视它；而身体也同样因被灵魂赋予活力而喜乐。然而，另一方面，如果有人主张身体是恶者的作品，因为它如此易朽、陈旧、无价值，那么接着就会推出，它无法承受精神的德能或灵魂的活动，以及那同一者的最辉煌的创造。因为，就像人将一片新布补在旧衣服上，破绽反而更糟糕；同样，身体若在这样的处境下，与灵魂那最为灿烂的作品结合在一起，也必会朽坏。或者，再用另一个比喻：就像一个人将灯的光带入黑暗之地，黑暗立刻被驱散，不留踪影；我们同样应当理解，灵魂一进入身体，黑暗便被立即驱逐，一个单一的本性随之实现，一个人也按一个种类构成。因此，与之相合，便应承认新酒被装入新皮囊，新布不补在旧衣服上。由此我们能够表明，在这两个\OriginalTerm{实体}{substance}之中，即身体的实体和灵魂的实体之中，存在着能力的合一；对于这种合一，圣经中最伟大的导师保禄曾谈到，他告诉我们：\BibleQuote{但如今天主却按自己的旨意，把肢体一一都安排在身体上了。}\par

19. 但若你们理解这一道理似乎有困难，并且你们不接受这些说法，我至少可以试着通过提供比喻来使它们成立。请把人类比作一种殿宇，依照\OriginalTerm{圣经}{Scripture}的比喻：在人内的灵，因此可以比作居住在殿中的像。那么，一个殿宇若未先确认为殿宇的承租者，便不能建立；另一方面，承租者若没有建起殿宇的结构，也不能安置在殿宇中。如今，既然这两者，即承租者与结构，是同时被祝圣的，它们之间如何能找到任何敌对或相反？它们岂不更应当显得都是和谐一心之主体的产物吗？为使你们知道事实如此，这些主体在团契和血缘上确实是合一的，那知晓并听见\Emph{一切}的那位作出了这样的回应：\BibleQuote{让我们造人}，等等。因为建造殿宇的询问塑造像的，并仔细询问有关大小、宽度和体积的尺寸，为要按照这些尺寸标出地基的空间；并且没有人会不先使自己熟悉安置像所需的尺寸，便着手徒劳地建造殿宇。同样，因此，身体的形式和尺寸也成为询问的对象，为使灵魂能适当地被万有之造作者天主安置其中。但若有人说塑造身体的那位与创造我灵魂的天主为敌，那么，为何这两方以敌意相视，却没有败坏这工程，要么使建造殿宇的将殿宇造得过小，以至于无法容纳安置其中的东西；要么使塑造像的带来如此巨大沉重之物，以致将其引入殿宇时，建筑物立刻倒塌呢？倘若并非如此，那么，根据这些，让我们按照我们所知的敌对者的目标和意图来思考它们。但若一切都应依同样的尺寸和同样的公平来安排，并以同样的荣耀展现，我们对此主题还当存何疑虑呢？若蒙允许，我们再增加一个比喻。人好似一艘由造船者建造并投入深海中的船，然而没有舵便无法航行，借由舵才能保持指挥，并转向舵手想要航行的任何方向。再者，舵与船的整个船身需要同一造作者，这是无可置疑的事；因为没有舵，船的整个结构——那巨大的身体——便只是一块不动的物体。因此，我们说灵魂是身体的舵；而且这两者都受我们所拥有的、相当于舵手的判断和情感的自由所支配；当这两者借由结合而合一，从而拥有一种适用于各样工作的功用和谐一时，其自身运作所产生的一切，都见证它们有着同一的创始者和制作者。\par

20. 听到这些论证，在场的群众极为喜悦；甚至几乎要动手拿住\Person{摩尼}；\Person{阿基莱乌斯}好不容易才制止住他们，将他们拦下，使他们恢复安静。法官们说：\Quote{阿基莱乌斯已向我们充分证明了人的身体和灵魂是同一手的作品这一事实；因为一个事物不能以作为一手的作品之恰当和谐与统一而存在，如果设计计划中有任何不和谐的话。但如果宣称一个人不可能足以发展这两样，\Emph{即身体和灵魂}，这不过是暴露制作者的无能。因为这样，即便有人承认灵魂是一位善神的造物，就人而言，它亦将只成为一件无用功，除非它也将身体归属自己。而如果反过来，身体被认为是一位恶神的造物，这工程亦同样是无用的，除非它接受灵魂；并且，事实上，除非灵魂借由混合与适当的引入而与身体和谐一致，使两者相互连接，人便不会存在，我们也不能谈论他。因此，我们认为，阿基莱乌斯已借众多比喻证明了整个人只有同一位造作者。}\Person{阿基莱乌斯}说：摩尼，我不怀疑你明白这一点，即，出生和被造的人被称为生他或造他者的儿子。但如果那邪恶者造了人，那么按照本性，他就应当是人的父亲。那么，主耶稣是用以下这些话教导人祈祷时，祂是对谁说的呢：\BibleQuote{当你们祈祷时，应当说：\InnerQuote{我们在天的父}；} 还有，\BibleQuote{应当向你们那在隐秘中的父祈祷？} 但祂说\BibleQuote{祂看见撒殚像闪电一样从天跌下}，是指撒殚说的；所以没有人敢说祂教我们向他祈祷。当然耶稣并非怀着聚集人、使人同撒殚和好的目的从天降下；相反，祂将牠交付，使其被祂忠信者们的脚下踏碎。然而，就我而言，我要说，那些外邦人反倒更为有福，他们确实引进众多神明，但至少认为这些神明全都同心合意，彼此和睦；而这个人，虽然只引进两位神，却不知羞耻地在他们之间设立仇敌和不合的情感。而且，事实上，如果那些外邦人以那种情状引进他们虚假的神明，我们确实将有机会目睹一场类似于他们之间进行的角斗比赛，带着他们无数的性情和分歧的情感。\par

21. 但现在，关于内里的人和外的人，我必须说的，可以用救主对那些吞下骆驼、穿戴伪君子外袍、满口阿谀奉承之人所说的话来表达。\Person{耶稣}正是对这些人在说话，祂说：\BibleQuote{祸哉，你们经师和法利塞假善人！因为你们洗擦杯盘的外面，里面却满是贪欲和污秽。你们不知道那造外面的，也造了里面吗？} 那么，祂为什么提到杯和盘呢？说这些话的人难道是玻璃匠，或制陶器的窑匠吗？祂岂非最明显地在论身体和灵魂吗？因为法利塞人确实注重\BibleQuote{捐献十分之一的茴香和莳萝，却放过了法律上最重要的公义、仁爱与信义}；他们对外在之事极其用心，却忽略了关乎灵魂救恩的事。因为他们也看重\BibleQuote{街市上的致敬}，和\BibleQuote{筵席上的上座}；主\Person{耶稣}知道他们将遭丧亡，就宣告说，他们只注意外在的事，却鄙视内在的事，视之为陌生之物，且不明白造身体的，也造了灵魂。有谁如此冥顽不灵、心智愚钝，竟看不出\Emph{我们主}的那些话足以应对一切情况呢？再者，\Person{保禄}的话与这些话完全和谐，当他解释律法上所写的一些话时，他说：\BibleQuote{牛在打场的时候，不可笼住它的嘴。难道天主所顾念的是牛吗？岂不是全为我们说的吗？} 但我们何必在这主题上再浪费更多时间呢？不过，我还要从许多可提供的东西中略加数点。假定现在有两个\OriginalTerm{非受生的本原}{unbegotten principles}，且我们为它们指定固定的场所：那么，如果天主\Emph{被假定}在某个地点之内，而不是无处不在，祂就因此被分割了；结果，祂\Emph{会被描绘}成远低于祂所处的那个地点，\Emph{因为包含者总是大于被包含者}：这样，天主就被造成与包含祂的那个地点的大小相对应，就像一个人在屋子里那样。再者，理性会问，是谁将它们分开，或为它们划定了界限；于是，两者都将被证明远远不及人自己的力量。因为\Person{利西马科斯}和\Person{亚历山大}曾掌握全世界的帝国，能征服一切外邦和整个人类；以致在整个那段时期，天下除了他们以外别无其他掌权者。有谁会冒昧地说，天主——祂是真光，永不受遮蔽，祂的国是圣洁的、永存的——不是无处不在的呢？这乃是那个最邪恶的人的方式，他在其不虔敬中，甚至拒绝将与人相等的能力归于全能的天主！\par

22. \Emph{法官们说：}我们知道，光照亮的是整个房屋，而非仅仅某一部分；正如耶稣所言：\BibleQuote{人点灯，并不是放在斗底下，而是放在灯台上，照耀屋中所有的人。}那么，如果天主是光，则此光（若耶稣可信）必须照耀整个世界，而非仅仅部分。若此光占据整个世界，那么何处还能有非受造的黑暗？或者，除非黑暗仅仅是某种偶然之物，否则怎能说黑暗存在呢？ \Emph{\Person{阿基劳斯}说：}诚然，因为你们对福音之言的理解，远胜于这位自称为\OriginalTerm{护慰者}{Paraclete}的人——尽管我宁愿称他为\OriginalTerm{寄生虫}{parasite}而非护慰者——我要告诉你们黑暗是如何出现的。当光已遍布各处，天主开始构造宇宙，由天地起始；在此过程中，出现了一个问题，即中部——就是大地所在、被阴影覆盖之处——由于受造物被造成时的介入，变得幽暗，因而需要将光引入这位于中部的地方。因此，在\BibleBook{}{创世纪}中，\Person{梅瑟}记述世界的构造时，他既未提及黑暗是被造的，也未提及它是非受造的。但他对此事保持沉默，留待那些能恰当注意的人去发现其解释。其实，这并非十分艰巨困难的任务。谁不能明白：我们的太阳在东方升起，向西方运行，是可见的；但当它行至地下，深入希腊人称之为\OriginalTerm{球体}{sphere}的构造中时，便因物体阻隔而被黑暗笼罩，于是不再显现？当太阳如此被遮蔽，大地之体正对其时，便生出一层阴影，从而产生黑暗；如此持续，直至夜间行经下方空间后，它又滚向东方，再次从惯常之处升起。因此，阴影和黑夜的原因就在于地球的坚实实体——事实上，人从自身身体投下的影子这个事实就能明白。因为在天地及所有形体的受造物出现之前，光一直恒常不变，没有亏缺或蚀变，因为没有任何物体能通过相对或介入而产生阴影；因此必须说，那时无处有黑暗，也无处有黑夜。举例来说，倘若那掌握万有者乐意除去西方区域，那么，因太阳不再向那方运行，便无处出现傍晚或黑暗；太阳将永在途中，永不下沉，几乎总是占据天宇的中部，永不停止显现；借此，整个世界将被最清澈的光照亮，任何部分都不会被阴影遮蔽，同一之光的力量均匀遍及各处。但另一方面，当西方区域保持其位置，太阳在世界的三部分运行其轨道时，那些处于太阳之下的人会被更明亮地照耀；我几乎可以说，属于不同地区的人还在睡梦中时，前者已拥有了白日的开始。但正如东方人比居住西方的人更早迎来光明，他们也更早被黑暗笼罩；唯有定居于地球中部的人，总是见到均等的光。因为当太阳占据天宇正中时，没有任何地方显得比别处更亮或更暗，世界各处都平等、无私地被太阳的光辉照亮。那么，如我们上文所言，若西方区域被除去，与之相邻的部分便不再受遮蔽。我本来可以更简明地阐述这些事，正如我也可描述\OriginalTerm{黄道带}{zodiacal circle}一样；但我目前不拟深入探讨。因此，我将不谈论这些，而是回到对手提出的主要反对意见，即他极力声称黑暗是非受造的；然而，这一立场已在我们力所能及的范围内被驳斥了。\par

23. \Emph{法官们说：}如果我们考虑到光在受造物进入之前就已存在，并且没有任何处于相对位置能够产生阴影的物体，那么必然得出结论：光那时遍布各处，一切地方都被其光辉照亮，正如你刚才陈述所表明的；我们看出那里面给出了真正的解释，因此我们将胜利归于\Person{阿尔克劳斯}的肯定。因为如果宇宙被清楚地划分，仿佛有一道墙穿过其中央，如果光居于一侧，黑暗居于另一侧，那么必须理解，这黑暗是因世上所立物体产生的阴影而偶然发生的；因此我们必须再问：是谁在这两个区划之间建造了这道墙？哦，\Person{摩尼}，倘若你确实承认有这样一道墙存在。但是，如果我们必须在假设没有建造这样一道墙的前提下考虑此事，那么再次得到理解：宇宙仅仅形成一个区域，毫无例外，并置于同一权下；若如此，黑暗便绝不能具有\OriginalTerm{非受生本性}{ungenerated nature}。\Emph{阿尔克劳斯说：}让他也解释以下议题，以关照所提出的论点。如果天主坐在自己的国度里，而恶者同样坐在自己的国度里，那么谁能建造他们之间的墙呢？因为没有任何事物能分开两个\OriginalTerm{实体}{substance}，除非是大于两者之一的事物，正如在\BibleBook{创}{世纪}中所说的：\BibleQuote{天主将光与黑暗分开。}因此，这道墙的建设者也必须具有类似的能力：因为墙标志着双方的界限，正如在乡村居住的人中，人们通常用一块石头来划定各家的地界；然而，这一习俗，若我们将这种划分特别视为对兄弟之间产业分配的标记，或许能更好地理解此事。但就目前而言，尽管这些事情可能看起来很关键，我还不需要谈论它们。因为我们所寻求的是对问题的回答：谁能建造那条为标示这两者各自国度的界限所必需的墙？没有人给出答案。让这个背信的家伙不要犹豫，现在承认他的\OriginalTerm{二元}{duality}实体已被再次归结为\OriginalTerm{统一}{unity}。让他指出任何能建造那道隔墙的人。当一方在建造时，这两方中的另一方可能在忙什么呢？他在睡觉吗？还是他对此事一无所知？还是他无法抵挡这一企图？还是他被收买了？告诉我们他当时在做什么，或者告诉我们全宇宙中谁是建造那工程的人。我向你们呼吁，哦，法官们，天主以最充分的智慧将你们派到我们这里；请判断，这两者中谁能建造这构造，或者当一方在建造时，另一方一直在做什么。\par

24. \Emph{审判官们说：}请告诉我们，\Person{摩尼}，是谁划定了各自王国的疆界，又是谁筑了中间的墙？因为\Person{阿基劳斯}恳请在本次讨论中重视问答的实践。\Emph{摩尼说：}那至善、与恶毫无共通之处的天主，将\OriginalTerm{穹苍}{firmament}放置在中间，为要显明那恶者是外于祂的。\Emph{阿基劳斯说：}你何等可怖地污蔑了那名字的尊严！你确实称呼祂为天主，但只是名义上如此，你使祂的神性类似人的软弱。你时而说祂从虚无中，时而又说祂从基础物质中——其实这物质在祂以先就已存在——建造了这架构，如同人间的工匠惯常所做的那样。你有时又说祂是忧虑的，有时又说祂是易变的。然而，做合乎天主之事乃是天主的本分，做合乎人之事乃是人的本分。那么，如果天主如你所说建了一道墙，这位天主便自显为忧虑且毫无刚毅。因为我们知道，凡怀疑有陌生人暗中预备谋害自己，且惧怕敌人阴谋的人，惯常在城周围筑墙，借此他们既在无知中自保，又显露出其能力的薄弱。然而在此，也有我们不当沉默放过，而应加以强调的事；好叫我们凭着所持真理的帮助，借对此问题的大量陈述，使对手的诸多诡计归于无有。我们姑且承认，这墙的建造是为着区分两个王国而用的；因为若没有这一分隔，任何一个王国都无法拥有自己的专属国度。但若承认了这一点，那么进一步可得，同样地，那恶者也不可能越过自己的界限去侵入那良善\Emph{君王}的疆域，因为墙立在那里作为障碍，除非它先被推倒；因为我们曾听闻敌人做过这样的事，而且我们确实亲手亲眼，就在近期看见这等事成功实行。当一位君王攻打一座被坚固城墙环绕的堡垒时，他首先使用\OriginalTerm{弩炮}{ballista}和投射物；然后试图用斧头劈开城门，用\OriginalTerm{攻城锤}{battering-rams}摧毁城墙；当他最终得以进入，占领那地方，他便随心所欲而行：或乐意将居民掳去为奴，或将堡垒及其内一切全然毁灭；又或者，在另一方面，他也可愿意因被征服者的谦卑恳求而赦免那被攻取的坚固城。那么，我的对手对此比喻作何言说？难道就没有什么对手实质地——也就是故意地——推倒了那筑在两方之间的\OriginalTerm{壁垒}{muniment}吗？因为在他之前的陈述中，他曾断言那黑暗越过了自己的界限，并侵入那善天主的国度。那么，在那一位能够如此越界之前，是谁推倒了那壁垒呢？因为当壁垒坚立时，恶者是不可能找到任何入口的。你为何沉默？你为何犹豫，\Person{摩尼}？然而，即使你退缩，我也将自行继续这任务。因为若我们假设你说天主摧毁了它，我就必须问：是什么促使祂如此摧毁了祂自己先前为着那恶者的逞强、且为着保持两者之间的分隔而建造的这东西？是在怎样的情绪发作，或是感到怎样的伤害之下，祂才这样做来与祂自己争战？还是祂贪图那恶者的某些产业？但若这些都不是导致天主摧毁祂很久以前为使那恶者与祂隔绝而建造之物的真正原因，那么即使天主也开始喜爱与那恶者相交，也就不必被视为奇事了；因为，按你的假设，那为保守天主免受其扰而设立的壁垒，如今竟好像被除去，只因如今那恶者不再被视为仇敌，而是友人。而另一方面，你若断言那墙是被恶者摧毁的，那么请告诉我们，善天主的工程如何竟能被恶者所胜。因为若那是可能的，那么恶的本性将被证明比天主更强。此外，那纯全的黑暗本身，既见那光照耀，怎能突袭并攫取那光呢？而福音作者却向我们作证说：\BibleQuote{光在黑暗中照耀，黑暗却没有胜过它吗？}这盲目的如何能武装？黑暗怎能与光明的王国争战？因为正如天主的造物在此无法以健全的眼承受太阳的光线，那存在也不能承受光明国度的清晰异象，而是永远对之为陌生的、局外的。\par

25. \Emph{\Person{摩尼}说：}并非所有人都领受天主的话，只有那些蒙赐能认识天国奥秘的人才领受。即使现在，我也知道谁是属我们的；因为祂说：\BibleQuote{我的羊，听我的声音。}为了那些属于我们的，以及那些蒙赐领会真理的人，我要用比喻说话。恶者好比一头狮子，企图偷偷袭击善牧的羊群；牧人看见后，就挖了一个大坑，从羊群中取出一只小山羊丢进坑里。狮子渴望得到它，情欲勃发，要吞吃它，就奔到坑前，掉进了坑里，发现自己无力再爬出来。于是牧人捉住它，小心关在洞中，同时确保了那只与它同在坑中的小山羊的安全。恶者就这样被削弱了——可以这样说，狮子不再有能力做任何有害的事了；这样，所有灵魂的族类都将得救，那曾经失落的，终将归回他本有的羊群。\Emph{\Person{阿尔赫劳斯}说：}如果你把恶者比作狮子，把天主比作真正的牧人，请告诉我们，我们该把绵羊和小山羊比作什么？\Emph{\Person{摩尼}说：}我看绵羊和小山羊本性相同：它们被用作灵魂的象征。\Emph{\Person{阿尔赫劳斯}说：}好，那么，当天主把那灵魂放在坑中的狮子面前时，祂就是把它交给了丧亡。\Person{摩尼}说：绝非如此；远非如此。但祂是出于某种安排，将来祂会拯救那另一个——\Emph{那灵魂}。\Emph{\Person{阿尔赫劳斯}说：}各位听众，如果一个牧人害怕狮子来袭，却把他素常抱在怀里的羊羔暴露给那野兽吞食的烈怒，然后又说日后他要拯救那小东西，这岂不是极为荒谬的事？这难道不是极其可笑吗？是的，这毫无道理。\Emph{因为在你的比喻所隐含的假设中}，天主这样把一个灵魂交给撒殚，好让他捉去毁灭。但牧人什么时候做过那样的事？\Person{达味}岂不是从狮子或熊的口中救出羊来吗？我们提到这一点，正是因为那\Emph{从狮子口中}的说法；因为按照你的理论，那将意味着牧人能够从狮子口中，甚至从它肚子里，取出它已经吞下的东西。但也许你会回答说，我们谈论的是天主，祂凡事都能。但请听我对此的回应：那你为什么不更主张祂真实的能力，肯定祂能凭自己的大能、靠天主纯全的力量，而不需任何诡计，也不需把一只小山羊或羊羔丢进坑里，就能制伏狮子呢？也请告诉我这一点：如果设想狮子在牧人没有羊的时候来袭，后果会怎样？因为这里所谓的牧人被假定为\OriginalTerm{非受生的}{unbegotten}，这里所谓的狮子也是\OriginalTerm{非受生的}{unbegotten}。因此，当人还未存在时——换句话说，在牧人拥有羊群之前——如果狮子在那时来袭牧人，那会有什么后果，既然在小山羊存在之前，狮子根本无物可吃？\Emph{\Person{摩尼}说：}狮子确实无物可吞，但当他奔驰在山岭峰巅时，他仍对所能遇见的一切施展他的恶行；若有时他需要食物，就捕拿一些属他国中的野兽。\Emph{\Person{阿尔赫劳斯}说：}那么，这两类——恶者国中的野兽，与善天主国中的小山羊——是同一\OriginalTerm{实体}{substance}吗？\Emph{\Person{摩尼}说：}远非如此；绝非如此：它们本身以及各自相关的属性都没有任何共同之处。\Emph{\Person{阿尔赫劳斯}说：}狮子吃食物时，用途只有一种，且是相同的。尽管它有时从属自己的野兽取食，有时从属善天主的野兽取食，但就所供的肉食而言，它们之间并无差别；由此显明，那些是同一\OriginalTerm{实体}{substance}。另一方面，如果我们说二者之间有很大的差别，我们就是把无知归给了牧人，因为他没有把适合狮子吃的食物放在它面前，反而放的是异类的肉。或者，也许你为了掩饰自己的真实立场，会对我说那狮子什么也没吃。好吧，假设如此，那么天主就以这种方式，叫那个根本不知如何吞吃的活物去吞吃一个灵魂吗？而且，天主岂不是只靠那个坑来欺骗他吗？——如果这确实配得天主去做，或设计骗局的话。这就像一个君王，当有人向他开战时，他完全不信任自己的力量，反而因自己软弱的恐惧而瘫痪，把自己关在城墙内，四围筑起壁垒和其他防御工事，装备齐全，却不依靠自己的手臂和勇力；相反，如果他是一个勇士，那么处于这地位的君王就会从自己的领土长途跋涉，到那里迎敌，奋力作战，直到征服并擒获他的对手。\par

26. \Emph{法官们说}：如果你声称牧人将山羊羔或绵羊羔暴露给狮子，而那头狮子正谋划攻击那不受生者，那么案件已经了结。因为显然，那牧羊人和羔羊的牧者自己被证实对它们犯了过错，那么他还能对哪一受造物进行审判呢？如果那只因牧人软弱而被交出的绵羊羔未能抵抗狮子，并且因此它不得不做狮子所乐意做的任何事情，又会如何？或者，另举一例，这就像一个主人由于恐惧而将一名奴隶赶出家门，或交给对手，而他后来却无法凭自己的力量将其收回。又或者，假设奴隶碰巧被收回了，主人又有什么合理的理由对他施以酷刑呢？如果事实证明那个奴隶对敌人强加给他的一切都俯首顺从，因为正是主人亲手将他交给敌人，就像山羊羔被交给狮子一样。你还断言，牧人预先就明白了全部情况。那么，当然，当羊羔受到鞭打，并被牧人质问为何在这些事情上服从了狮子时，它大概会这样回答：\Quote{是你亲手把我交给了狮子，你没有对他进行任何抵抗，尽管你深知并预见到我的命运如何，当我不得不屈从于他的命令时。} 而且，我们不必在这上面过多铺陈，可以说，\Emph{通过这个比喻}，既没有显示出天主是一位完美的牧人，也没有表明狮子尝到了不属于它的肉；因此，在真理本身的教导下，显然我们应该把优胜给予阿尔凯劳斯所提出的推理。\Emph{阿尔凯劳斯说}：考虑到在所有我们至今讨论的各点上，法官们的深思熟虑给我们提供了最充分的余地，我们最好对其他问题保持沉默，留待日后讨论。因为正如，若有人一旦踩碎蛇的头，便无需再砍断它身体的其他任何部分；所以，一旦我们解决了这个二元论的问题——我们已尽力而为——与之相关的其他论点也随之一并瓦解。不过，我仍要对这位持这种观点的人，也就是现在在我们面前的这位，说几句话，哪怕是几句，以便让所有人都彻底了解他是谁，来自何处，以及他表现出的究竟是何等人物。因为他自称是那位\OriginalTerm{护慰者}{Paraclete}，即耶稣在离别时承诺要派遣给人类，以拯救信众灵魂的那一位；他做出这种自称，就好像他甚至比保禄还更胜一筹，而保禄却是\OriginalTerm{被拣选的器皿}{elect vessel}、\OriginalTerm{蒙召的宗徒}{called apostle}，他正因此，在宣讲真道时曾说：\BibleQuote{你们岂是在寻求基督在我内说话的凭证吗？} 然而，我要说的可以通过以下比喻变得更加清楚：——有一个人将极大量的谷物收进他的仓房里，以至于那里完全堆满了。他令人满意地关闭了仓房并贴上封条，吩咐人严加看守。然后主人自己便离开了。然而，过了很长一段时间，另一个人来到仓房，声称他是那个锁闭并密封此处的人所派遣的，受命也要在同一处收集和存放一些麦子。看守仓房的人看见他，便向他索要凭证，即出示印鉴，好让他们确定是否有权为他打开仓房并服从他，如同服从那位密封此处者所派遣的人一样。由于他既拿不出钥匙，也不能出示印鉴作为凭证，\Emph{因为实际上他根本无权}，他便被看守人推出去，被迫逃走了。因为他非但不是他所自称的那类人，反而被他们识破是个盗贼和劫匪，并且通过以下事实被定罪和揭穿：尽管他看起来是要在那预先确定了的时间之后许久才出现，然而他既不能向看守人出示钥匙、印鉴或任何信物，也不能显示出对仓中储粮数量有任何了解——所有这些情况都是确凿的证据，证明他并非由真正的物主所差遣；于是，理所当然地，他被看守人禁止进入。\par

27. 如果你认为好的话，我们可以再举一个例子。有一个人，是一位家主，拥有大量财富。他打算外出旅行一段时间，并向他的儿子们许诺，将派一个人来代替他，将他们应得的财产均分给他们。事实上，不久之后，他确实派了一个可靠、正直且诚实的人给他们。那人到达后，接管了全部财产。他首先努力整理并管理这些财产，自己辛苦奔波，甚至亲手勤劳工作，像仆人一样为产业的利益辛劳。后来，那人感到自己死期将近，便立下遗嘱，将继承权转让给亲属及所有近亲；并将他的印章交给他们，逐一按名将他们召聚起来，嘱咐他们保存遗产，照料财产，并按照他们所领受的那样，正确地管理，享用其中的货财及果实，如同他们自己就是它们的主人和继承人一样。而且，如果有人请求允许享用这片田地的果实，他们应当对他宽容。但另一方面，如果有人宣称自己与他们同享继承权，并以此为由提出要求，他们就要远离他，宣布他为外人；而且\Emph{他们应当认为}，想要被他们接纳的人，正因如此更当工作。那么，假定所有这些事都妥善且正确地安排和解决了，并且这种状况延续了非常长的时间之后，对于这样一个在三百年后出现，并提出继承要求的人，我们该如何对待呢？我们难道不把他赶走吗？我们岂不应当公正地宣布这样的人为外人——他无法证明自己属于\Emph{与我们的主有关系的人}，他从未在我们离世的主病危时刻陪伴在旁，从未走在被钉十字架者的送葬行列中，从未站在坟墓旁，对祂离开的方式或性质一无所知，总之，如今却渴望进入粮仓，却无法出示那位将它上锁加封者所给的任何凭证？我们难道不像驱赶强盗和小偷一样将他赶走，并用尽一切方法将他逐出我们中间吗？然而，这个人如今就在我们面前，却未能出示任何我们已概述过的凭证，却宣称他就是耶稣预言将要派遣的\OriginalTerm{护慰者}{Paraclete}。因着这个宣称，或许出于无知，他将使耶稣自己成为说谎者；因为那曾说过不久之后就要派遣护慰者的那一位，如果我们在三百年多之后，接受他为自己所作的见证，就会被证实仅仅派遣了这个人。那么，在审判的日子，那些从那时期直到现在离世的人，要对耶稣说什么呢？他们岂不是要用这样的话质问祂吗：\Quote{如果我们没有做祢的工，求祢不要严厉惩罚我们。因为祢曾应许要在\Person{提庇留·凯撒}在位时派遣护慰者，为叫我们知罪又知义，但为何祢竟只在罗马皇帝\Person{普罗布斯}在位时才派遣他，使我们成为孤儿，尽管祢说过：\BibleQuote{我必不留下你们为孤儿}\BibleRef{若 14:18}，而后祢又明确保证祢离去之后，立刻就派遣护慰者？我们这些孤儿能做什么呢，既没有监护人？我们并没有过错；是祢欺骗了我们。}我们万不可对每个灵魂的救主、我们的主耶稣基督作如是想。因为祂并不限于空许；当祂曾说：\BibleQuote{我去到我父那里，并给你们派遣护慰者，}\BibleRef{若 14:12,16}祂便立即赐下那恩赐，即护慰者，分施给祂的门徒们——却更丰盛地赐予了\Person{保禄}。\par

28. \Emph{摩尼说：}你已被你自己提出的指控所困住。因为你一直在反对你自己，却没有意识到，当你试图将羞辱抛向我时，你却使自己陷入更大的过错。我请你告诉我这一点：如果如你所声称的那样，从提庇留时代直到普罗布斯时代去世的人会对耶稣说：\Quote{不要审判我们没有行你的工程，因为你没有派遣\OriginalTerm{护慰者}{Paraclete}给我们，尽管你应许要派遣祂；}那些从梅瑟时代直到基督亲自来临期间离世的人，岂不更会如此述说吗？那些已离世的人岂不更有理由用这样的话表达自己：\Quote{不要把我们交于折磨，因为我们没有得到任何关于你的知识？}而且，难道只有那些在祂来临之前死去的人才有理由提出这样的指控吗？那些从亚当时代直到基督来临之间去世的人，难道不也是这样吗？因为他们中没有任何人得到关于\OriginalTerm{护慰者}{Paraclete}的知识，也没有接受耶稣教义的教导。但只有这最后一代人，按你的说法，从提庇留以来的人，才能得救：因为正是基督\Quote{由法律的咒骂中赎出了他们；}\BibleRef{迦 3:13} 正如保禄也给出了这些进一步的见证：\Quote{文字叫人死，而不叫人活，}\BibleRef{格后 3:6} 以及\Quote{法律是死亡的职务，}\BibleRef{格后 3:7} 和\Quote{罪的势力。}\BibleRef{格前 15:56} \Emph{阿尔科劳说：}你错了，既不明白圣经，也不明白天主的能力。因为在基督来临之后直到如今这一时期，许多人也已丧亡，如今仍有许多人丧亡——就是那些不致力于义德工作的人；而只有那些接受祂，并且仍在接受祂的人，\Quote{获得了成为天主子女的权能。}\BibleRef{若 1:12} 因为圣史并没有说所有人都获得了那权能；然而另一方面，他也没有限定时间。但他的表述是：\Quote{凡接受他的。}\BibleRef{若 1:12} 再者，从世界的创造以来，祂一直与义人同在，并从不停止\Emph{从恶人手中}追讨他们的血，从义人亚伯尔的血直到匝加利亚的血。那么，当还没有梅瑟法律，也还没有兴起先知履行预言职责之时，义人亚伯尔以及所有那些被列在义人之中的后继贤者，他们的义德从何而来呢？岂不是因为他们践行了法律，\Quote{人人都显示出法律的精华已刻在他们心上，他们的良心也为此作证吗？}\BibleRef{罗 2:15} 因为当一个人\Quote{没有法律的，顺着本性去行法律上的事，他虽然没有法律，自己对自己就是法律。}\BibleRef{罗 2:14} 现在你想一想，在那些过着正直生活的义人们中间存在的大量法律，他们有时发现刻在自己心中的天主法律，有时从父母那里学得，有时又通过古人和长老们进一步受教。但是，由于为数不多的人能通过这种方式达到义德的高峰，也就是说，靠父母的传统，那时还没有成文的法律，天主便怜悯了人类，乐意藉梅瑟赐给人们一部成文法律，因为自然法律的公平确实未能完美地保留在他们心中。因此，与人起初的受造一致，一部成文的立法被预备好，通过梅瑟赐予，为了许多人的得救。因为如果我们推想人成义不靠法律的行为，如果亚巴郎被算为义人，那么那些满全了包含有益于人之事的法律的人，岂不更将获得义德吗？既然你只提及了那三段经文，就是宗徒所宣称的：\Quote{法律是死亡的职务，}\BibleRef{格后 3:7} 和\Quote{基督由法律的咒骂中赎出了我们，}\BibleRef{迦 3:13} 以及\Quote{法律是罪的势力，}\BibleRef{格前 15:56} 你现在尽管再提出其他类似意思的章节，并尽量找出那些看似反对法律的经文吧。\par

29. \Emph{摩尼说}：耶稣对门徒讲论那些人是不信者时，所说的那句话，不也有同样的意思吗？\BibleQuote{你们是出于你们的父亲魔鬼，并愿意追随你们父亲的欲望。}\BibleRef{若 8:44}祂的意思其实是，这邪恶的世界首领所渴望的、所贪求的一切，他都通过梅瑟记录下来，并藉此交给人们去行。因为\BibleQuote{他从起初就是杀人犯，不站在真理上，因为在他内没有真理；他说谎话，是出于他自己，因为他是说谎者，也是说谎者的父亲。}\BibleRef{若 8:44} \Emph{阿赫劳斯说}：你对已经引证的内容满意吗？还是另有其他陈述要说？\Emph{摩尼说}：确实，我还有许多话要说，而且比这些更重要。但我就满足于这些了。\Emph{阿赫劳斯说}：当然可以。现在让我们从你自称支持你的陈述中选出一个例子；这样，如果发现这些陈述得到妥当处理，其他问题也可与之并列；如果情况相反，我就接受法官们的定罪，也就是说，我将蒙受失败的耻辱。 那么，你说法律是属死的职务，并且你承认死亡，即这世界的首领，从亚当起，直到梅瑟，作了王；因为圣经的话是这样说的：\BibleQuote{连那些没有犯罪的人也要受它管辖。}\BibleRef{罗 5:14} \Emph{摩尼说}：毫无疑问，死亡确实如此作王，因为存在二元对立，这两个敌对的力量无非都是\OriginalTerm{非受生的}{unbegotten}。 \Emph{阿赫劳斯说}：那么请告诉我这一点——一个非受生的死亡怎能有一个特定时间的开始呢？因为圣经说的是\BibleQuote{从亚当}，而不是\BibleQuote{在亚当之前}。 \Emph{摩尼说}：但我反过来问你，死亡是如何在义人和罪人身上都取得了王权的呢？\Emph{阿赫劳斯说}：等你先承认它从某个确定的时间而非永恒拥有那王权，我就告诉你。摩尼说：经上记载，\BibleQuote{死亡从亚当到梅瑟作了王。}\BibleRef{罗 5:14} \Emph{阿赫劳斯说}：因此它有终结，因为它在时间中有开始。而那句话说也是真实的：\BibleQuote{死亡在胜利中被吞灭了。}\BibleRef{格前 15:54}那么显而易见，死亡不可能是非受生的，因为它被证明既有开始也有终结。 \Emph{摩尼说}：但那样岂不是说天主是它的创造者。\Emph{阿赫劳斯说}：绝非如此；不可有这种想法！因为经上说：\BibleQuote{天主并未造死亡，也不乐意生灵灭亡。}\BibleRef{智 1:13} \Emph{摩尼说}：天主没有造它，但正如你承认的，它毕竟造成了。那么请告诉我们，它是从谁获得统治权，或是被谁创造的。 \Emph{阿赫劳斯说}：如果我充分证明死亡不可能具有非受生性的实体，你是否会承认只有一个天主，且是非受生的天主？ \Emph{摩尼说}：你继续讲吧，因为你的目的是要用巧言辩驳。\Emph{阿赫劳斯说}：不，是你把那些主张提出来的，仿佛它们可以为你证明一个非受生的根源。然而我们上面讨论的论点可能已经足够了，因为我们藉此充分表明，两个非受生性体的实体不可能并存。\par

30. \Emph{法官们说}：阿赫劳斯，请就他现在提出的那些论点发言。\Emph{阿赫劳斯说}：他用\Quote{世界的首领}、\Quote{那邪恶者}、\Quote{黑暗}和\Quote{死亡}来指同一个东西，并声称法律是由那位赐下的，理由是圣经称之为\BibleQuote{属死的职务}\BibleRef{格后 3:7}，以及他用来反对法律的其他一些理由。那么，我说，正如我们上面所解释的，那原本刻在人心的法律未能小心保持对邪恶事物的记忆，而且长老们之间没有足够稳固的传统，因为敌对的遗忘总是伴随着记忆，并且一个人由老师传授\Emph{对那法律的知识}，另一个人则靠自己，于是自然铭刻的法律的过犯就轻易发生了，并且由于违反诫命，死亡在人间取得了王权。因为人类的本性如此，需要用铁杖由天主来管辖。因此死亡得胜了，并以全部权能从亚当直到梅瑟作了王，甚至管辖了那些没有犯罪的人，正如我们所解释的：确实管辖了罪人，因为他们是它适当的对象并臣服于它——像加音和犹达斯那样的人；但也管辖了义人，因为他们拒绝顺从死亡，反而抵抗它，抛弃了邪恶和贪欲——像那些从亚伯尔到匝加利亚不时出现的人；因此死亡一直延续，直到\Emph{梅瑟}时期，将类似的人都置于其下。\par

但梅瑟出现以后，将法律赐给以色列子民，使他们记起法律的一切要求，以及人在法律之下所当谨守遵行的一切；当他只将那些违犯法律的人交付死亡时，死亡便被切断，不能再统治所有的人；因为那时死亡只统治罪人，正如法律对它说的：\Quote{勿触犯遵守我诫命的人。}因此，梅瑟就这言论对死亡执行了这职务，同时他把所有其他违犯法律的人交付毁灭；因为梅瑟的到来，并非为了死亡不再在任何领域内统治，实际上在梅瑟之后，无数人确实仍被死亡的权势所辖制。而法律被称为\Quote{\OriginalTerm{死亡的职务}{ministration of death}}，是因为那时只有违犯法律的人受惩罚，而那些遵守法律，顺从并遵行法律中所记载的人却不然，正如亚伯尔所做的，而加音，他成了恶者的器皿，杀害了亚伯尔。然而，就是在这些事以后，死亡仍想破坏藉梅瑟所立的盟约，重新统治义人；为此，它确实攻击了众先知，杀害并用石头砸死那些由天主派遣的人，直到匝加利亚。但我的主耶稣，为维护梅瑟法律的正义，因死亡违犯盟约和整个那职务而发怒，屈尊就卑，以人的形体显现，为要报仇，不是为自己，而是为梅瑟，以及为他以后连续不断遭受死亡暴力压迫的人。然而，那恶者由于不明白这种\OriginalTerm{安排}{dispensation}的\Emph{意义}，进入了犹达斯，想藉那人的手杀害他，正如他先前曾杀害义人亚伯尔一样。但他进入犹达斯之后，后悔不已，便上吊自尽了；为此，经上也说：\BibleQuote{死亡！你的胜利在哪里？死亡！你的刺在哪里？}\BibleRef{格前 15:55}又说：\BibleQuote{死亡被胜利吞灭了。}\BibleRef{格前 15:54}因此，正是由于这个缘故，法律被称为\Quote{死亡的职务}，因为它把罪人和违法者交付死亡；但对那些遵守它的人，它保护他们免于死亡；这些人也靠着我们的主耶稣基督的助佑，建立在光荣中。\par

31. 也请听我对于所引用的另一句话要说的，即，\BibleQuote{基督把我们由法律的诅咒中赎出来}。我对这段经文的看法是，天主的那位杰出的仆人\Person{梅瑟}，将象征性的法律和真正的法律，交给了那些愿意有正确看见的人。举例来说，当天主在六天内创造了世界及其中的一切之后，祂在第七天停工休息了；我这样说，并不是要断言祂因疲乏而休息，而是因为祂完成了祂决意要创造的每一受造物。然而后来，\Emph{新法律却说}：\BibleQuote{我父到现在一直工作，我也工作。}那么，这是否意味着祂仍在创造天、太阳、人、动物、树木或诸如此类的东西呢？不；意思是，当这些可见的事物被完全造成之后，祂就停止了那种工作；然而，祂仍继续以一种内在的活动方式，在不可见的事物上工作，并拯救人。同样，立法者也希望我们每一个人都不断地献身于这种工作，就像天主自己一样；因此，他命令我们持续地从世俗的事务中休息，不从事任何世俗的工作；而这被称为我们的安息日。他在法律中还加上这一点，即不可做任何糊涂的事，而应谨慎行事，按照正义与公义引导我们的生活。这法律悬在人之上，对违犯它的人发出最严厉的诅咒。但因为它的子民也只不过是人，而且，正如我们中间也常发生的那样，争辩与伤害兴起，法律也立即以最严格的公正，使任何不义回归到作恶者自己的头上；因此，举例来说，如果一个穷人在安息日想捡一捆柴，他就被置于法律的诅咒之下，面临立即死亡的刑罚。因此，那些与埃及人一同长大的人，被法律的限制力量严厉地压迫，他们无法承受法律的刑罚和诅咒。但是，那永远救主，我们的主\Person{耶稣基督}，来了，他把那些人从这些法律的痛苦和诅咒中解救出来，赦免了他们的过犯。祂确实不像\Person{梅瑟}那样对待他们，\Person{梅瑟}将法律的严厉加诸于人，且对任何过犯都不予宽免；但祂宣告，如果有人受邻人伤害，他要宽恕他，不只一次，甚至不是两次或三次，也不只七次，而是七十个七次；但是，另一方面，如果在这之后，那冒犯者仍继续行这样的不义，那么，作为最后的方法，他就应该被置于\Person{梅瑟}法律之下，而且对于这样坚持作恶的人，即使已被宽恕七十个七次，也不应再给予宽恕。祂不仅赦免具有那种性质的过犯者，甚至赦免得罪人子的。但是，如果一个人这样对待圣神，祂就使他承受两个诅咒，即\Person{梅瑟}法律的诅咒和祂自己法律的诅咒；\Person{梅瑟}法律的诅咒诚然在今生，而祂自己法律的诅咒在审判之时：因为祂的话是这样：\BibleQuote{无论在今世或在来世，都不得赦免。}因此，有\Person{梅瑟}的法律，在今生不给予\Emph{这样的}人赦免；还有\Person{基督}的法律，在将来惩罚。所以，由此请注意祂如何坚固法律，不但不废除，反而成全。这样，祂就把他们从属于今生的法律的诅咒中赎了出来；\Quote{法律的诅咒}这个称呼即由此而来。这就是关于那种说话方式\Emph{需要给出的}全部解释。再者，为什么法律被称为罪过的权势，我们将尽我们所能，立即简短地解释。经上记着说：\BibleQuote{法律不是为义人立的，乃是为不法和不服的，不虔敬和犯罪的。}那么，在这些时代，在\Person{梅瑟}之前，没有为犯罪者设立的成文法律；因此，\Person{法郎}也不知道罪过的权势，以不义的重担压迫以色列子民的方式犯罪，并且藐视天主，不仅他自己，还有所有与他在一起的人。但是，为了不作绕弯的陈述，我将简要地如下解释。当\Person{梅瑟}的百姓在旷野中受他管理时，有一些埃及族裔的人混在\Person{梅瑟}的百姓中；当\Person{梅瑟}上到山上，为要领受法律时，那些没有耐性的人，我不是指真正的以色列人，而是那些与埃及人混在一起的人，按照他们崇拜偶像的古老习俗，树立了一头牛犊当作他们的神，以为用这样的方式，他们就可以永远不必为他们的不义承受应得的惩罚。这样，他们完全不知道他们罪过的权势。但当\Person{梅瑟}从山上回来，发现了那事，他就下令用刀剑处死那些人。从那时起，藉着\Person{梅瑟}法律的工具，这些人开始正确认识到罪过的权势；因此，法律被称为\Quote{罪过的权势}。\par

32. 此外，关于福音中所记载的这句话：\BibleQuote{你们是出于你们的父魔鬼}等等，我们简要地说，有一个魔鬼在我们内工作，其目的乃是凭着他自己意志的力量，使我们像他一样。因为天主所造的一切受造物，祂都造得甚好；祂赐给每一个体自由意志的感觉，按照这个标准，祂也设立了审判的法律。犯罪是我们的事，而我们不犯罪是天主的恩赐，因为我们的意志被构造为能选择犯罪或不犯罪。\Person{玛内斯}，这一点你自己无疑非常了解；因为你知道，即使你将你所有的门徒都聚集起来，劝诫他们不要犯任何过犯或行任何不义，他们每个人仍可能不顾审判的法律。的确，凡愿意的，都能遵守诫命；凡藐视诫命，转而违背它们的人，无疑仍必须面对这审判的法律。因此，某些天使拒绝服从天主的命令，反抗了祂的旨意；其中一个确实如同闪电一般坠落在地上；而其他的，被龙所困扰，在与人的女儿交合中寻求幸福，从而给自己带来了应得的永火之罚。至于那被扔到地上的天使，由于再也找不到进入天界任何区域的入口，如今便在人间四处招摇，欺骗他们，引诱他们成为像他一样的违背者，直到今日他仍是天主诫命的敌对者。然而，他堕落和毁灭的榜样，并非人人都会效法，因为每个人都被赐予了意志的自由。也正因如此，他获得了\Emph{魔鬼}的名称，因为他已从天上之地越过，在地上显现为天主诫命的诽谤者。然而，因为是天主首先颁布了诫命，主耶稣自己对魔鬼说：\BibleQuote{撒殚，退到我后面去！}无疑，走到天主后面乃是作祂仆人的标志。祂又说：\BibleQuote{你要朝拜上主你的天主，惟独侍奉祂。}因此，由于有些人倾向于顺从他的意愿，救主就用这样的话对他们说：\BibleQuote{你们是出于你们的父魔鬼，并愿意追随你们父亲的欲望。}最后，当他们被发现确实在行他的旨意时，就对他们说：\BibleQuote{毒蛇的种类！谁指教你们逃避那即将来临的忿怒？那么，就结出与悔改相称的果实吧！}从这一切，你就应当看出，为人类具有自由意志是何等重大的事。然而，请我的对手在这里说一说，对虔敬者和不虔敬者有没有审判？\Emph{玛内斯说}：有审判。\Emph{阿凯拉奥斯说}：我认为，我们关于魔鬼所说的话，既包含不小的理性，也包含不小的虔敬。再者，每个受造物都有其本身的秩序；人类有一个秩序，动物有另一个秩序，天使又有另一个秩序。而且，只有一位不可改变的实体，即天主的实体，永恒而不可见，这是众所周知的，也由以下圣经所证实：\BibleQuote{从来没有人见过天主，只有那在父怀里的独生子，祂给我们详述了。}因此，所有其他受造物，必然都是可见的——如天、地、海、人、天使、总领天使。但如果天主从未在任何时间被任何人看见过，那么祂与那些受造物之间能有什么同体性呢？因此，我们认为，一切事物按照它们各自的秩序，在各自的位置上，都有其本有的实体。另一方面，你声称每一个活动的生物都是由一（质）造成的，并且你说每一个对象都从天主那里领受了同样的实体，这实体能够犯罪并受审判；你又不愿接受那宣称魔鬼原是一位天使，他因过犯而堕落，并且他不是与天主体性相同的声明。从逻辑上讲，你应该摒弃任何关于审判教义的认可，而这样一来，我们当中谁有错误就清楚了。实际上，如果由天主所造的天使不可能因过犯堕落，那么灵魂作为天主的一部分，又怎能犯罪呢？然而，反过来，如果你说对犯罪的灵魂有审判，并且还主张这些灵魂与天主体性相同；而且，即使你坚持它们具有天主性体，你仍断言，尽管如此，它们却不遵守天主的诫命，那么，即便基于这些理由，我的论证也非常有力，即主张魔鬼首先堕落了，因为他未能遵守天主的诫命。他确实不是天主的性体。他堕落，与其说是为了伤害人类，不如说是为了被人类所轻蔑。因为祂\BibleQuote{赐给我们权柄，可以践踏蛇蝎，并制伏仇敌的一切势力。}\par

33. \Emph{判官们说}：关于魔鬼的起源，他已经提供了充分的说明。并且既然双方都承认将会有审判，那么这一承认必然包含着每个人都被显示拥有\OriginalTerm{自由意志}{free will}；这一点既已清楚显明，毋庸置疑，每个人在运用自身意志的正当能力时，都可以随心所欲地决定自己的道路。\Emph{梅内斯说}：如果正如你们声称的，善来自（你们的）天主，那么你们就使耶稣自己成了说谎者。\Emph{阿基劳斯说}：首先，要承认我们所引述的叙述是真实的，然后我将就\Quote{他的父亲}向你们提出证明。\Emph{梅内斯说}：如果你向我证明他的父亲是说谎者，却又表明你无论如何都不归咎于天主任何（这样的恶念），那么在各方面你都会取得信任。\Emph{阿基劳斯说}：诚然，一旦关于魔鬼的全面叙述被呈现出来，并且\OriginalTerm{分施}{dispensation}被阐明，任何人只要具有一般敏锐的理解力，只需在脑中仔细思考，便能领会这儿所说的\Quote{魔鬼之父}是谁。但尽管你自称是\OriginalTerm{护慰者}{Paraclete}，你却远不及常人的智慧。因此，既然你暴露了自己的无知，我就告诉你\Quote{魔鬼之父}这个表述的含义。\Emph{梅内斯说}：我这样说……；\Emph{他补充说}：任何事物的奠基者或制造者都可以被称为他所造之物的父亲，\Emph{父母}。\Emph{阿基劳斯说}：好吧，我真的惊讶于你对我所说的话做出了如此正确的承认，并且没有隐藏你对该断言的理解，也没有隐藏其实质。现在，由此可知这位魔鬼之父是谁。当他从\Place{天国}堕落时，他来到地上居住，并留在此处，时刻留意寻找某人，可以依附于他，通过与自己结盟，也使对方成为自己邪恶的同伙。那么，在人类尚未存在时，魔鬼从未连同他的父亲被称为杀人者或说谎者。但是后来，当人一旦被造，并且被魔鬼的谎言和诡诈所欺骗，而魔鬼也进入蛇的身体，那蛇是百兽中最狡猾的，从那以后，魔鬼就与他的父亲一同被称为说谎者，并且诅咒不仅停留在魔鬼自己身上，也停留在他的父亲身上。因此，当蛇接纳了他，并且确实完全将他容纳于自身之内，可以说，蛇就怀孕了，因为它担负了魔鬼极大的邪恶；它如同怀胎的妇人，经受着分娩的阵痛，力图排出他恶毒暗示的搅扰。因为蛇嫉妒第一个人的荣耀，便进入了\Place{乐园}；在自身中孕育着这分娩的痛苦，它开始发出虚伪的言辞，并为那些由天主所造、并领受了生命恩赐的人带来死亡。然而，魔鬼未能藉着蛇完全彰显自己；他将其完美保留了一段时间，好借着\Person{加音}——他藉由加音被完全生出——来展示。因此，一方面通过蛇，他向\Person{厄娃}展示了他的伪善和欺骗；另一方面通过\Person{加音}，他实施了杀人的开端，将自己引入那\Quote{果实}的初熟之物中，而此人管理得如此糟糕。由此，魔鬼从起初就被称为杀人者，也是说谎者，因为他欺骗了那些他对其说\BibleQuote{你们将如同天主一样}\BibleRef{创 3:5}的人；因为他假意宣称要成为神的那等人，后来却被逐出了\Place{乐园}。所以，那在腹中怀他、生他、并把他带到日光之下的蛇，被立为魔鬼的第一位父亲；而\Person{加音}成为他的第二位父亲，他借着孕育邪恶而产生痛苦和弑亲：因为夺去人的性命实乃施行不义、邪恶与不虔的总和。此外，所有接纳他并行他私欲的人，都成了他的兄弟。\Person{法郎}是他的完美父亲。每个不敬虔的人都成了他的父亲。\Person{犹达斯}成了他的父亲，因为他确实怀了他，却流产了：因为他未能带来完美的分娩，因为实际上通过\Person{犹达斯}受攻击的是一位更尊贵者；因此如我所说，那是一次流产。因为正如妇人接受男人的种子，从而感受到内里的日渐成长，\Person{犹达斯}同样每天都在邪恶中长进，那些机会是由那恶者像种子一样提供给他的。他里面邪恶的第一粒种子，就是贪财；其增长便是偷窃，因为他偷取了放在钱袋中的钱。其后代，更包括那些小小的纷争，与法利塞人的勾结，以及肮脏的讨价还价；然而在可怕的绞索中他结束生命所见证的，是流产，而非诞生。这情形正与你相同：如果你在自己的行为中将那恶者显露出来，并行他的私欲，你就怀了他，并要被称为他的父亲；但另一方面，如果你珍惜补赎，并卸下你的重担，你就如同一个生产的人。因为，如同在学校练习中，如果有人从老师那里得到主题材料，然后自己创作并产生整篇论述，他便被称为他所生出之作品的作者；同样，凡从原初之恶中领受了哪怕一点点邪恶的酵母的人，也必然被称为那从起初就抵抗真理的恶者的父亲和创造者。事实上，那些献身于美德的人也可能有同样的情形；因为我听到最勇敢的人对天主说：\BibleQuote{上主啊，因敬畏祢，我们怀孕，经历产痛，且生出了拯救之灵。}\BibleRef{依 26:18} 因此，那些相对恶者的畏惧而怀孕，并生出邪恶之灵的人，也必须被称为其父亲。这样，一方面，只要他们仍顺服于他的服事，他们就被称为那恶者的儿子；但另一方面，如果他们达到邪恶的极致，就被称为父亲。因为我们的主正是基于此意对法利塞人说：\BibleQuote{你们是出于你们的父亲魔鬼}\BibleRef{若 8:44}，这样当他们仍然被他搅扰，并在心中为善者图谋以恶报善时，就使他们成为他的儿子。因此，当他们心中如此盘算，并且他们邪恶的计谋归于他们自身责任时，\Person{犹达斯}，作为万恶之首，并将其不义的计谋付诸实施达于极致的人，被立为此罪行之父，他收取他们给予的三十块银钱作为其不敬虔之残忍的酬报。因为\BibleQuote{在吃那片饼之后，撒殚进入了他}\BibleRef{若 13:27}，完全地进入。但是，如我们所说，当他的子宫扩大，产期临到，他只是生出了不义孕育中的流产之担，因而他不能被完全地称为完美的父亲，除非仅在孕育尚在腹中之时；此后，当他自缢于绞刑架，显出他未能达成完全的出生，因为悔恨随之而来。\par

34. 我想你也不可能不明白这一点，即\Quote{父亲}这个词确实只是一个单一的词语，却可以按不同的方式理解。因为有人被称作父亲，是因为他是那些由他自然生育的子女的生身之父；另有人被称作父亲，是因为他充当了他所抚养的子女的监护人；还有人在时间或年岁上具有优越地位，因而也被称为父亲。因此，我们的主耶稣基督自己也被说成有多种父亲：因为\Person{达味}被称为祂的父亲，\Person{若瑟}也被算作祂的父亲，然而这两人中没有一个是在自然本质意义上作祂父亲的。\Person{达味}被称为祂的父亲，是在时间与年岁的尊荣方面；\Person{若瑟}被称为祂的父亲，是在抚养的法律方面；但天主自己才是祂本性上唯一的父亲，祂乐意藉着祂的圣言在短时间内向我们彰显一切。而我们的主耶稣基督，毫不迟延，在一年之内便使众多病人恢复健康，并使死者重获生命的光明；祂确实以自己圣言的能力拥抱了万有。那么，祂究竟在哪里迟延了，竟使我们不得不相信祂等待了如此之久，\Emph{甚至直到这些日子}，才切实地派遣了\OriginalTerm{护慰者}{Paraclete}呢？恰恰相反，正如上文已经说过的，祂立即证明了祂与我们同在，并且极其丰沛地将自己通传给了\Person{保禄}，对于保禄的见证我们同样相信，他说：\BibleQuote{这恩宠唯独赐给了我。}\BibleRef{弗 3:8} 因为这人从前是天主的教会的迫害者，但后来在众人面前公开显为护慰者的忠信仆人；藉着他，祂那无与伦比的仁慈得以彰显于众人，甚至我们这些曾一度毫无希望的人，也得享祂恩赐的厚惠。我们中有谁能指望迫害教会、与教会为敌的\Person{保禄}，竟会成为教会的捍卫者和守护者呢？不但如此，他还要成为教会的治理者，教会的建立者和建造师呢？因此，在他之后，以及在与祂同在的人——即门徒们——之后，我们就不应再期盼任何其他（此类）人物来临，这是依据圣经；因为我们的主耶稣基督论及这位护慰者时说：\BibleQuote{祂要由我领受。}\BibleRef{若 16:14} 因此，祂拣选了他作为悦纳的器皿；祂在圣神中派遣\Person{保禄}到我们这里来。圣神被倾注在他内；正如那圣神不能居住在所有人身上，而只能居住在那由天主之母\Person{玛利亚}所生的人身上一样，那护慰者的圣神也不能进入任何别的人，只能降临在宗徒们和\Person{圣保禄}身上。因为祂说：\BibleQuote{他是我拣选的器皿，为把我的名带到君王和外邦人面前。}\BibleRef{宗 9:15} \Person{保禄}宗徒自己也在他的第一封书信中表明了同样的事，他说：\BibleQuote{依天主赐给我的恩宠，使我为外邦人成了耶稣基督的使臣，天主福音的司祭。}\BibleRef{罗 15:15-16} 他又说：\BibleQuote{我在基督内说实话，并不说谎，有我的良心在圣神内与我一同作证。}\BibleRef{罗 9:1} 又说：\BibleQuote{因为我不敢提别的什么事，只说基督藉我，以言语和行动所成就的事。}\BibleRef{罗 15:18} \BibleQuote{我原是宗徒中最小的一个，不配称为宗徒。然而，因天主的恩宠，我成为今日的我。}\BibleRef{格前 15:9-10} 他愿意去对付那些寻求那位在他内说话的基督之证据的人，正是因为护慰者在他内；并且，他已蒙受了恩宠的赠礼，充满了极大的荣耀，他说：\BibleQuote{为这事，我曾三次求主，叫这刺离开我。但主对我说：有我的恩宠为你够了，因为我的德能在软弱中才全显出来。}\BibleRef{格后 12:8-9} 再者，我们的主耶稣基督在福音中表明，那在\Person{保禄}内的正是护慰者自己，祂说：\BibleQuote{如果你们爱我，就要遵守我的命令；我也要求父，祂必会赐给你们另一位护慰者。}\BibleRef{若 14:15-16} 在这话中，祂指的就是护慰者自己，因为祂说的是\Quote{另一位}护慰者。因此，我们信赖了\Person{保禄}，并听从了他所说的话：\BibleQuote{你们寻查那在我内说话的基督的证据吗？}\BibleRef{格后 13:3} 以及他以类似方式表达的话，这些我们在上面已经提到过了。就这样，他也为我们如同忠信的继承者那样，印上了他的遗嘱；他如同一位父亲，在致\Place{格林多}人的书信中用这些话对我们说：\BibleQuote{我当日把我所领受的，又传给你们了，其中首要的是：基督照经上记载的，为我们的罪死了，被埋葬了，且照经上记载的，第三天复活了。并且显现给\Person{刻法}，以后又显现给那十一位宗徒；此后，又一起显现给五百多弟兄，其中多半到现在还在世，有些已经安眠。然后显现给\Person{雅各伯}，以后又显现给众宗徒；最后也显现了给我这个像流产儿的人。我原是宗徒中最小的一个。}\BibleRef{格前 15:3-8} 因此，\BibleQuote{总之，不拘是我，或是他们，我们都这样传了，你们也这样信了。}\BibleRef{格前 15:11} 还有一次，在把他最先获得的产业交付给他的继承者时，他说：\BibleQuote{但我恐惧，唯恐你们的心意败坏，失去在基督内所当有的纯洁，就像那蛇以诡计诱惑了\Person{厄娃}一样。因为如果有人来，给你们宣讲另一位基督，不是我们所宣讲过的，或是你们领受另一神，不是你们所领受过的，或是另一福音，不是你们所接受过的，你们竟然容忍了！我原想我一点也不在那些超等的宗徒之下。}\BibleRef{格后 11:3-5}\par

35. 这些事，他之所以这样说，是为了指示我们，在他以后来的，都将是假宗徒、欺诈的工人，他们冒充基督的宗徒。这不足为奇，因为连撒殚自己也冒充光明的天使。所以，如果他的仆役也冒充正义的使者，这有什么大不了的呢？——他们的结局必按他们的行为而定。他进而指出这些人是怎样的人，并说明他们被谁操纵。当迦拉达人想要背离福音时，他对他们说：\BibleQuote{我真奇怪，你们竟这样迅速地脱离了那藉基督的恩宠召叫你们的天主，而归向了另一福音；其实，并没有另一福音，只是有些人扰乱你们，企图改变基督的福音而已。但是，无论是我们，或是从天上来的天使，若给你们宣讲的福音，与我们以前给你们宣讲的福音不同，当受诅咒。}\BibleRef{迦 1:6-8} 他又说：\BibleQuote{我原是宗徒中最小的，竟蒙受了这恩宠；}\BibleRef{弗 3:8} 以及：\BibleQuote{如今我在为你们受苦，反觉高兴，因为这样我可在我的肉身上，为基督的身体——教会，补充基督的苦难所欠缺的。}\BibleRef{哥 1:24} 在另一处，他又宣称自己比众人更劳苦作基督的仆役，好像在他之后，再无可期待的人了；因为他甚至命令，连从天上来的天使也不可这样接待。那么，我们怎能相信这个来自波斯、自称是\OriginalTerm{护慰者}{Paraclete}的摩尼的宣告呢？正因如此，我反而认出他是那些冒充者之一，正如宗徒保禄，那\OriginalTerm{特选的器皿}{elect vessel}，十分清楚地指示我们的，他说：\BibleQuote{圣神明明地说：在最后的时期，有些人要背弃信德，听信迷惑人的邪神和魔鬼的训言，这训言是出于那些伪善的说谎者，他们的良心已经麻木，犹如被热铁烙了一般。他们禁止嫁娶，且\Emph{命人戒避}天主所造、为给那些信奉并认识真理的人，以感恩的心所享用的食物。因为天主所造的样样都好，如以感恩的心领受，没有一样是可摒弃的；}\BibleRef{弟前 4:1-4} 圣神也在圣史\BibleBook{玛窦}{福音}中，谨慎地记下了我们的主耶稣基督的这些话：\BibleQuote{你们要谨慎，免得有人欺骗你们，因为将有许多人假冒我的名字来说：我就是基督；他们要欺骗许多人。那时，若有人对你们说：看，基督在这里，或说：在那里，你们不要相信。因为将有假基督和假先知出现，显大奇迹和大异迹，以致如果可能，连被选的人也要受欺骗。看，我预先告诉了你们。那么，若有人对你们说：看，他在旷野里，你们不要出来；或说：看，他在内室里，你们不要相信。}\BibleRef{玛 24:4-5,23-26} 然而，尽管有了这一切的指示，这个人，既没有显任何奇迹或异兆给人看，也没有显出任何关系，甚至从未在门徒之列，从未跟随过我们已故的主（我们因他的产业而喜乐）——我说，这个人，尽管从未在我们的主软弱时侍立在他身旁，从未作过他遗嘱的见证人，甚至从未接近过那些在他患病时服侍他的人，总之，他得不到任何人的见证，却要我们相信他自称为\OriginalTerm{护慰者}{Paraclete}的宣告；然而，即使你能行奇迹异事，按圣经所说，我们仍要视你为假基督、假先知。因此，我们最好更谨慎行事，遵照圣宗徒在写给哥罗森人的书信中所给的警告，他说：\BibleQuote{只要你们在信德上站稳，根基稳固，不偏离福音的望德。这福音就是你们听过的，而且已传给普天之下一切受造之物；}\BibleRef{哥 1:23} 又说：\BibleQuote{你们既然接受了基督耶稣为主，就该在他内行事为人，在他内扎根建造，在所教给你们的信德上站稳，满怀感恩之情。你们要小心，免得有人以哲学和虚伪的妄言，按照世俗的传授，而不按基督，把你们夺去。因为在天主性的一切圆满，都具体地居住在他内；}\BibleRef{哥 2:6-9} 在这一切都仔细说明之后，这位真福宗徒如同父亲对子女说话一般，又加上了以下的话，作为他遗嘱的印记：\BibleQuote{这场好仗，我已打完；这场赛跑，我已跑到终点；这信德，我已保持了。从今以后，正义的冠冕已为我预备下了，就是主，公正的审判者，到那一日必要赏给我的；不但赏给我，也赏给一切爱慕他显现的人。}\BibleRef{弟后 4:7-8}\par

36. \Person{摩尼}啊，你无法使任何人成为迦拉达人，你也无法用这种方法使我们偏离对基督的信德。是的，即使你行神迹奇事，即使你复活死人，即使你将\Person{保禄}本人的肖像展示给我们，你仍然是被诅咒的。因为我们预先已受教导关于你的事：圣经已经警告我们并武装我们来对付你。你是假基督的器皿；实际上，不是尊贵的器皿，而是卑贱的器皿，被他用来正如任何野蛮人或暴君可能做的那样：当他试图入侵一个生活在法律正义之下的民族时，会预先派一个特选的器皿，仿佛是注定要死的，以便探知那位合法君王及其国民所拥有的力量的确切规模和性质。因为那人太害怕自己完全不加防备地入侵，也缺乏胆量派遣任何自己身边亲近的人去执行这样的任务，唯恐他可能遭受某种伤害。因此，你的君王假基督，以类似的角色派遣你，仿佛是注定要死的，来到我们这些处于善而圣的君王治理之下的人民这里。我这样说并不是不加考虑或未经适当调查；而是因为我看到你并未行任何奇迹，我认为自己有资格对你怀有此类看法。因为我们预先得知，魔鬼自己会变成光明的天使，他的仆役也会以类似的假象出现，他们会行神迹奇事，以致于，如果可能的话，连选民也会受欺骗。可是，请问，你是谁呢？你的父\Person{撒殚}并没有分配给你这样的亲属地位。你曾使谁从死里复活？你曾止住了谁的流血？你曾用泥抹瞎子的眼睛，从而使他们看见吗？你曾何时用几个饼使饥饿的群众得以饱足？你曾在何处行走在水面上？或者，居住在\Place{耶路撒冷}的人中有谁曾见过你？波斯野蛮人啊，你从来不能通晓希腊人、埃及人、罗马人或其他任何民族的语言；你只通晓加色丁语，而这种语言确实不是在广大民众中普遍使用的语言，你也无法听懂任何其他民族的人说话。圣神却不是这样：绝非如此；祂分施给众人，通晓各种语言，明了万事，向什么样的人，就作什么样的人，以致人心的心思意念都无法逃脱祂的知晓。因为圣经说什么呢？\BibleQuote{众人都听见宗徒们藉着圣神，也就是\OriginalTerm{护慰者}{Paraclete}，用自己的语言说话。} 但是，我何必在这个问题上多说呢？\Person{密特拉}的蛮族司祭和狡猾的助手啊，你只会是太阳神\Person{密特拉}的崇拜者，按你的看法，他是玄秘之处的光照者，也是自觉的神明；也就是说，你会像他的崇拜者那样嬉戏，并且，尽管可以说没那么雅致，你也会庆祝他的神秘仪式。但是，我何必如此愤慨呢？难道这一切不都是合适的吗？你应当像莠子一样繁殖，直到你那位大能的父来临，他复活死人——\Emph{正如他将声称要做的那样}——并迫害几乎所有拒绝顺从他命令的人，几乎要到地狱里去，他用他那种自大的恐怖来辖制群众，并用威胁对付其他人，通过改变容貌和欺诈行为来戏弄他们。然而，他不能超越此限；因为他的愚昧必向众人显露，正如\Person{雅纳尼}和\Person{杨布勒}的情形一样。\Emph{法官们说}：正如我们从你这里听到的，以及\Person{保禄}本人似乎也告诉我们，并且我们也同样从福音较早的叙述中得知，宣讲、教导、传福音或说预言的入门，至少在此生中，并不以同样的条件给予宗徒时代之后的任何人：如果相反的情形似乎曾经出现，那人只能被视为假先知或假基督。如今，既然你声称\OriginalTerm{护慰者}{Paraclete}在\Person{保禄}内，并且祂在\Person{保禄}内证实一切，\Person{保禄}自己怎么说呢？\BibleQuote{我们现在所知道的，只是一部分；我们作先知所讲的，也是一部分；及至那完全的来到，这部分的就必被废去。} 他说这些话时，等候的是谁呢？因为如果他声称自己正在等候某个完全的，如果必然有人要来，请告诉我们他指的是谁；免得他的话似乎把我们引向这个人，就是\Emph{\Person{摩尼}}，或者派他来的那位，也就是按你的断言，\Person{撒殚}。但如果你承认那完全的尚未来到，这就排除了\Person{撒殚}；如果你等候\Person{撒殚}的来临，那又排除了完全的。\par

37. \Emph{\Person{阿尔该劳斯}说}：圣保禄所提出的那些言论，并非没有天主的指引而说出的，因此可以确定，他向我们宣告的是：我们要仰望我们的主耶稣基督为\OriginalTerm{那完全者}{the perfect one}，祂是唯一认识圣父的那一位，除了祂所选择要将圣父启示给他的那人外，没有别人知道圣父，如同我能从祂自己的话所证实的。但请注意，经上说，当那完全的来到时，那局部的就要被废弃了。现在这个人（\Person{摩尼}）自称他就是那完全者。那么，让他向我们指明，他废除了什么；因为要被废除的，乃是存在于我们内的无知。因此，让他告诉我们，他废除了什么，又给\Emph{我们知识的领域}带来了什么。如果他能做任何这类事情，就让他现在去做，好让他可以被人相信。保禄的这些话，只要有人能充分理解其含义的力量，就只会使我的陈述完全可信。因为在\BibleBook{格林多}{前书}中，保禄论到那将要来的完全时，这样说：\BibleQuote{如有先知之恩，终必消失；如有语言之恩，终必停止；如有知识之恩，终必消灭。因为我们现在所知道的，只是局部的；我们作先知所讲的，也只是局部的；及至那完全的来到时，这局部的就要消失了。}\BibleRef{格前 13:8-10}现在，请注意那完全者本身拥有什么德能，又有何种等级的成全。那么，让这个人告诉我们，他废除了犹太人或希伯来人的什么预言；或者他使何种语言停止了，无论是希腊人的，还是其他崇拜偶像之人的语言；或者他摧毁了哪种异端教义，无论是\Person{瓦伦提努斯}派的、\Person{马西盎}派的、\Person{塔齐安}派的、\Person{撒贝利乌斯}派的，还是其他任何为自己编造了一套独特知识体系的人。让他告诉我们，他已经废除了这一切中的哪一样，或者作为这完全者，他何时还要废除它们中的任何一样。或许他正在寻求某种休战——是吗？但那位\OriginalTerm{真正完全者}{the truly perfect one}，即我们的主耶稣基督，祂来临的方式绝不会这样微不足道，绝不会这样晦暗卑下。不，就像一位君王临近自己的京城时，首先派遣他的侍卫、旌旗、军旗与旗帜，以及他的将帅、首领与总督走在前面，随后一切事物都被激发并震动，以不同的方式，有的因预见君王降临而恐惧，有的则欢欣鼓舞；同样，我的主耶稣基督，那位真正完全者，在祂降临时，也将先派遣祂的荣耀，\Emph{以及}一个无玷无瑕国度的圣洁先锋走在祂前面：随后，整个受造界将被撼动和扰乱，发出祈祷与恳求，直到祂将其从奴役中解救出来。必然的，人类种族那时将因自己所犯的众多过犯而处于恐惧与剧烈的激动中。那时，唯有义人将欢欣，因为他们期待着所应许给他们的事物；而此世的事务将不再维持，万物都将被毁灭：无论是各种预言，或是先知的著作，都要消失；无论是全人类的语言，都要停止；因为人将不再需要为生活必需品而忧虑或操心；无论是知识，无论是由何种教师所掌握，也都要被毁灭：因为这一切事物中，没有一样能承受那位大能君王的来临。因为正如一点火星，若被拿来与太阳的光辉相比，就立即从眼前消失，同样，整个受造界，一切预言，一切知识，一切语言，正如我们上面所说的，都将被毁灭。然而，普通人性之能力全然不足以用这样寥寥数语，且如此软弱又如此极度贫乏的言辞，来描述这位天上君王之来临——的确，或许只有那些圣洁且极高贵的人，才有特权尝试就此主题发表任何论述——但就我而言，\Emph{能够说出}我现在就此主题所阐述的这些，乃是出于单纯的必然——我这样做，实在是被此人的固执所逼，且仅仅是为了向你们表明，他是何等样的人。\par

38. 说实话，我认为\Person{马吉安}、\Person{瓦伦提努}、\Person{巴西里德}和其他异端者，与这人相比，都算得上是圣人了。因为他们确实显示出某种智慧，并且他们确实自以为能够理解全部圣经，并因此自封为那些愿意听从他们之人的领袖。尽管如此，他们中没有一个敢宣称自己是天主，或是基督，或是\OriginalTerm{护慰者}{Paraclete}，而这个人却这样做了。他总是在争论，有时论及永世，有时论及太阳，以及这些事物是如何被造的，好像他自己超越于它们之上；因为任何阐述某个事物如何被造的人，就是在把自己置于所讨论的对象之上，并且比它更古老。然而，谁敢谈论天主的本体呢？或许唯有我们的主耶稣基督才能这样做。而且，我并不是仅凭自己言语的权威作出这项陈述，而是用那曾教导我们的圣经的权威来证实它。因为宗徒对我们说了以下的话：\BibleQuote{为使你们在世界上成为光，持守生命的话，好让我在基督的日子得荣耀，因为我并非徒然奔跑，也非徒然劳苦。}\BibleRef{斐 2:15-16}我们应当理解这句话的力量和意义；因为这话语或许适用于领袖，但有效的工程则属于君王。正如一个期待其君王来临的人，努力使所有在他管理下的人成为顺服、预备好、可称赞、可爱、忠心，并且同样无可指摘、充满一切美善，好让自己能从君王那里获得嘉许，并被认为配得更大的尊荣，因他妥善治理了所托付给他的省份；同样，真福\Person{保禄}也用这些话让我们明白我们的处境：\BibleQuote{为使你们在世界上成为光，持守生命的话，好让我在基督的日子得荣耀。}因为这句话的意思是，我们的主耶稣基督来临的时候，会看到祂的道理在我们身上产生了效益，并且发现他，这位\Emph{宗徒}，并没有徒然奔跑，也没有徒然劳苦，就将酬报的冠冕赐给他。在同一封书信中，他也警告我们不要思念地上的事，并告诉我们应当以天堂为家乡；我们也正是从那里期待救主，我们的主耶稣基督。因为知道末日的日期对我们来说并不稳妥，他就在写给\Place{得撒洛尼}人的书信中，就这一点为我们作出了声明，如此说：\BibleQuote{弟兄们，论到时期和时候，不需要给你们写什么；因为你们原确实知道，主的日子要像夜间的盗贼一样来到。}\BibleRef{得前 5:1-2}那么，这个人怎么还敢挺身而出，试图说服我们接受他的观点，纠缠他所遇到的每一个人成为摩尼教徒，四处奔走，潜入人家，尽力欺骗那些背负罪孽的人呢？但我们不持这种看法。相反，我们更愿将事实本身摆在你们所有人面前，如果你们愿意，可以将它们与\Emph{我们所知的}完美护慰者进行比较。因为你们注意到，他有时用询问的语气，有时用哀求的语气。但在我们救主的福音中记载，那些站在君王左边的人会说：\BibleQuote{主啊，我们什么时候见祢饿了，或渴了，或赤身露体，或作客，或在监里，而没有服事祢呢？}\BibleRef{玛 25:44}他们就这样恳求祂宽待他们。但那位公义的审判者和君王将对他们作出怎样的答复呢？\BibleQuote{你们这些作恶的人，离开我，到永火里去吧！}\BibleRef{玛 25:41}祂把他们投入永火，尽管他们不停地向祂哀求。那么，\Emph{\Person{摩尼}}啊，你看见那位完美君王的来临注定是怎样的情形了吗？你难道没有发觉，它并不像你所宣称的那样是一种完美，\Emph{或圆满}吗？但如果那伟大的审判之日要在那位君王之后才被期待，那么这人肯定远远低于祂。但如果他是低下的，他就不能是完美的。如果他不是完美的，那么宗徒所说的就不是他。但如果宗徒所说的不是他，而他却仍然谎称\Emph{宗徒的话}是指着他自己说的，那么他就必定被判定为假先知。还有许多话可以说到同样效果。但如果我们想详细陈述所有类似能引用的论点，时间将不足以完成如此庞大的任务。因此，我认为从许多事中只提请你们注意到这一些，已是绰绰有余了，而把这类讨论的其余部分留给那些有意深入探究的人。\par

39. 在场的人听了这些事，便大大光荣天主，将相称于祂的赞颂归于祂。他们也向\Person{阿尔赫劳斯}本人致以许多敬意。随后\Person{玛尔凯路斯}站起身，脱下外氅，抱住\Person{阿尔赫劳斯}，亲吻他，拥抱他，紧紧依偎着他。接着，那些恰好聚集在周围的孩子们也开始向\Person{玛内斯}投掷东西，要把他赶走；其余的众人也跟随他们，激动地四处走动，旨在迫使\Person{玛内斯}逃跑。\Person{阿尔赫劳斯}见了，便如号角般在一片喧嚣中高声呐喊，急于制止群众，对他们说：\Quote{住手，我亲爱的弟兄们，免得在审判的日子我们身上背负流血的罪责；因为论到这样的人，经上写着：\BibleQuote{在你们中间原该有异端，好叫那些经得起考验的人显明出来。}}他说了这话，群众便又安静下来。如今，因为\Person{玛尔凯路斯}乐意让这场辩论得以记载并描述下来，我无法违逆他的愿望，便相信读者们会宽仁体谅，相信他们若嫌我的叙述略显拙朴或粗鄙，也会宽宥；因为我们始终以之为重的，乃是务必让所有愿意了解此事的人都能获知个中情由。此外，必须补充，\Person{玛内斯}一经逃离，就再也没有\Emph{在那里}出现过。不过他的随从\Person{图尔博}却被\Person{玛尔凯路斯}交给了\Person{阿尔赫劳斯}；\Person{阿尔赫劳斯}祝圣他为执事\OriginalTerm{执事}{deacon}后，他便留在了\Person{玛尔凯路斯}的侍从中。但\Person{玛内斯}在逃亡中来到一个离城相当远的村庄，名叫\Place{迪奥多鲁斯}。那地方也有一位长老\OriginalTerm{长老}{presbyter}，名字同样叫\Person{迪奥多鲁斯}，性情安静温和，无论因他的信德还是因他卓越的人品，都享有美誉。某一天，\Person{玛内斯}聚拢了一群听者，向他们高谈阔论，把一些全然与教父传统相悖的外道言论摆在在场的人们面前，全未察觉他们之中可能会有人反对他，\Person{迪奥多鲁斯}见他正以邪说产生某种影响，便决意给\Person{阿尔赫劳斯}送去一封信，措辞如下：——\par

\Person{迪奥多鲁斯}向\Person{阿尔赫劳斯}主教\OriginalTerm{主教}{bishop}问安。\par

40. 最虔诚的神父，我希望您知道，在这些日子里，我们这里来了一个名叫 \Person{摩尼} 的人，他宣称自己是要来完成 \OriginalTerm{新约}{New Testament} 的教义的。在他所作的言论中，的确有些内容可能与我们信仰相符；但他也有某些断言似乎与我们由教父的传统所传承下来的相去甚远。因为他以怪异的方式解释某些教义，将他自己的一些见解强加其上，在我看来，这些见解与信仰完全格格不入并且对立。基于这些事实，我现在被促使给您写这封信，因为我知道您在教义上的智慧既完全又丰盛，并且确信这一切事都无法逃过您的认知。因此，我怀有坚定的希望，相信您不会因任何怨恨而不向我们解释这些事。至于我自己，确实不可能被引入任何新奇教义；然而，为了所有蒙昧的人，我被引导以您的权威请教一句话。因为事实上，这个人表现出自己具有非凡的人格力量，无论在言语上还是行为上；确实，甚至他的外貌和服饰也证实了这一点。但我要在此写给您一些要点，是我从他阐述的所有题目中所能记忆下来的：因为我知道，即使通过这些少数几点，您也会了解其余的。您无疑清楚地知道，那些试图建立任何新信条的人，习惯于非常轻易地把他们想从圣经中取用的任何证据，曲解为符合他们自己的观念。然而，对此早有先见，宗徒的话语标记出这样的情况：\BibleQuote{若有人给你们宣讲的福音，与你们所领受的不同，他就该受诅咒。}因此，除了宗徒们一次而永远托付给我们的之外，基督的门徒不应接受任何新的教义。但为了不把我必须说的事拖得太长，我直接回到眼下的主题。那么，这个人主张，简而言之，\OriginalTerm{梅瑟法律}{the law of Moses}不是来自美善的天主，而是来自\OriginalTerm{邪恶之王}{the prince of evil}；并且它与\OriginalTerm{基督的新法律}{the new law of Christ}没有任何亲属关系，反而与之对立和敌视，二者互为直接的仇敌。当我听到有人提出这种观点时，我向民众重复了福音中的那句话，我们的主耶稣基督论到他自己说：\BibleQuote{我来不是为废除法律，而是为成全。}然而，那人断言他根本没有说过这句话；因为他认为，当我们发现他确实废除了那同样的法律时，我们首要关注的，应当是他实际所行的事，胜过其他所有考量。接着，他开始引述法律中的许多不同章节，也引述许多来自福音和宗徒保禄的章节，这些章节似乎彼此矛盾。所有这一切，他同时以十足的自信宣讲出来，毫无犹豫或畏惧；以至于我确实相信，他有那条永远与我们为敌的蛇作他的助手。于是，他宣称，\Emph{在法律中}天主说：\BibleQuote{我使人富贵，也使人贫穷；}然而\Emph{在福音中}耶稣却称贫穷的人是有福的，又加上，除非人舍弃他所有的一切，就不能作他的门徒。他又主张，梅瑟在人民逃出埃及时，从埃及人那里拿走了金银；而耶稣却颁布了诫命，叫我们不可贪图属于我们邻人的任何事物。然后他断言，梅瑟在法律中规定了，应以眼还眼、以牙还牙作为惩罚；而我们的主却吩咐我们，连另一边面颊也转向那打我们一边面颊的人。他也告诉我们，梅瑟命令，凡在安息日作任何工，以及不持守法律所写的一切事的人，都要受惩罚并被石头砸死，正如，由于无知而在安息日捡了一捆柴的那个人所遭受的；而耶稣却在安息日治好了一个瘫子，并命令他那时也拿起他的床。更进一步，他没有阻止他的门徒在安息日掐麦穗，用手搓着吃，而这在安息日原是不许可做的事。我何必再提其他事例呢？因为借着许多类似性质的不同断言，他这些教义被以最大的精力和最炽烈的热忱提了出来。同样，借着一位宗徒的权威，他试图确立这样的立场：梅瑟法律是死亡的法律，而耶稣的法律，相反地，是生命的法律。因为他把那个断言建立在这段经文上：\BibleQuote{祂也使我们能够做新约的仆役，不是凭着字句，而是凭着圣神，因为字句叫人死，圣神却叫人活。那以字句刻在石头上的死亡的职务，尚且有光荣，甚至以色列子民为了梅瑟面貌上易于消逝的光荣，不能注视他的面容；如果那易于消逝的尚且有光荣，那属神的职务岂不更该有光荣吗？如果定罪的职务有过光荣，那么成义的职务更该多么充满光荣！那先前有过光荣的，因了这更超越的光荣，已算不得光荣了；因为如果那易于消逝的曾一度有过光荣，那常存的光荣更该多么大呢！}正如您也十分清楚的，这段经文出现在\BibleBook{致格林多}{后书}中。除此以外，他又加上取自\BibleBook{致格林多}{前书}的另一段经文，根据这段经文，他断言\OriginalTerm{旧约}{Old Testament}的门徒属于地和血肉；并据此认为，血肉不能承受天主的国。他还主张，保禄亲口说了这话：\BibleQuote{如果我重新建筑我所拆毁的，我就自认是罪犯。}而且，他断言，同一位宗徒最显然地就\OriginalTerm{肉身的复活}{resurrection of the flesh}这一主题作出了这一声明，他也说：\BibleQuote{唯在外表上做犹太人的，不是真犹太人；在外表上肉身上的割损，也不是真割损；}并且根据字句，法律毫无益处。他又引用了这话：\BibleQuote{亚巴郎有了光荣，但不是在天主面前；}以及\BibleQuote{借着法律而来的只是对罪的认识。}他又提出许多其他事情，目的是贬损法律的尊荣，理由是法律本身就是罪；这些言论持续提出，使得较为单纯的人多少受了影响；根据这一切，他又利用了这个断言：\BibleQuote{法律和先知直到若翰为止。}然而，他宣称，若翰宣讲的是\Emph{真正的}天国；因为他确实认为，藉着他的头被斩下，就象征所有在他之前、比他优先的人都要被斩除，而唯独他以后所要来的应当存立。因此，就这一切事，哦，最虔诚的\Person{阿基劳斯}，请写一封简短的回信给我们：因为我听说您钻研这类事务非同寻常；而\Emph{您所拥有的能力}是天主的恩赐，因为天主要将这些恩赐赐予那些相配的人、祂的朋友，以及那些在志向与生活上表明与祂联合的人。因为我们的本分是预备自己，亲近那恩慈而慷慨的心，这样我们就会立刻从它那里领受最丰厚的恩赐。因此，既然我所有、用来讨论此类论题的学识，不能满足我的渴望和目的的要求——因为我承认自己是个没有学问的人——我已向您写信，正如我不止一次说过的，希望从您手中获得对这个问题最充分的解答。愿您安好，无与伦比且可敬的神父！\par

41. 接到这封信后，\Person{阿尔凯劳斯}对这人的大胆感到惊讶。但同时，由于情况需要迅速回复，他立即发出一封信，回应\Person{狄奥多鲁斯}的陈述。那封信内容如下：——\Person{阿尔凯劳斯}向可敬的儿子、\Person{狄奥多鲁斯}长老致以问候。\par

收到你的信使我极其欢喜，我亲爱的朋友。此外，我得知，这个人前些日子来到我这里，试图在此引入一种与宗徒的与教会的知识不同的新知识，他也到了你那里。我没有给这人容身之地：因为我们曾在一起辩论，他当场就被驳倒了。我现在希望能为你把我在那次辩论中所使用的全部论据抄录下来，好让你由此了解那人信仰的实质。但由于那需要我有闲暇才能办到，我认为鉴于当前的急需，有必要就你写给我的有关他所提出的主张一事，给你写一封简短的回信。我了解，他的主要努力旨在证明梅瑟的法律与基督的法律不相一致；他试图以我们的\WorkTitle{圣经}的权威作为这一论点的依据。而我们这一方面，不仅用同一部\WorkTitle{圣经}确立了梅瑟的法律以及其中所写的一切，而且还证明了全部\WorkTitle{旧约}与\WorkTitle{新约}一致，完全和谐，它们实在构成一体，就像一个人可以看到同一件衣服由纬线和经线织成一样。因为事实就是，正如我们在衣服上辨认紫线，我们若可如此表达的话，就能在\WorkTitle{旧约}的纹理中分辨出\WorkTitle{新约}；因为我们看见主的荣耀映照在其中。所以，我们不应当抛弃这面镜子，既然它忠实地、真实地显示了事物本身的真像；相反，我们应当愈加珍视它。你真的以为，那个在幼时由他的\OriginalTerm{训蒙师傅}{paedagogue}带到教师面前的小男孩，长大成人后，就因为不再需要他的服务，无需那位陪伴者协助就能自己去学校、迅速找到讲堂，便应当不敬重那位训蒙师傅吗？或者，再举一个例子，一个起初吃奶的孩子，长大能接受较硬的食物后，若以怨恨之心弃绝乳母的胸怀，并对之感到厌恶，这是对的吗？不，他应当敬重珍爱它们，并承认自己受惠于它们的服务。如果你愿意，我们还可以再举一个例子。有一个人，一次看见一个婴孩被丢弃在地上，已受尽苦楚，便把他抱起来，带回家中抚养，直到他长成年少，并承受了所有抚养孩子的人通常会有的辛劳和焦虑。然而，过了一段时间，孩子的生父前来寻找这个男孩，并在抚养他的人那里找到了他。这孩子得知这是他的生父后，应该怎么做？我所说的，当然是一个品行端正的孩子。他岂不会设法让抚养他的人得到丰厚的酬报，然后，考虑到自己应有的产业，再去跟从自己的生父吗？同样，我认为我们必须设想，那位天主杰出的仆人梅瑟，以某种类似的方式，发现了一个受埃及人压迫的民族；他把这个民族带到自己身边，在沙漠中像父亲一样养育他们，像教师一样教导他们，像官长一样管理他们。他也保护这民族，直到那本属于这民族的主来到。过了相当长的时间后，那位父亲来了，领回了他的羊。那么，那位监护人难道不会因交付羊群而在一切事上受尊敬吗？那些因他而得保全的人，岂不会荣耀他吗？那么，我亲爱的\Person{Diodorus}啊，谁会愚蠢到说那些彼此结盟的人、彼此替代对方说预言的人、以及显出了同等同质、彼此类似的奇迹异事的人，是彼此互为异己的呢？或者，按实话说，这些彻头彻尾属于同一种类的？的确，梅瑟首先对百姓说：\BibleQuote{上主我们的天主，要从你们弟兄中间，给你们兴起一位像我一样的先知；} 而耶稣后来说：\BibleQuote{因为梅瑟论及了我。} 你看，这两者如何相互握手，虽然一位是先知，另一位是爱子，虽然我们在一位身上认出忠信的仆人，而在另一位身上则是主自己。另一方面，我可以指出一个事实，即从前有人想要不经训蒙师傅而去学校，老师不收他。因为老师说：\Quote{除非他接受训蒙师傅，我绝不收他。} 这个比喻中所说的人物是谁，我将简要解释。从前有一个富人，按外邦人的方式生活，天天奢华宴乐；又有一个穷人，是他的邻居，连每日的食粮都得不到。这两个人都离了世，都进入了坟墓，那穷人被带到安息之所，等等，正如你所知道的。而且，那富人还有五个兄弟，他们像他一样生活，对于他们从这样一位老师在家中所学得的教训毫不怀疑。那富人于是恳求，让这些人一起立刻接受高超的教义。但亚巴郎知道他们仍需要训蒙师傅，就对他说：\BibleQuote{他们有梅瑟和先知。} 因为，如果他们不接受这些，从而使自己的道路受他——即\Emph{梅瑟}——如同训蒙师傅的指引，他们就不能接受那更高超师傅的教义。\par

42. 但我也要尽我所能，对上述其他一些言论进行阐释；也就是说，我要指出\Person{耶稣}的所言所行，并没有任何与\Person{梅瑟}相悖之处。首先，关于\BibleQuote{以眼还眼，以牙还牙}这句话——这是\Emph{正义的体现}。至于祂的吩咐，就是人若被打了一边脸颊，便应将另一边也转给他，那是\Emph{良善的体现}。那么，正义与良善互相反对吗？绝非如此！这不过是从单纯的正义，进到了积极的良善。再者，我们也有这样的说法：\BibleQuote{工人自当有他的工价。}但若有人意图从中施以欺诈，那么他被要求吐出其以欺诈所得之物，无疑是最为正义的，尤其是当工价可观之时。我这样说，是指当埃及人借着被派来监督他们造砖的督工苦待以色列子民时，\Person{梅瑟}立时之间，一并将惩罚加诸其身，全数追讨。但这是否便可被称为不义呢？绝非如此！实际上，当一个人只是适度地使用真正必需之物，而放弃一切超出其上的，这确然是良善的自持。再者，让我们来看看，在旧约中我们发现这样一句话：\Quote{我造了富人和穷人，}而\Person{耶稣}却称贫穷的人是有福的。其实，\Person{耶稣}在那句话里，所指的并非仅仅是世间财物匮乏的穷人，而是指那些神贫的人，意即那些不因骄傲而膨胀，而是怀有温良谦卑的性情，不高看自己的人。然而，这个问题我们的对手并未正确地提出。因为我在此看出，\Person{耶稣}也欣然地注视着富人投入银库的献仪。同时，若唯有富人将献仪投入银库，那也实在是太过贫乏；因此，那穷寡妇的两个小钱，同样也被欣然收纳了；在那样的献仪中，确然展现出了某些超越\Person{梅瑟}关于收纳银钱所规定之事。因为\Person{梅瑟}从那些拥有者收取献仪；而\Person{耶稣}却甚至从那些没有者收纳。可是，这个人又进一步说，经上写着：\BibleQuote{人若不撇下一切所有的，就不能作我的门徒。}然而，我再次注意到，那百夫长，这位拥有极大财富且世俗势力雄厚的显赫之人，却怀着超越全以色列的信心；以至于，即使真有人撇下了一切，那人的信心也被这百夫长所超越。但如今也许有人会这样与我们辩论：因此，放弃财富并非好事。我回答：对于有能力这样做的人，那是好的；但同时，将财富用于行正义与怜悯的事工，其受悦纳的程度，并不亚于立刻全然舍弃。再者，关于安息日已被废除的断言，我们否认祂已明令废除了它；因为祂自己也是安息日的主。而且，这\Emph{法律与安息日的关系}，就好比一位负责照管新郎新房的仆人，他悉心打点一切，不容许任何外人扰乱或触碰，将其完好保存以待新郎来临；这样，当新郎到来时，便可按己意使用，或者准许那些他吩咐一同进去的人使用。而主\Person{耶稣基督}本人，也为我们所肯定的事作了见证，当时祂以属天的声音说：\BibleQuote{新郎还伴同在的时候，伴郎岂能禁食呢？}再者，祂也并没有实际上废弃割损；我们更应该说，祂在自己身上，并代替我们，承受了割损的缘由，借着祂自己受的苦解救了我们，不容我们白白地遭受任何痛苦。因为，若一个人虽然行了割损，内心却对近人怀着最恶毒的意念，这于他又有何益呢？故此，祂更愿意借着一条简捷的途径，向我们开启最圆满生命的道路，以免我们在行过了自己漫长的路程之后，却发现白昼过早地被黑夜吞没，也以免我们虽然外在地在世人的眼中显得光鲜，内里却只堪比择肥而噬的豺狼，或如同粉饰的坟墓。因为，即便一个人身穿褴褛破旧的衣衫，只要他心中没有隐藏任何对近人的恶念，就远比那种品格的人高超得多。因为惟有内心的割损能带来救恩；而那仅仅肉身的割损，若没有灵性的割损与之同在，于人并无益处。请听圣经对此的论述：\BibleQuote{心里洁净的人是有福的，因为他们要看见天主。}那么，既然我已得知这简捷的生命之道，并且知道只要我内心洁净，这便为我所有，我何须还要辛苦\Emph{并受苦}呢？而这完全符合我们如今所领受的真理，也就是说，若有人谨守这两条诫命，他便成全了全部的法律和先知。此外，宗徒之长\Person{保禄}，在阐述了这一切之后，更给了我们对此事越发清晰的指教，他说：\BibleQuote{或者你们愿意寻求那在我内说话的基督的证据吗？}那么，既然我在未受割损时即可成义，割损于我有何干呢？因为经上记载：\BibleQuote{有人受过割损吗？就不要恢复未受割损的样子。有人未受割损吗？就不要受割损。因为这些都算不得什么，惟有遵守天主的诫命才是重要的。}因此，割损既不能拯救任何人，便无甚可要求的，尤其是当我们看到，若有人蒙召时未受割损，而后来却愿意受割损，他即刻便成了法律的违犯者。因为，我若受割损，也是为着能得救而遵行法律的诫命；但我若未受割损，且保持未受割损，便更因遵守诫命而获得生命。因为我在圣神内领受了内心的割损，而非那仅凭墨迹的字句的割损，前者所获的称赞，并非来自人，而是来自天主。因此，谁也不要将这类指责加诸于我。这正如一个拥有众多金银财宝的富人，他家里一切需用之物皆由此等贵重金属制成，故在其家业中无需展示任何陶器；然而，这并不得出结论说，陶匠的产品或制陶技艺应被他憎恶。同样，我因天主的恩宠而成为富足的，并获得了内心的割损，便无论如何也不需要那全然无益的\Emph{属肉身的}割损；然而，尽管如此，我也并不会因此就说割损是恶的。绝对不可！不过，若有人愿意对这些事接受更精确的教导，他在宗徒的第一封书信中会发现对此极尽详尽的讨论。\par

43. 我现在要极其简略地谈谈\Person{梅瑟}的帕子和属死的职务。因为我不认为这些事至少能给法律带来多大的贬损。相关的经文如此说：\BibleQuote{但若那以文字刻在石头上属死的职务，尚且是在光荣中进行的，以致以色列子民因\Person{梅瑟}面容的光荣，不能注目凝视他的脸；而这光荣原是容易消逝的；}等等。所以，这段经文至少承认\Person{梅瑟}面容上有光荣存在，而这无疑是对我们立场有利的事实。即使这光荣要消逝，并且在诵读时还有一个帕子，只要其中仍有光荣，这一点并不使我烦恼或不安。也不是说，凡要消逝的东西，就在任何情形下都沦为耻辱的境地。因为当圣经论及光荣时，也向我们表明，它知道光荣是有等次的。于是它说：\BibleQuote{太阳的光荣是一样，月亮的光荣是另一样，星辰的光荣又是另一样：因为星辰与星辰在光荣上也有差异。}因此，虽然太阳的光荣比月亮大，但并不能据此说月亮就沦为耻辱的境地。同样，虽然我的主\Person{耶稣基督}在光荣上超越\Person{梅瑟}，如同主人超越仆人，但也不能据此说\Person{梅瑟}的光荣就应受轻视。因为我们也能这样满足我们的听众，因为这话语本身的性质带着说服力，我们既然依据圣经本身的权威来肯定我们所主张的，又借着取自圣经的例证使我们陈述的证明更为清晰。这样，虽然一个人在夜间点着一盏灯，但太阳一旦升起，他便不再需要他那灯盏微弱的亮光，因为太阳的光辉普照全世界；然而，尽管如此，这人并不会轻蔑地将他的灯扔掉，好像它是与太阳完全对立的东西似的；相反，当他一旦发现灯的用途后，他反而会更加小心地保存它。正是如此，\Person{梅瑟}的法律便如同灯盏，作百姓的某种监护者，直到真正的太阳——即我们的救主——升起，正如宗徒也对我们说：\BibleQuote{基督必要光照你。}但我们还必须看下文所说的：\BibleQuote{他们的心智被蒙蔽了：因为直到今天，诵读旧约时，这同样的帕子仍然存在；它未被揭去，因为是在基督内才被消去的。因为直到今天，诵读\Person{梅瑟}时，帕子还蒙在他们的心上。然而，当他们对主转过头来，帕子就要被揭去。主就是那神。}那么，这是什么意思呢？难道\Person{梅瑟}到今天仍与我们同在吗？难道他从未安眠，从未安息，从未离开人世吗？这里所说的\BibleQuote{直到今天}这个说法究竟如何解释？那么，只需注意他所说帕子的所在，就是他们诵读时蒙在他们心上的帕子。因此，这就是针对以色列子民的责备的话，因为他们诵读\Person{梅瑟}却不理解他，并且拒绝转向主；因为主正是\Person{梅瑟}预言将要来的那一位。所以，这就是蒙在\Person{梅瑟}脸上的帕子，这也是他的\OriginalTerm{遗嘱}{testament}；因为他在法律中如此说：\BibleQuote{权杖不离犹大，柄杖不离他腿间，直到那应得令牌者来到；他是万民的期望：他要把小驴拴在葡萄树上，把驴驹拴在精选的葡萄树上；他要在酒中洗他的外衣，在葡萄汁中洗他的长袍；他的眼睛因酒润泽，他的牙齿因乳洁白；}等等。此外，他还指示了他是谁，以及他要从哪里来。因为他说：\BibleQuote{上主你的天主，要从你中间，从你兄弟中，为你兴起一位像我一样的先知，你们应听从他。}然而这显然不能理解为指农的儿子\Person{若苏厄}说的。因为在他身上并没有见到这事的应验。而在他之后，犹大仍有君王出现；因此，这预言绝不适用在他身上。这就是蒙在\Person{梅瑟}身上的帕子；因为这帕子并非如某些无学问的人所想的，是什么亚麻布帕子，或遮盖他面容的什么皮子。但宗徒也留意向我们说明这一点，他告诉我们说，帕子仍蒙在诵读旧约上，因为那些自古以来就被称为以色列的人，仍然等待基督的来临，却没有察觉权杖已离开犹大，柄杖已离开他腿间；正如我们目前所见，他们臣服于君王和元首，并向他们纳贡，自身再无审判或惩罚的权力，而这样的权力犹大确曾有过，因为他在定\Person{塔玛尔}的罪之后，仍能宣判她为义。\BibleQuote{但你将亲见你的性命悬于你眼前。}\par

44. 这话也有\OriginalTerm{帕子}{veil}。因为直到\Person{黑落德}的时代，他们确实似乎还拥有某种王国；是\Person{奥古斯都}在他们中间进行了第一次人口登记，他们开始纳税、被征税。也正是从我们的\Person{主耶稣基督}开始被预言和期待的时候起，\Place{犹大}出了君王和民众的首领；而这些，又恰好在祂来临时消失了。那么，若那在他们诵读时所蒙的帕子被揭去，他们就会明白\OriginalTerm{割损}{circumcision}的真实意义；他们也会发现，我们所宣讲的那一位的世代、祂的十字架，以及在我们主的历史中所发生的一切事，正是那位先知所预言的那些事。我确实希望逐一查考每段这样的圣经经文，指出其应被理解的意义。但既然如今另有急务，这些经文留待闲暇时再讨论。因为目前，我所说过的可能已足以表明，在诵读旧约时，某些人心上蒙着帕子并非没有理由。但那些归向主的人，帕子就被揭去了。然而这一切事究竟有多大的力量，我留给那些有健全理智的人去领悟。现在让我们回到\Person{梅瑟}的那句话，他说：\BibleQuote{上主你的天主，要从你们弟兄中间，给你兴起一位像我一样的先知。}我在这句话中看到，仆人\Person{梅瑟}传达了一个伟大的预言，他知道那将要来的的确拥有比他更大的权威，却要遭受与他相似的事，并显示相似的征兆和奇事。因为在那里，梅瑟出生后，被他的母亲放在一个蒲草箱里，放在河岸边；在这里，我们的主耶稣基督，由祂的母亲\Person{玛利亚}所生后，借着天使的帮助逃往\Place{埃及}。在那里，梅瑟带领他的百姓从\Person{埃及人}中间出来，拯救了他们；在这里，耶稣带领祂的百姓从\Person{法利塞人}中间出来，把他们转移到永远的救恩中。在那里，梅瑟藉着祈祷寻求食物，并从天上领受，为在旷野中喂养百姓；在这里，我的主耶稣，以自己的能力在旷野中用五个饼饱饫了五千人。在那里，梅瑟在山上受试探，禁食四十天；在这里，我的主耶稣被圣神引到旷野，受魔鬼试探，也同样禁食四十天。在那里，由于\Person{法郎}的背叛，所有\Person{埃及人}的长子在梅瑟眼前都死了；在这里，在耶稣降生时，由于\Person{黑落德}的背叛，所有犹太人男婴突然都死了。在那里，梅瑟祈祷，希望\Person{法郎}和他的百姓免受灾祸；在这里，我们的主耶稣祈祷，希望\Person{法利塞人}得赦免，当祂说：\BibleQuote{父啊，宽恕他们吧！因为他们不知道他们做的是什么。}在那里，梅瑟的面容因主的荣光而发光，致使以色列子民因他面容的荣光，不能注视他的脸；在这里，主耶稣基督发光如太阳，门徒因祂面容的荣光和极度耀目的光辉，不能注视祂的脸。在那里，梅瑟用刀剑击杀了那建造金牛的人；在这里，主耶稣说：\BibleQuote{我来是为把刀剑投在地上，并使人与近人不和，}等等。在那里，梅瑟无畏地进入载水的乌云黑暗中；在这里，主耶稣以大能履海。在那里，梅瑟向海发命；在这里，主耶稣在海上起来，向风和海发命。在那里，梅瑟受攻击时，伸出手与\Person{阿玛肋克}交战；在这里，主耶稣，当我们受攻击并因那现今在义人中活动的错谬之神的力量而丧亡时，在十字架上伸出了祂的双手，赐予我们救恩。然而还有许多其他的这类事，我必须略过，我极为亲爱的\Person{狄奥多若}，因为我急于尽快将这小册子寄出；这些省略之处，你自己可以通过你的悟性极容易地加以补充。请给我写信，告知这位作仇敌事业之仆人将来可能做的一切。愿全能的天主保全你灵魂与精神完全无恙！\par

45. 收到这封信后，\Person{狄奥多鲁斯}掌握了信的内容，然后便开始与\Person{摩尼}辩论。他进行得如此充满热忱，以致于所有人都大大称赞他，因为他谨慎而令人满意地论证了两个约之间以及两部法律之间存在着相互的关系。他为自己发现了更多的论据，因而能够提出许多切中要害且有力的论点来反对那人，并为真理辩护。他也基于词语的理由，以令人信服的方式反驳了他的对手。例如，他以下列方式与他辩论：——你说过有两个约吗？那么，你要么说有两个旧约，要么说有两个新约。因为你断言有两个属于同一时间，或更确切地说属于同一永恒的\OriginalTerm{非受生者}{unbegotten}：而如果这样就有两个，那么或者应当有两个旧约，或者应当有两个新约。然而，如果你不承认这一点，却反过来肯定只有一个旧约，并且还有另一个新约的重申，这便只能再次证明两者只有一位作者；而且这种次序本身也表明，旧约属于那一位新约也归属的主。我们可以用一个例子来说明：一个人对另一个人说，\Quote{把你的旧房子租给我。}因为这样的说法，岂不表明那人也是新房子的主人吗？或者，反过来说，如果他对他说，\Quote{给我看看你的新房子！}用那个词，岂不也指定了他也是旧屋的拥有者吗？此外，还需要考虑这一点：既然有两个具有非受生性体的存在者，由此就必须假定每一个都有（必须称为）旧约之物，这样就会出现两个旧约；如果你确实肯定这两个存在者都是古老的，而且两者都无起始的话。但我并没有学到过这样的教义；\WorkTitle{圣经}也不包含此说。然而，你声称\Person{梅瑟}法律来自恶者之王，而非出自善的天主，那么请你告诉我，那些当面反对\Person{梅瑟}的人是谁——我是指\Person{雅乃斯}和\Person{杨布勒}吗？因为，凡反对之物，所反对的不是自己，而是另外的某一位，或更善或更恶；正如\Person{保禄}在\BibleBook{弟茂德}{后书}中所写的，也使我们明白：\BibleQuote{就如雅乃斯和杨布勒反抗梅瑟，照样这等人也反抗了真理。这些人，他们的心志已败坏，在信德上是不可靠的。但他们再不能往前了，因为他们的愚昧将在众人前暴露，如那两个人一样。}\BibleRef{弟后 3:8-9} 你注意到了吗，他将雅乃斯和杨布勒比作心地败坏、在信德上背弃的人，而另一方面，他却将梅瑟比作真理？但是，圣\Person{若望}，这位最伟大的圣史，也告诉我们恩宠上加恩宠的赐予和散布；因为他确实指出，我们从\Person{基督}的满盈中领受了梅瑟法律，他的意思是，藉着\Person{耶稣基督}，那一个恩宠，这另一个恩宠在我们内得以圆满。我们的主\Person{耶稣基督}亲自用下面的话表明了这一点：\BibleQuote{不要想我要在父面前控告你们；有一位控告你们的，就是你们所寄望的梅瑟。若是你们信了梅瑟，必会信我，因为他是指着我写的。如果你们不信他的著作，怎么会信我的话呢？}\BibleRef{若 5:45-47} 除了所有这些话之外，还可以从宗徒\Person{保禄}和\WorkTitle{福音}中引出许多其他的章节，藉此我们能够证明，旧法律不属于别人，只属于那位新约也归属的主，而且我们最好将这些章节引证出来，并以令人满意的方式加以运用。然而现在，傍晚使我们不能这样做了；白天即将结束，我们理当结束这场辩论。但明天会给你们机会，就你们喜欢的任何问题向我们提问。说完这些话，他们就走了。\par

然而，次日清晨，在任何人尚未出门活动之前，\Person{亚尔基老斯}突然出现在\Place{狄奥多鲁斯}所住的寓所。因此，\Person{摩尼}全然不知\Person{亚尔基老斯}再次来到现场，便公开挑战\Person{狄奥多鲁斯}与他进行辩论；他意在用言辞炫耀来压垮对方，因为他看出\Person{狄奥多鲁斯}秉性单纯，对圣经问题并无深究。因为他已领略过\Person{亚尔基老斯}的教训。因此，当群众再次聚集在通常用于辩论的场所，当\Person{摩尼}刚开始推理时，\Person{亚尔基老斯}突然出现在他们中间，拥抱了\Person{狄奥多鲁斯}，并以圣洁的吻问候他。这时，\Person{狄奥多鲁斯}和所有在场的人，确实对天主眷顾的安排充满了惊奇，因为\Person{亚尔基老斯}恰在问题\Emph{恰好}被提出之时来到他们中间；因为实际上，必须承认，尽管\Person{狄奥多鲁斯}十分虔诚，他对这场冲突还是有些惧怕。但当\Person{摩尼}见到\Person{亚尔基老斯}时，他立刻收起了侮辱的态度；他的骄傲被大大打击，他清楚地表明自己乐意逃离这场较量。然而，听众们将\Person{亚尔基老斯}的来临视如宗徒的降临，因为他显得装备齐全，敏捷而准备好用言语为\Emph{真理}辩护。于是，\Person{亚尔基老斯}举起右手示意民众安静——\Emph{因为不小的骚动已经发生}——之后，开始作如下致辞：——尽管我们中间有些人已获得了智慧的尊荣和荣耀的奖赏，但我仍请求你们，将我来之前所说过的话的见证存记于心。弟兄们，我知道并确信，我现在接替\Person{狄奥多鲁斯}的位置，并非由于他有什么不足，而是因为我先前已认识眼下这个人，那时他通过那位声名显赫的\Person{马塞勒斯}的帮助，带着他邪恶的计划闯入我居住的地区，他试图使\Person{马塞勒斯}背离我们的教义和信德，目的就是让他成为这不虔敬教训的有效支持者。尽管如此，尽管他言辞动听，却未能动摇他，也未能使他在任何一点上偏离信德。因为这位最虔敬的\Person{马塞勒斯}，被证实犹如那建造在磐石上的房屋，根基极其稳固；雨水降下，洪水涌来，风吹袭，撞击那房屋，它却屹立不倒：因为它建基于最坚固、不可动摇的根基上。而这位如今站在你们面前的人所作的尝试，给他自己带来的不是荣耀，而是羞辱。此外，在我看来，如果他证明对未来之事一无所知，那是不可原谅的；因为他当然应该预先知道谁是支持他的人：如果那\OriginalTerm{护慰者}{Paraclete}的圣神确实住在他内，他必定拥有这等程度的知识。但既然他实际上是一个被无知的黑暗所蒙蔽的人，他跑到\Person{马塞勒斯}那里去是徒然的，他不过是表明自己如同观星者，忙于描述天上的事物，却对自己家中发生的事情一无所知。但为了避免让人觉得我通过这些说词而搁置了当前的问题，我现在就停止这类谈论。并且我也将给这人特权，让他选择他认为最合适的任何要点，作为对这个主题和问题的讨论的开端。至于你们，如我已说过的，我只请求你们作公正的审判者，以便将适当的荣誉和胜利的棕榈枝归给讲说真理的人。\par

47. 于是，当所有人安静下来后，\Person{摩尼}这样开始发言：\Person{阿基劳斯}，你像其他人一样，用最伤人的话语打击我，尽管我对天主的看法是正确的，我对基督也持有正确的观念；然而，宗徒的行列更是以承受一切、忍耐一切为特征，即使有人用辱骂和诅咒攻击他们。如果你们想迫害我，我已做好准备了；如果你想要我受惩罚，我决不退缩；是的，即使你想要杀死我，我也不怕：\BibleQuote{因为我们应当害怕的，只有那能使灵魂和肉身都陷于地狱中的那位。}\Emph{\Person{阿基劳斯}说}：绝对不要这样想！这不是我的意图。因为你曾在我，或与我们同心的人手中受过什么苦呢？即使在你贬低我们、伤害我们，在你诋毁我们先祖的传统，在你存心使那些在真理中得坚固、并以极谨慎明辨之心被保守之人的灵魂死亡的时候，——而事实上，全世界的财富都不足以作为报偿，你又曾受过什么苦呢？然而，你凭什么采取这个立场呢？你有什么可展示的呢？告诉我们这个——你有什么救恩的记号能摆在我们面前呢？因为仅靠夸口的话，既不足以使在座的众人满意，也不足以使他们有资格辨识我们哪个更正确地掌握了真理的知识。因此，既然你得到了先发言的机会，请先告诉我们，你希望我们将辩论引向哪个主题的哪一点。\Emph{\Person{摩尼}说}：如果你不再一次无理地反对我将以完全恰当的方式陈述的论点，我就会与你交谈；但如果你仍要表现出我上次所看到的那种态度，我就转向\Person{狄奥多罗斯}发言，避免你的搅扰。\Emph{\Person{阿基劳斯}说}：我已经表明了我的想法，即我们仅仅用空话来浪费时间简直就是在滥用这个机会。如果任何一方有意提出不公正的反对，就留给审判者决定吧。但现在，请告诉我们你必须提出的论点。\Emph{\Person{摩尼}说}：如果你不打算再一次仅仅为了反对我以完全正确的方式陈述的立场而反对，我就开始。\Emph{\Person{阿基劳斯}说}：\Quote{如果不这样}和\Quote{如果不那样}，这样的说法正是无知之人的标志。因此，你对未来的事是无知的。但至于你所宣称的仍是未来的这件事，反对或不反对完全在我自己。那么，你寄予信赖如同一面最坚固盾牌的那两个树的论点，又将如何成立呢？因为如果我属于对立面，你如何要求我顺服呢？反过来，如果我有顺服的意向，你又为何如此怕我反对你呢？因为你主张恶永远为恶，善永远为善，全然不知自己话语的力量。\Emph{\Person{摩尼}说}：难道我雇了你做我言论的辩护人，以至于你也能决定与我知识相称的智慧吗？当你连自己的事都说不清楚的时候，又如何能解释别人的事呢？但如果\Person{狄奥多罗斯}现在承认自己被驳倒了，那么我的论证就会针对你。然而，如果他仍坚持并且准备发言，我恳请你让开，不要妨碍真理的证实。因为你是一只陌生的羊；不过，日后你将被领入同一羊群，正如\Person{耶稣}的声音也表明的——那位\Person{耶稣}，即确实以人形显现，却并非人的那一位。\Emph{\Person{阿基劳斯}说}：那么，你不认为祂是由童贞\Person{玛利亚}所生吗？\Emph{\Person{摩尼}说}：天主不允许我承认我们的主耶稣基督是通过女人自然的母胎降到我们这里的！因为祂亲自给我们作证说，祂是从父的怀抱中来的；再者祂说：\BibleQuote{接受我的，就是接受派遣我来的；}又说：\BibleQuote{我来不是为执行我的旨意，而是为执行派遣我来者的旨意；}又说：\BibleQuote{我被派遣，只是为了以色列家失落的羊。}还有无数其他类似含义的章节，都指出祂是\Emph{来}的，而不是\Emph{生}的。但如果你比祂更大，如果你比祂更知道何为真理，我们又为何还信祂呢？\Emph{\Person{阿基劳斯}说}：我既不比祂大，因为我是祂的仆人；我也不可能与我主同等，因为我是祂无用的仆人；我是祂话语的门徒，我相信祂所说的一切，并且肯定这些话语是不可改变的。\Emph{\Person{摩尼}说}：有一个类似你的人曾对祂说：\BibleQuote{你的母亲和你的兄弟站在外面；}祂却没有善意回应，而是责斥了说这话的人，说：\BibleQuote{谁是我的母亲？谁是我的兄弟？}并且指出，那些遵行祂旨意的人，就是祂的母亲和兄弟。然而，如果你硬要说\Person{玛利亚}就是祂的生母，你就使自己陷入极大的危险。因为毫无疑问，按同样原则，必定会证明祂也有由她所生的兄弟。现在告诉我，这些兄弟是\Person{若瑟}所生，还是同一圣神所生？如果你说他们是同一圣神所生，那就会有许多的基督了。如果你说他们不是同一圣神所生，却又断言祂有兄弟，那么毫无疑问，我们将不得不认为，继圣神之后、在\Person{加俾额尔}之后，至洁无玷的童贞女实际与\Person{若瑟}结了婚。但如果这也完全荒谬——我是指她与\Person{若瑟}有任何交合的假设——告诉我，那时祂有兄弟吗？你这样是要把奸淫的罪也加在她身上吗，最聪明的\Person{马塞洛}？然而若这些假设无一符合这位无玷童贞女的地位，你又如何能断言祂有兄弟呢？如果你无法清楚地向我们证明祂有兄弟，那么按照那个胆敢写道\BibleQuote{看，你的母亲和你的兄弟站在外面}的人所说的话，你若要证明\Person{玛利亚}是祂的母亲，岂不就更难了么？然而，尽管那人胆敢如此对祂说话，没有谁能比显示母亲或兄弟的这位本人更有力或更伟大。不，祂甚至不配听到有人说祂是\Person{达味}之子。可是，宗徒\Person{伯多禄}，所有门徒中最卓越的一位，却能在那个场合承认祂，当时众人都提出各人对祂的看法：因为他说：\BibleQuote{你是基督，永生天主子；}祂立即称他为有福的，对他说：\BibleQuote{因为我的天父启示了你。}注意耶稣所说的这两句话之间的区别。对于说\BibleQuote{看，你的母亲站在外面}的人，祂回答说：\BibleQuote{谁是我的母亲？谁是我的兄弟？}但对于说\BibleQuote{你是基督，永生天主子}的人，祂则以福分和祝福回报。因此，如果你坚持认为祂是由\Person{玛利亚}所生，那么这就意味着，祂自己所说的话，不亚于\Person{伯多禄}的话，现在被证明是虚假的了。然而，如果\Person{伯多禄}说的是真的，那么毫无疑问，之前的那人是错误的。如果之前那人是错误的，那么问题就要归到那作者头上了。所以，我们知道只有一位基督，正如宗徒\Person{保禄}所说，他的话语至少与祂的来临相符，我们相信这些话语。\par

48. 听了这些话，聚集的群众大受震动，仿佛觉得这些论证给出了关于真理的正确说明，而阿基劳斯无力反驳他们；这从他们中间升起的骚动就可以看出来。但当听众静下来后，阿基劳斯便这样回答：的确，没有人能证明自己比我们的主耶稣基督的声音更有力，也没有任何名字能与祂的名相比，如经上所载：\BibleQuote{为此，天主举扬了祂，赐给祂一个名号，超越其他一切名号。}再者，在做见证方面，也没有人能比得上祂；因此我只需引证祂自己声音的见证来答复你们——首先，当然是为了解决你们所提出的那些难题，免得你们像惯常那样说，这些事与祂本人不一致。你们主张，那位把祂母亲和兄弟的消息带给耶稣的人，被祂斥责了，仿佛他犯了错误，如同那位执笔者犯了错误一样。而我断言，无论是那位给祂带去关于祂母亲和兄弟消息的人并未受到斥责，伯多禄也没有唯独被宣告为有福，高于那人；而是这两个人各自从祂那里得到了恰如其分、由他们各自的言词所激发的回答，正如随后的论述将要证明的。当人还是孩子时，他的思想像孩子，说话像孩子；但当他长大成人，就该把属于孩子的事除去：换句话说，当人奋力向前进取时，就会忘记背后的。因此，当我们的主耶稣基督致力于教导和医治人类，为使他们不致全然丧亡，并且所有聆听祂的人的心神都专注于这些事时，那位传讯者进来提醒祂有关祂母亲和兄弟的事，实在是极不合时宜的打扰。那么，你们自己来评判，祂该不该丢下那些祂正在医治和教导的人，去跟祂的母亲和兄弟说话呢？做这样的假设，岂不是立即贬低了祂本人的品格吗？再者，当祂拣选那些被罪恶重压所累的人，赐予门徒的荣耀，数目共十二位，并称他们为宗徒时，祂给了他们这项命令：要离开父母，为能配得上我；意思是，从此以后，对父母亲的记挂不该再削弱他们内心的坚贞。另一次，当有一个人对祂说：\BibleQuote{主，请许我回去安葬我的父亲。}祂回答：\BibleQuote{跟随我，任凭死人去埋葬他们的死人吧！}看，我的主耶稣基督就这样在一切必要的事上教导祂的门徒，并按各人相宜的，向每个人分赐祂的神圣训言。同样，在这另一场合，当有人冒失地来报告有关祂母亲的消息时，祂并没有乘机丢下祂圣父的使命，即便为了与母亲在一起也不去顾。但为了更清楚地向你们证明这才是事情的真相，让我提醒你们，伯多禄在得到祂那祝福的宣告之后，有一次对耶稣说：\BibleQuote{主，千万不要这样！这事绝不会临到你身上！}他这样说，是在耶稣向他宣告人子必须上耶路撒冷去，受死，并在第三天复活之后。那时祂回答伯多禄说：\BibleQuote{撒殚，退到我后面去！因为你不体贴天主的事，只体贴人的事。}既然你们认为，那位报告祂母亲和兄弟消息的人被耶稣斥责了，而稍早前说 \BibleQuote{你是默西亚，永生天主之子。} 的人得到了祝福的话语，那么请留意，耶稣（可以说）反而更偏爱那个人，就是祂屈尊给予更仁慈、更宽容答复的那位；至于伯多禄，虽然在得到祝福之后，如今却没有得到任何表示宽容的称呼，因为他未能慎思明辨向他所做的宣告的性质。因为那传讯者所犯的错误立即被答复的要旨纠正了；而这位宗徒领受的迟钝却被更严厉地斥责了。由此你们可以明白，主耶稣观看针对这些询问所宜和适时者，便给予每个人相称而恰当的答复。可是，假设如你们所说，伯多禄被宣告为有福，是因为他说了真实的话，而那位传讯者被责备是因为他所犯的错误，那么请告诉我，为什么当魔鬼承认祂，并说：\BibleQuote{我们认识你是谁：你是天主的圣者。} 祂却斥责它们，命它们不要作声？如果祂确实喜悦那些承认祂的人为祂所作的见证，那么为什么祂不用祝福来回报它们，如同祂对吐露真言的伯多禄所做的那样呢？但如果那是个荒谬的假设，那么我们只能这样理解：祂所说的话总是按照地点、时间、人物、主题，并对环境加以适当的考虑。因为只有当这种方法能救我们免于犯下鲁莽判断祂话语的错误，从而使我们免受应得的惩罚：并且这也将帮助我更加清楚地向你们说明，那位报告祂母亲消息的人反倒是更受尊荣的人。然而，你们忘记了我们讨论所提案的主题，转到别的话题上去了。但还是请你们听我片刻。因为如果你们确实愿意更仔细地思考那些话，我们就会发现主耶稣对前述两人中的前一位表现了极大的宽仁；我将用适合你们理解力的比喻来向你们证明这一点。有一位国王，拿起兵器出征迎敌，正在努力思考筹谋如何克敌制胜那些敌对的外邦军队。正当他满心怀着种种忧虑和焦思，在他冲入敌军之后，进而开始俘获他们的时候，焦虑的念头又压迫着他，如何保障那些与他一同辛劳作战、身负战争重担之人安全和利益，这时有个传讯者不合时宜地闯进来，开始向他提醒家务事。他对那人的大胆和不合时宜的提示感到惊讶，就想着要将这样的人处死。若非那传讯者能够藉着报告家中安好、一切顺利成功，而触动他内心最深的柔情，那么那应得的惩罚便会立即临到他了。因为，在战争持续的时日里，国王所关心的，除了设法保障他省份人民的安全，并料理军事之外，还能有什么呢？同样，那位传讯者也不合时宜地来到我的主耶稣基督前，不适时地带来了关于祂母亲和兄弟的消息，那时祂正在与攻击心灵堡垒的疾病争战，正在医治那些长久以来被各种软弱所困的人，并正竭尽全力确保众人的救恩。真的，那人本来可能受到像对伯多禄那样，甚至更严厉的判决。但一听祂母亲和兄弟的名字，就唤起了祂的宽仁。\par

49. 然而，除了以上所论述的一切之外，我还愿提出进一步的证据，好使所有人都能明白你这断言中包含著何等的不虔敬。因为如果你的主张是真的，即祂未曾诞生，那么无可怀疑地，祂也就未曾受难；因为未曾诞生的人是不可能受难的。但如果祂未曾受难，那么十字架的名义就被废除了。而如果十字架未曾被承受，那么\Person{耶稣}也就没有从死者中复活。如果\Person{耶稣}没有从死者中复活，那么其他的人也不会复活。如果没有人会复活，那么就不会有审判。因为确实地，如果我不会复活，我就不能受审判。如果没有审判，那么遵守天主的诫命就毫无意义，也没有节制的理由了；我们不如就说：\BibleQuote{我们吃喝吧，因为明天就要死了。} 因为当你否认祂由\Person{玛利亚}所生时，这一切后果便随之而来。但如果你承认祂由\Person{玛利亚}所生，那么祂的受难就必然随之而来，祂的复活也继受难而来，审判则随著复活而来：这样，圣经的训谕对我们便有了其应有的价值。因此这并不是一个无聊的问题，而是有最重大的问题系于此言。因为正如全部法律和先知都总括于两句话中，同样，我们所有的望德也都取决于真福\Person{玛利亚}所生者。因此，请回答我将向你提出的这几个问题。我们怎能摆脱宗徒这些如此重要而又如此明确的话呢？这些话表达如下：\BibleQuote{但当时期一满，天主就派遣了自己的儿子来，生于女人；} 又说：\BibleQuote{基督，我们的逾越节羔羊，已为我们作了牺牲；} 又说：\BibleQuote{天主既使主复活了，也要以自己的能力使我们与祂一同复活？} 此外还有许多意义相似的章节；例如，下面这一段：\BibleQuote{你们中怎么有人说，死人复活是没有的事呢？如果没有死人复活的事，基督也就没有复活；如果基督没有复活，我们的宣讲便是空的，你们的信德也是空的。并且，我们还被视为天主的假见证人，因为我们相反天主作证，说天主使基督复活了，其实并没有使祂复活；因为如果死人不复活，基督也就没有复活；如果基督没有复活，你们的信德便是假的，你们还在你们的罪中。那么，那些在基督内长眠的人也就丧亡了。如果我们只在今生寄望于基督，我们就是众人中最可怜的了。但是，基督从死者中实在复活了，做了长眠者的初果；} 等等。那么，我问，谁能如此鲁莽大胆，竟不让自己的信德与这些神圣的话相吻合呢？这些话中毫无保留，也毫无可疑之处。我请问你，愚昧的\Place{加拉太}人，谁迷惑了你呢？正如那些被迷惑的人，\BibleQuote{耶稣基督被钉在十字架上，已活现地摆在你们眼前。} 从这一切看来，我认为这些证据足以证明审判、复活和受难；而由\Person{玛利亚}所生，自然也立刻包含在这些事实之中。虽然你拒绝承认这一点，又有什么关系呢？因为圣经已最无误地宣告了这事实。不过，我要再向你提出一个问题，请你给我一个答复。当\Person{耶稣}为\Person{若翰}作证，说：\BibleQuote{在妇女所生的人中，没有兴起一位比洗者若翰更大的；但在天国里最小的，也比他大，} 请告诉我，在天国里有一个比他更大的，这是什么意思。难道在天国里\Person{耶稣}比\Person{若翰}小吗？我说，\Emph{天主不许！}那么，请告诉我这该怎样解释，你就必定超越你自己了。毫无疑问，\Emph{意思是说，} \Person{耶稣}在由妇女所生的人中比\Person{若翰}小；但在天国里，祂却比他大。所以，\Person{摩尼}啊，也请告诉我这一点：如果你说\Person{基督}不是由\Person{玛利亚}所生，而只是显为像人，却并非真的人，这显现是由在祂内的能力所造成和产生的，那么我再问你，\OriginalTerm{圣神}{Holy Spirit}像鸽子一样降在谁身上呢？那位受\Person{若翰}洗礼的是谁？如果祂是完美的，如果祂是\OriginalTerm{圣子}{Son}，如果祂是\OriginalTerm{能力}{Power}，圣神就不能进入祂内；正如一个国不能进入另一个国一样。并且，那从天上发出的声音，为祂作的见证：\BibleQuote{这是我的爱子，我所喜悦的，} 又是谁的声音呢？来吧，告诉我，不要拖延；是谁领受了这一切，行了这一切？回答我：你竟这样大胆地用亵渎的话当作理由，并且想为它找到立足之地吗？\par

50. \Emph{\Person{摩尼}说}：没有人会因为能够回答你刚才的指控而担心犯下亵渎之罪，反而应当被视为完全配得一切称赞。因为一位真正的技艺大师，当有任何事情引起他注意时，理当谨慎准备答复，并让所有人清楚了解所讨论或有所质疑的问题；尤其对尚未受过教导的人，更应如此。既然你对我们教义的说明并不满意，那么，请像一位彻头彻尾的技艺大师那样，也以合理的方式为我解答这个问题。因为在我看来，说天主子在实行祂降临世上之事时，完全不需要任何事物，这似乎是虔诚的；祂根本不需要鸽子，也不需要圣洗，不需要母亲，不需要弟兄，甚至也许不需要父亲——然而按你们的看法，这位父亲就是\Person{若瑟}——而是祂完全独自降下，随自己的美意，将自己变作\Emph{人的样子}，正如\Person{保禄}所说的那句话：\BibleQuote{形状也一见如人}\BibleRef{斐 2:7}。因此，请告诉我，那位能使自己变换成各种样貌的，还可能需要什么呢？因为当祂选择如此行时，祂又将这人的形貌和样子变成了太阳的模样。但是，如果你再一次反驳我，拒绝承认我正确表达了信仰，那么请听我对你所持立场的界定。因为如果你说祂只是\Emph{生于}\Person{玛利亚}的人，并且在受洗时领受了圣神，那么其结果就会是，祂被视作子，是由于渐进，而不是由于本性。然而，如果我允许你说祂是按渐进而成为子，并且被造成一个人，那么你的看法就是祂真是一个有血有肉的人。但这样一来，就必然推导出，那如同鸽子显现的圣神，无非就是一只自然的鸽子。因为这两种表述是相同的——即\Quote{如人}和\Quote{如同鸽子}；因此，无论你对使用\Quote{如人}这一表述的经文持何种观点，你也理应对使用\Quote{如同鸽子}这一表述的另一处经文持同样的观点。显然，必须把这两处归为一类，因为惟有如此，我们才能发现圣经中关于祂所写之事的真正含义。\Emph{\Person{阿基劳斯}说}：正如你仍不能为自己做到多少，不像一位彻底的技艺大师，那么我也无需费心去纠正这个问题，并百般忍耐地阐明它，给出明显的解答，除非是为了那些在场聆听我们的人。因此，为了这个缘由，我也要尽可能恰当地解释应当给予这个问题的答复。那么，你们似乎并不认为说耶稣有玛利亚作为母亲是虔诚的；对于你们现在以我完全不愿重复的措辞讨论的某些其他观点，你们也持类似的看法。然而，有时任何技艺的大师都会被对手的无知所迫，不得不说和做那些时间会\Emph{使他}拒绝的事情；因此，出于对在场众人的考虑，我不得不简要回答你那些错误百出的陈述。那么，现在让我们假设你的指控是：如果我们按照自然的过程，将耶稣理解为由玛利亚成胎的人，并因此认为祂有血有肉，那么也就必然要认定圣神是一只真实的鸽子，而不是神。好吧，可是，一只真实的鸽子怎能进入一个真实的人，并住在他内呢？因为肉体不能进入肉体。恰恰相反，只有当我们承认耶稣是真人，并且同时将那在那里说成是如同鸽子的那位认定为圣神，我们才能基于理性对双方都作出正确说明。因为，按照正确的理性，\Emph{可以这样说}圣神居住在人内，降临到他身上，并住在他内；而这些事的确已在完备而适宜的方式下发生过，并且其发生始终是可能的，正如你自己\Emph{所承认的，因为你}先前曾声称是天主的\OriginalTerm{护慰者}{Paraclete}——你这燧石，我这样称呼你，而不是人——常常连自己所声称的事都忘记。因为你宣称耶稣所应许要派遣的圣神降临到了你身上；而祂若非从天降下，又能从何而来呢？如果圣神如此降在那堪当领受的人身上，那么我们岂非要想象真有鸽子降在你身上？那么，我们倒不如认定你是一个偷卖鸽子的商人，为捕鸟而布置圈套和网罗。因为你实在配得上被讥讽的言语戏弄。然而，我放过你，唯恐我或许因这种表达而冒犯了听众，更是因为我的目的并非要把你当受的那些话全数抛出针对你。不过，让我回到正题。因为我记得你所谓的那个转变——据此你说天主已将祂自己转变成了\Emph{人的样式}或\Emph{成为}太阳——你藉此主张试图证明我们的耶稣只是在外形和外表上成了人；愿天主拯救任何信徒不发出这样的断言。现在，至于其余部分，你那观点会把整个事情缩减为一个梦幻，就我们而言，成为纯粹的空想；不仅如此，就连\Quote{降临}这个名称本身也会被取消：因为如果祂如你所说，是神而不是真人，那么祂本可以仍坐在天上而随己意行事。但事实并非如此，\BibleQuote{他贬抑自己，取了奴仆的形体}\BibleRef{斐 2:7}；我这是指那位由玛利亚成为人的祂而言。为什么呢？难道我们就不能提出一些类似你刚才所讨论的事，而且甚至更轻易、更宽松地提出吗？但我们绝不偏离真理\BibleQuote{一撇或一画}\BibleRef{玛 5:18}。因为那位由玛利亚所生的是子，祂自选了承担这场大能的斗争——即耶稣。这是天主的基督，祂降临到了那位由玛利亚而来者身上。然而，如果你甚至拒绝相信那从天上发出的声音，那么你所能提出的一切都只是你自己的一些妄断；即使你对此自我声言，也没有人会相信你。因为耶稣立刻被圣神领到旷野，为受魔鬼的试探；魔鬼对祂缺乏正确的认识，就向祂说：\BibleQuote{你若是天主子}\BibleRef{玛 4:3}。除此之外，魔鬼并不理解这位\Emph{由玛利亚}所怀的圣子，祂宣讲天国的福音，祂也确实拥有一个伟大的帐幕，一个无人能预备的帐幕：由此，那位被钉在十字架上的，从死者中复活后，就被接升到那里，天主子基督统治的地方；所以当祂开始施行审判时，那些未曾认识祂的人，将要\BibleQuote{瞻望他们所剌透的}\BibleRef{若 19:37}。但为了让你信服，我向你提出这个问题：为什么祂的门徒与祂同住一整年，除了那一次祂的面容如太阳般发光的时候外，没有一人俯伏在祂面前，正如你刚才所说的？岂不是因为那\Emph{为祂}由玛利亚所造的帐幕吗？因为正如除了门徒和圣\Person{保禄}之外，无人有足够的能力承受护慰者的重担，同样，除了那位由玛利亚所生、超乎所有圣人之上者——即耶稣——之外，也无人能承受那从天降下的圣神降临，父的声音也藉着祂作证说：\BibleQuote{这是我的爱子}\BibleRef{玛 3:17}。但现在，就我向你提出的这些问题，请你给出答复。如果你认为祂仅在外貌和外形上如此，那么祂怎会被那些由男女所生的人——即法利塞人——捉拿并拖去审判呢？因为一个属灵的身体无法被物质性的肉体所捉拿。但是，如果你至今仍未对我提出的论证作出答复，如今却有任何话可回应我所引证的言辞和命题，那么，我请你，至少给我取一把或一些你的日光来。然而，那太阳本身，因其更为精微的身体，足以遮蔽包裹你；而你，即便将它踩在脚下，也不能加给它丝毫伤害。然而，我的主耶稣，如果祂曾被捕，那是作为一个被人所捕的人。如果祂不是人，祂就不会被捕。如果祂未被捕，祂就既没有受苦，也没有受洗。如果祂没有受洗，那么我们之中就没有任何人受洗。如果没有圣洗，也就没有罪过的赦免，每个人都将死在自己的罪中。\Emph{\Person{摩尼}说}：那么，圣洗是为了罪过的赦免而设立的吗？\Emph{\Person{阿基劳斯}说}：当然。\Emph{\Person{摩尼}说}：那么，基督既受了洗，岂不就推论出祂犯了罪吗？\Emph{\Person{阿基劳斯}说}：万万不可！恰恰相反，祂是替我们成了罪，亲自承担了我们的罪。为此缘故，祂由女人所生，也为此缘故，祂接受了圣洗，好使祂能获得这部分的净化，从而使祂所取的身体能够承载那以鸽子的形状降下的圣神。\par

当\Person{阿尔刻劳斯}结束这篇讲话时，群众对他的教义之真实感到惊奇，以大声的呼喊表达了对这位人物的热烈赞扬；他们如此尽力，甚至要阻止他离去。然而，随后他们还是散去了。过了一些时候，当他们再次聚集时，\Person{阿尔刻劳斯}说服他们同意他的请求，安静地聆听道理。听众人中不但有与\Person{狄奥多鲁斯}在一起的人，还有从本省和邻近地区来的所有人。当安静下来后，\Person{阿尔刻劳斯}开始对他们这样谈论\Person{摩尼}：\Quote{你们确实已经听到了我们所传教义的性质，也对我们信仰获得了一些证明；因为我已在你们众人面前讲解圣经，完全按照我自己研究时所领悟的观点。但现在我恳请你们静静听我，我尽可能简短地讲，好让你们明白这位出现在我们中间的人是谁，他从哪里来，以及他有什么品性，正如他的一个同伴、名叫\Person{西西尼乌斯}的人向我指明的；如果你们愿意，我也准备好传召此人来为我将要陈述的事作证。事实上，即使在\Person{摩尼}在场时，此人也不拒绝肯定我们现在所引证的同样事实；因为上述的这个人成了我们教义的信者，与我在一起的另一个人、名叫\Person{图尔博}的也是如此。因此，这些人在他们的见证中传达给我的一切，以及我们自己在\Person{摩尼}身上发现的一切，我都不会向你们隐瞒。}\par

于是，群众更加激动，聚集起来要听\Person{阿尔刻劳斯}讲话；因为说实话，他的陈述给他们带来了极大的乐趣。因此，他们急切地催促他讲出他乐意讲的一切，以及他心中所想的一切；并且宣称他们已准备好就在此时此地听他，并保证即使到晚上，直到点灯时分也会留下来。\par

因此，\Person{阿尔刻劳斯}受到他们热诚的鼓舞，充满信心地开始这样讲论：\Quote{我的弟兄们，你们确实已经听到了关于我主耶稣的主要缘由——我是指那些由律法和先知而决定的缘由；关于我们的救主、我主耶稣基督的辅助缘由，你们也不是不知道。我何必多说呢？出于对救主的切爱，我们被称为基督徒，正如全世界所证实的，也正如宗徒们所清楚宣告的。不但如此，那位最杰出的工师、\Person{保禄}本人，已经立好了我们的根基，也就是教会的根基，并将律法托付给我们，在其中设立了\OriginalTerm{圣职人员}{ministers}、\OriginalTerm{长老}{presbyters}和主教，并分别在指定的地方描述了天主的圣职人员应以何种方式、何种品格行事，长老应具有何种声望，应如何被设立，以及那些渴望主教职务的人应成为怎样的人。所有这些制度，曾一度为我们妥善而正确地确立，至今在我们中间保持着良好的地位和秩序，并且这些规则的正常施行仍在我们中间继续。但是，至于这个名叫\Person{摩尼}的家伙，他目前从\Place{波斯}省嚣张地突然出现在我们中间，我和他之间已经进行了第二次辩论，我要详细告诉你们他的家系，我也要非常清楚地指出他的教义从何源头而来。这个人既不是这类教义的第一个、也不是唯一的创立者。相反，一个属于\Place{斯基泰}、名叫\Person{斯基提亚努斯}的人，生活在宗徒时代，乃是这个教派的创始人与领袖，正如许多其他背教者自命为创始人与领袖一样，他们不时出于妄图攫取更高地位的野心，以谎言取代真理，将较单纯的民众引入自己的私欲，然而你们目前不容我们详述他们的名字与背叛。就是这个\Person{斯基提亚努斯}，引进了这种自相矛盾的\OriginalTerm{二元论}{dualism}；对此，他本人是受惠于\Person{毕达哥拉斯}，正如这一切教义的所有其他追随者一样，他们都持守二元论的观念，偏离圣经的正道：但他们在其中将不会获得任何进一步的成就。}\par

52. 然而，从未有人像这位\Person{斯基提亚努斯}那样，在宣扬这些教义时如此厚颜无耻地推进。因为他引入了两个\OriginalTerm{未受生者}{unbegottens}之间纷争的观念，以及由此类立场而产生的所有其他幻想。\Person{斯基提亚努斯}本人属于\Place{萨拉森}人的族系，并娶了一位来自\Place{上底比斯}的女俘为妻，她说服他居住在\Place{埃及}而不是沙漠中。但愿他从未被那个行省接纳过，在那里，他住了一段时间，获得了学习\Place{埃及}人智慧的机会！因为，说真的，他是一个才华横溢的人，而且家境殷实，正如那些认识他的人在我们所收到的记述中所证实的那样。此外，他有一个名叫\Person{特雷宾托斯}的门徒，为他写了四本书。他给第一本书定名为\WorkTitle{奥秘}，第二本为\WorkTitle{头}，第三本为\WorkTitle{福音}，最后一本为\WorkTitle{宝藏}。他拥有这四本书，以及这位名叫\Person{特雷宾托斯}的门徒。那时，这两个人决定单独一起居住相当长一段时间，\Person{斯基提亚努斯}便想前往\Place{犹太}远足，目的是拜会所有在当地享有盛誉的教师；但此后不久，他突然去世，未能有任何建树。而且，那位与他同住的门徒不得不逃跑，前往\Place{巴比伦尼亚}，该行省目前由\Place{波斯}人控制，现在离我们这里大约六天六夜的路程。到达那里后，\Person{特雷宾托斯}成功地让自己一篇奇妙的叙述流传开来，宣称自己充满了\Place{埃及}人的一切智慧，且他现在真正的名字不是\Person{特雷宾托斯}，而是另一个叫\Person{布达斯}的，这个名号已经被加在他身上。他还进一步声称，他是某位贞女之子，并由一位天使在山上抚养长大。然而，有一位名叫\Person{帕库斯}的先知，以及\Person{密特拉}的儿子\Person{拉布达库斯}，指控他撒谎，并且日复一日地，他们就此事进行了激烈而高深的争论。但我又何必详细讲述这些呢？尽管他屡遭责备，但他仍然继续向他们宣讲关于世界之前之事、天体、以及两个光体的问题；此外，还探讨灵魂去向何处、以何种方式离去，以及以何种方式重返身体的问题；他还发表了许多此类性质的断言，以及比这些更糟糕的其他说法——例如，在元素中与天主挑起战争，以便这位先知自己能够被相信。然而，由于他为诸如此类的断言所迫，他便带着他的四本书投奔了一位寡妇：因为他在那个地方没有收任何门徒，唯独只有一位老妇与他过从甚密。后来，有一次，清晨天刚破晓，他上到某座房屋的顶部，开始呼唤某些名字，据\Person{图尔博}告诉我们，只有那七位被选者才学过这些名字。那么，他上到房顶，目的是进行某种宗教仪式，或施展自己的某种法术；他独自上去，以免被任何人发现：因为他认为，如果他被定以玩弄或轻视民众宗教信仰之罪，就会受到那地方的真正王公们的惩罚。正当他在心中反复思量这些事时，天主以其至公义决定，他应被一个灵推入地下；立刻，他从屋顶被抛下；他的身体毫无生气地坠落在地，被上面提到的那位老妇怜悯地收殓，安葬在惯常的墓地。\par

这事以后，斯基提亚诺斯从\Place{埃及}带来的所有物品都归她所有。她因他的死大为欢喜，这有两个理由：第一，因为她对他的法术并不满意；第二，因为她获得了如此一笔遗产，价值巨大。但她独自一人，便考虑找个人来侍候她；为此她找了一个大约七岁的男孩，名叫\Person{科耳比基乌斯}，她立即给了他自由，并教他读书识字。这男孩到了十二岁时，老妇人去世了，把她所有的财产都留给了他，其中包括\Person{斯基提亚诺斯}撰写的那四本书，每本书的篇幅都不长。女主人埋葬后，\Person{科耳比基乌斯}开始自行使用所有留给他的财产。他离开旧居，搬到城市中心，\Place{波斯}国王的住所就在那里；他在那里改了名，称自己为\Person{马内斯}，不再叫\Person{科耳比基乌斯}，或者更准确地说，不是\Person{马内斯}，而是\Person{马尼}：因为这是\Place{波斯}语言中的一种屈折形式。当这男孩长到将近六十岁时，他已精通当地教授的各门学问，我几乎可以说他在这些方面超越了所有其他人。尽管如此，他对这四本书中包含的教义钻研得更加勤奋；他还收了三个门徒，名字叫\Person{多默}、\Person{阿达斯}和\Person{赫尔玛}。然后，他也取了这些书，如此抄录它们，以致于在其中加入了许多自己的新东西，这些东西只能比作老太婆的荒诞故事。这样，这三个门徒因此依附于他，成为他邪恶计划的知情参与者；而且，他把自己的名字加在书上，删除了前主人的名字，好像这些书都是他自己编写的一样。然后，他觉得最好是差派他的门徒，带着他写在书上的教义，前往那省的上部地区，经过许多城镇村庄，以争取追随者。于是\Person{多默}决定占领\Place{埃及}地区，\Person{阿达斯}占领\Place{斯基提亚}地区，而只有\Person{赫尔玛}选择留在那人身边。当他们出发上路后，国王的儿子得了一种病；国王非常希望看到他痊愈，便发布了一道法令，提供大量赏金，并承诺将赏给任何能治好王子的人。听到这个消息，\Emph{全然不加思索地}，就像那些习惯于玩立方游戏（即骰子的别称）的人那样，\Person{马内斯}来到国王面前，宣称他能治好那孩子。国王听到后，便礼貌地接见了他，并衷心欢迎。但为避免听众彻底厌倦于讲述他所做的许多事，我只说那孩子死在他的手下，或者更确切地说，被他夺去了生命。于是国王下令将\Person{马内斯}投入监牢，并给他戴上重达半英担的铁链。此外，他那两个被派往各城宣扬他教义的门徒也被搜捕以施惩罚。但他们逃跑了，然而从未停止在他们所到之处传入他们那种与信仰相悖、纯由\OriginalTerm{假基督}{Antichrist}所感的教导。\par

54. 但这些事之后，他们回到他们的师傅那里，报告了发生的一切；同时他们也了解到他遭遇到的无数祸患。因此，正如我所说的，当他们见到他时，他们敦促他注意他们在那一个个地区中所必须忍受的苦难；至于其他事项，他们力劝他现在应考虑安全问题；因为他们一直深恐他所遭受的任何灾祸会临到自己身上。但他劝他们什么也不要怕，并起身向他们发表了一篇演讲。然后，当他还在狱中时，就命令他们去获取基督徒经书的抄本；因为这些被他派遣到不同团体中的门徒，已为所有人憎恨，尤其是那些视基督徒之名极为尊荣的人。因此，领到少量钱款后，他们出发前往那些出版基督徒经书的地区；并假装自己是基督徒的信使，请求向他们出示这些书，以便他们能获得抄本。简而言之，他们就这样获得了我们圣经的全部经书，并带回到仍在狱中的他们的师傅那里。那个狡猾的家伙拿到这些抄本后，就开始搜寻我们书中所有看似支持他二元论观点的陈述；不过，那其实不是他自己的观点，而是\Person{斯基提亚努斯}的观点，斯基提亚努斯在他很久以前就公布了这一观点。正如他与我辩论时所做的，那时他也如此，通过拒绝我们圣经中的某些内容，篡改另一些内容，试图证明这些经文支持他自己的教义，只不过基督的名字附着其上。因此，他假借这个名字自称，以便各团体中的人们，听到基督的神圣尊名时，就不会有憎恶和骚扰他那些门徒的意图。此外，当他们看到我们圣经中论及\OriginalTerm{护慰者}{Paraclete}的那话时，他竟自以为他本人可能就是那位护慰者；因为他没有足够仔细研读，未曾注意到护慰者早已来了——也就是，在宗徒们仍在地上之时。因此，当他炮制了这些亵圣的发明之后，他也派遣他的门徒大胆地宣讲这些杜撰和错谬，并使这些虚假新奇的话传遍四方。但波斯国王闻知此事后，准备对他施以应得的惩罚。然而，\Person{摩尼}在梦中得到警告，得知了国王的意图，便逃出了监狱，并得以成功脱逃，因为他用巨款贿赂了看守之人。后来，他住进了\Place{阿拉比雍}城堡；并从那里通过\Person{图尔伯}之手寄出了他写给我们的\Person{马塞勒斯}的信，在信中他表明想去拜访他。他到达那里后，他与我之间进行了一场辩论，就如你们在这里所观察到和听到的辩论一样；在那场讨论中，我们尽力表明他是个假先知。我还可以补充说，让他逃脱的狱卒受到了惩罚，国王下令无论何处一经发现此人，就要将他缉拿归案。这些事都是我亲身确知的，因此我有必要也让你们知道，直到今天，波斯国王仍在追捕这个家伙。\par

55. 众人听到这话，就想抓住\Person{摩尼}，把他交给邻近的外族人，就是那些住在\Place{斯特兰加河}对岸的人，尤其是因为不久前，已有一些人前来寻索他；不过当时他正在逃亡之中，他们找不到他的任何踪迹，只得离去。然而，当\Person{阿基劳斯}作出这番声明时，\Person{摩尼}立刻逃跑，而且没等任何人追来，就成功地脱身了。因为众人都被\Person{阿基劳斯}叙述的话留住了，他们津津有味地听着；尽管如此，仍有少数人紧紧追赶他。但他再次沿来时的路逃走，过了河，回到了\Place{阿拉比翁}城堡。不过，后来他在那里被捕，并解送到国王面前；国王对他极其忿恨，又因切望为两条命报仇——就是他自己的儿子和掌狱官的死——就下令将他剥皮，挂在城门上，又把他的皮浸以某种药物，充气鼓胀；并下令把他的肉给飞鸟作食。这些事后来为\Person{阿基劳斯}所知时，他在先前的辩论之外又补充了\Emph{有关这些事的叙述}，好让所有的事实都为众人知晓，正如我——记述这些事的人——在前面已说明的那样。于是，所有基督徒都聚集起来，决议应当作出对他的裁决，以作为他生命结语的一种附录，这结语与他生平的其它情况非常相合。此外，\Person{阿基劳斯}又说了大意如下的话：——我的弟兄们，你们不要不相信我所说的：我指的是\Person{摩尼}并非这邪恶教义的首创者，而只是他在世上的某些地方将其公之于众而已。因为，一个人仅仅将某事物带到某个地方，并不能立刻算作那事物的创始者；唯有那发现者才配得那种荣誉。就如舵手接过别人建造的船，可以随己意驶往他国，但他与船的构造并无关系一样；这个人的情形也应如此理解。因为，实际上他从起初并没有赋予这事以起源；而只是把别人所发现的传递给人，我们从可信的见证得知，以此为基础，我们要向你们证明：这邪恶的发明并非出自\Person{摩尼}，而是源自另一个人，那个人确实还是个外族人，在他之前许久就已出现。再者，这教义一度未曾刊行，直到那些隐藏了一段时间的道理终于被他当作自己的，公开推了出来，而作者的署名已被删除，正如我在上面所表明的。在\OriginalTerm{波斯人}{Persians}中，也有一位宣传类似观点的人，名叫\Person{巴西里德斯}，年代更早，生活在我们宗徒之后的时期不久。此人天性机敏，他观察到当时所有其它论题已被占据，便决意主张\Person{斯库提亚努斯}所持的同一\OriginalTerm{二元论}{dualism}。最后，由于他没有什么真正自己的东西可以提出，就把他人的言论带到对手面前。他所有的书卷都包含着一些既艰深又极苛刻的内容。不过，他的\WorkTitle{讲论集}第十三卷仍存于世，开头如下：\Quote{在撰写我们的\WorkTitle{讲论集}第十三卷时，那整全的话语为我们提供了必需而多产的话语。}然后他阐述说，\Emph{善与恶之间的对立}，是在富之本原与贫之本原的比喻下产生的，后者本性无根无所，只是外加于事物之上。这就是这卷书所包含的唯一主题。那么，它岂不是包含了一种奇怪的话语吗？而且，正如有些人曾有这种看法，你们岂不也会对那卷书如此开头的本身感到反感吗？——但\Person{巴西里德斯}在有了大约五百行的引言之后，回到主题上，又继续说：\Quote{你们要放弃这些空虚而好奇的变异，不如去查考外族人对善恶的探究，以及他们对这一切所持的看法。因为他们当中有些人主张，万物有两个本原，他们把善和恶归于这两个本原，认为这些本原无始无生；这就是说，在事物的本源中有光和暗，它们自己存在，而不仅是宣称存在而已。当这些本原自己存在时，它们各自过着自己适意且力所能及的生活；因为在万有中，凡是适合某一物的，便与之相合，没有一物在自己看来为恶。但及至它们彼此认识，黑暗开始凝视光明之后，那时黑暗仿佛对某种更优越之物燃起了爱慕，便强求与光明交合。}\par

% structure: p0067 source=h2 target=section reason=relative-heading-rank; base=h2

\section{同一辩论的片段。}

\Emph{\Person{西里尔}以如下措辞介绍了这个片段：}——他，即\Person{摩尼}，从监狱逃出，来到了\Place{美索不达米亚}；但他在那里遇到了那位正义的盾牌，主教\Person{阿基劳斯}。为了在哲学法官面前考验他，\Person{阿基劳斯}召集了一群希腊听众，以免有人指责法官偏私，若他们是基督徒就可能如此。\Emph{然后事情的进展如下所述：}——\par

1. \Person{阿基劳斯}对\Person{摩尼}说：现在请你陈述你所宣扬的教义。——于是那人口如敞开的坟墓，立刻开始亵渎万物的创造者，说：\WorkTitle{旧约}的天主是邪恶的创造者，他曾这样论到自己：\BibleQuote{我是吞噬的烈火。}——但睿智的\Person{阿基劳斯}彻底驳斥了这亵渎。因为他说道：若按你的指控，\WorkTitle{旧约}的天主自称是烈火，那么那位说\BibleQuote{我来是为把火投在地上}的，又是谁的儿子呢？你若挑剔那位说\BibleQuote{主使人死，也使人活}的，又为何尊崇曾复活\Person{塔彼达}、却也处死\Person{撒斐辣}的\Person{伯多禄}呢？再者，你若因那位预备了烈火而挑剔他，又为何不挑剔另一位曾说\BibleQuote{离开我，到永火里去吧}的呢？你若挑剔那位曾说\BibleQuote{是我，天主，创造平安，也降灾祸}的，就请向我们解释，\Person{耶稣}又为何说：\BibleQuote{我来不是为带平安，而是带刀剑。}既然两人用同样的方式说话，那么二者必居其一：要么他们都是善的，因为用了同样的语言；要么，若\Person{耶稣}说了那样的话仍不受指责，你必须告诉我们，你为何责难那位在\WorkTitle{旧约}中用类似方式说话的了。\par

2. \Person{摩尼}这样回答他：那么，使人眼瞎的，到底是怎样的一位天主呢？因为\Person{保禄}说过这样的话：\BibleQuote{此世的神弄瞎了不信者的心意，免得福音的光照耀他们。}但\Person{阿基劳斯}打断了他，极有力地驳斥说：不过，请读一读这节前面的几句吧，即：\BibleQuote{但如果说我们的福音也被蒙蔽了，那只是为丧亡的人蒙蔽着。}你看，这福音是在丧亡的人身上蒙蔽了。\BibleQuote{你们不要把圣物给狗。}再者，难道只有\WorkTitle{旧约}的天主弄瞎了不信者的心意吗？难道\Person{耶稣}自己不是也曾说：\BibleQuote{为此，我用比喻对他们讲话，是因为他们看，却看不见。}那么，是因为他憎恨他们，才不愿他们看见吗？抑或\Emph{并非}因他们不配，因为他们自己闭上了眼？因为每逢邪恶是自行选择的，那里也就没有恩宠。\BibleQuote{因为凡有的，还要给他；凡没有的，连他自以为有的，也要从他夺去。}\par

3. 然而，纵使我们不得不接受某些人所倡导的解释——因为这题目并非全无注意的价值——并且不得不声称他确实弄瞎了不信者的心意，我们仍然必须肯定，他弄瞎他们是为了善，好叫他们恢复视力，能看见圣善的事物。因为经上并没有说他弄瞎了他们的灵魂，而只说他弄瞎了不信者的心意。而这种表达方式大致有这样的意思：弄瞎那嫖客淫荡的心意，这人便得救了；弄瞎那盗贼贪婪偷窃的心意，这人便得救了。难道你不愿这样理解这句话么？好罢，还有另一种解释。太阳使视力不良的人目眩；那些眼睛流泪的人，当遭受强光照射时也会目眩，然而这并不是因为太阳的本性使人目眩，而是因为眼睛自身的结构无法正常视物。同样地，那些内心沾染不信毛病的人，也不能仰观\Emph{光荣的}天主性的光芒。再者，经上并没有说：\BibleQuote{他弄瞎了他们的心意，免得他们听见福音，}反而是说：\BibleQuote{免得我们的主耶稣基督荣耀的福音之光，照耀他们。}因为聆听福音是向所有人交付的事务；但基督福音的光荣，只分施给真诚和纯正的人。为此，主对那些无法聆听的人用比喻讲话，却私下向他的门徒解释这些比喻。因为光荣的光照是为那些已蒙光照的人，而弄瞎却是为那些不信的人。这些奥迹，就是教会现在向你们这些从慕道者行列中迁移出来的人所宣讲的，她向来不会向外邦人宣讲。因为我们不会向外邦人宣讲有关圣父、圣子及圣神的奥迹；我们也不会在慕道者面前明明地宣讲这些奥迹；许多时候，我们以隐晦的方式表达，好叫那些有理解力的信徒能够领悟所指的真理，而那些没有这理解力的人则不致受害。\par

\SourceRecord{Translated by S.D.F. Salmond.；Ante-Nicene Fathers；Vol. 6.；Edited by Alexander Roberts, James Donaldson, and A. Cleveland Coxe.；Buffalo, NY: Christian Literature Publishing Co.,；1886.；<http://www.newadvent.org/fathers/0616.htm>.}
